% Ordered Associative Chiral Koszul Duality
% Integration-ready chapter file (stripped from standalone amsart draft).
% Uses only \providecommand for macros that may not be in main.tex preamble.

\providecommand{\Assch}{\mathrm{Ass}^{\mathrm{ch}}}
\providecommand{\Barch}{\overline{B}^{\mathrm{ch}}}
\providecommand{\Barchord}{\overline{B}^{\mathrm{ord}}}
\providecommand{\Cobar}{\Omega^{\mathrm{ch}}}
\providecommand{\coHoch}{\operatorname{coHH}}
\providecommand{\Cotor}{\operatorname{Cotor}}
\providecommand{\Coext}{\operatorname{Coext}}
\providecommand{\RHom}{R\!\operatorname{Hom}}
\providecommand{\Tot}{\operatorname{Tot}}
\providecommand{\KK}{\mathbb{K}}
\providecommand{\Dpbw}{D^{\mathrm{pbw}}}
\providecommand{\Dco}{D^{\mathrm{co}}}
\providecommand{\chotimes}{\mathbin{\otimes^{\mathrm{ch}}}}
\providecommand{\wt}{\widetilde}
\providecommand{\eps}{\varepsilon}
\providecommand{\susp}{s}
\providecommand{\coeq}{\operatorname{coeq}}
\providecommand{\id}{\mathrm{id}}
\providecommand{\op}{\mathrm{op}}
\providecommand{\cop}{\mathrm{cop}}
\providecommand{\HH}{\operatorname{HH}}
\providecommand{\Tor}{\operatorname{Tor}}
\providecommand{\Ext}{\operatorname{Ext}}
\providecommand{\gr}{\operatorname{gr}}
\providecommand{\MC}{\operatorname{MC}}
\providecommand{\Ob}{\operatorname{Ob}}
\providecommand{\eq}{\operatorname{eq}}
\providecommand{\cF}{\mathcal{F}}
\providecommand{\cL}{\mathcal{L}}
\providecommand{\cA}{\mathcal{A}}

\chapter{Ordered Associative Chiral Koszul Duality}
\label{ch:ordered-associative-chiral-kd}
\label{chap:ordered-associative}% phantom: referenced from frontier/hochschild chapters
\label{ch:ordered-chiral}%% alias resolving legacy \ref sites in sc_chtop_heptagon

\section{Introduction: the diagonal as the universal boundary}

\begin{remark}[$E_1$ primitive layer]
\label{rem:ordered-kd-e1-primitive}
This chapter is the $E_1$ primitive layer of the theory.
The ordered bar coalgebra $\Barch(A)$ with its deconcatenation coproduct
is the native object; the symmetric bar $B^{\Sigma}(\cA)$ of Volume~I's
modular theory is its $\Sigma_n$-coinvariant quotient under the
averaging map $\mathrm{av}\colon \mathfrak{g}^{E_1}_\cA \to
\mathfrak{g}^{\mathrm{mod}}_\cA$.
The $R$-matrix, Drinfeld associator, and higher Yangian coherences
computed here are the $E_1$ MC data whose $\mathrm{av}$-image is
the modular characteristic $\kappa(\cA)$ and the shadow obstruction tower.
In particular, $\kappa(\cA)$ is the degree-$2$ image of $\mathrm{av}$
(for abelian algebras $\kappa = \mathrm{av}(r(z))$;
for non-abelian Kac--Moody
$\kappa = \mathrm{av}(r(z)) + \dim(\fg)/2$).
See Volume~I, \S\textup{E$_1$ as primitive}
for the full statement.
\end{remark}

The open-color projection of $\Theta^{\mathrm{oc}}$ is the $E_1$-chiral Koszul duality: the ordered bar coalgebra $\Barch(A)$ with its deconcatenation coproduct, the diagonal bicomodule $C_\Delta$, and the duality involution that sends $A$ to~$A^!$. The commutative shadow of chiral Koszul duality is by now
classical. The genuinely associative story, the
open-color projection of the $\SCchtop$-algebra structure on
the Koszul dual, is different. It is not the
Francis--Gaitsgory equivalence between chiral Lie algebras
and cocommutative factorization coalgebras. It is a theory
in which ordered configuration spaces survive, shuffles do
not collapse under symmetric-group quotients, and after
linear duality chiral algebras appear again on the far side
(the duality involution of
Theorem~\ref{thm:duality-involution}).

Begin with the simplest case. The Heisenberg algebra~$\cH_k$
has ordered bar complex
$\Barchord(\cH_k)\simeq\bigoplus_{n\ge 1}
(s^{-1}\alpha)^{\otimes n}$, the cofree tensor coalgebra on the
desuspended generator, with differential
$d(s^{-1}\alpha\otimes s^{-1}\alpha)=\pm k\,s^{-1}(1)$
from the double-pole OPE
$\alpha_{(1)}\alpha=k$. The ordering is load-bearing:
the collision residue $r(z)=k/z$ (one pole absorbed by the
$d\log$ kernel, per the bar construction) changes sign under
reversal $z\mapsto -z$, and the $R$-matrix $R(z)=z^k$ is
nontrivial. The coproduct is ordered deconcatenation
$\Delta[a_1|\cdots|a_n]=\sum_{p=0}^{n}
[a_1|\cdots|a_p]\otimes[a_{p+1}|\cdots|a_n]$,
cutting words at a single point. Left and right
deconcatenation make~$\Barchord(\cH_k)$
into a bicomodule over itself, the coalgebra-side shadow of
the diagonal bimodule ${}_A A_A$ for $A=\cH_k$. This is
the structure that encodes boundary conditions.

To make the associative picture rigid, we return to first principles and ask a single question.

\medskip

\emph{What is the correct coalgebra-side avatar of the diagonal bimodule ${}_A A_A$?}

\medskip

If $C=\Barch(A)$ is the ordered bar coalgebra of an associative chiral algebra $A$, then the
correct boundary object is the \emph{diagonal bicomodule}
\[
C_\Delta,
\]
namely the coalgebra $C$ equipped with its deconcatenation coproduct acting on the left and on the right.

On the algebra side, interval composition is derived tensor over $A$ and the annulus is Hochschild
homology. On the coalgebra side, interval composition becomes derived cotensor over $C$, the
diagonal becomes the composition unit, and the annulus becomes the derived cotrace of the diagonal
against itself:
\[
\HH^{\mathrm{ch}}_\bullet(A)\simeq C_\Delta \square_{C^e}^{\mathbf R} C_\Delta.
\]
The annulus is the self-cotensor of the universal boundary condition.

\subsection*{Dictionary}
The new results of this chapter can be summarized by the following dictionary:
\[
A \longleftrightarrow C_\Delta,\qquad
A^e \longleftrightarrow C^e:=C\widehat\otimes C^{\cop},\qquad
-\otimes_A^{\mathbf L}- \longleftrightarrow -\square_C^{\mathbf R}-,
\]
\[
\HH^{\mathrm{ch}}_\bullet(A,M)
\longleftrightarrow
\coHoch^{\mathrm{ch}}_\bullet(C,\KK_A^{\mathrm{bi}}(M)),
\qquad
\HH_{\mathrm{ch}}^\bullet(A,M)
\longleftrightarrow
\coHoch_{\mathrm{ch}}^\bullet(C,\KK_A^{\mathrm{bi}}(M)).
\]
\emph{The associative dual already remembers the open boundary formalism.}

\subsection*{Ledger of status}
To keep the exposition honest, we separate three classes of statements.

\begin{enumerate}[label=(\roman*)]
\item \textbf{Repository inputs.}
The existence of the associative chiral bar--cobar equivalence, the one-sided module/comodule
Koszul duality, and the enveloping algebra formalism are taken as already established in the
repository.

\item \textbf{New theorems proved here.}
The ordered chiral shuffle theorem, opposite-duality, enveloping duality, the two-sided twisted
bicomodule, the identification of the tangent bar complex with the twisted bicomodule, the PBW-complete
bimodule/bicomodule equivalence, the diagonal correspondence, the Hochschild--coHochschild
dictionary, the ordered open trace formalism, the $R$-matrix as
ordered-to-unordered descent datum with the covering-space frame and descent identification,
the Yang--Baxter equation from $d^2=0$ via boundary-face analysis, the classical
Yang--Baxter equation as linearisation, the affine ordered-bar
realization of the dg-shifted Yangian with Drinfeld generators and
RTT presentation, the ordered AOS relations ($d^2=0$ from
associativity alone at degrees~$3$--$5$), the Euler-eta identity for
graded bar Euler characteristics, and bar cohomology concentration with
the depth spectral sequence.

\item \textbf{Conjectural frontier.}
The Francis--Gaitsgory commutator-shadow spectral
sequence. (The bordered configuration-space realization, including the annular closure,
is now proved: Theorem~\ref{thm:bordered-annular}. The associative modular Maurer--Cartan
class and BRST/Drinfeld--Sokolov compatibility are proved in
\S\ref*{sec:conjectures} of the frontier chapter.)
\end{enumerate}

\section{The strongly admissible regime and the foundational inputs}
\label{sec:setup}

We work over a field $k$ of characteristic $0$. Let $X$ be a smooth curve over $k$, and let
$\mathcal D(X)$ denote the relevant derived category of $D$-modules on $X$ (or on the Ran space of
$X$) endowed with the chiral tensor product
\[
\chotimes := \boxtimes_{D_X}.
\]
An $E_1$-chiral algebra means an algebra over the non-$\Sigma$ operad $\Assch$.

\begin{definition}[Strongly admissible associative chiral algebra]
An augmented $E_1$-chiral algebra $A$ on $X$ is called \emph{strongly admissible} if the following
hold.
\begin{enumerate}[label=(\alph*)]
\item $A$ is pro-nilpotent and equipped with an exhaustive separated descending PBW-type filtration
$F^\bullet A$ which is complete for the induced topology.

\item Every associated graded piece and every finite filtration quotient is bounded below and perfect
enough that completed tensor products, totalizations, and normalized chains behave exact\-ly on these
pieces.

\item The reduced ordered bar coalgebra
\[
C:=\Barch(A)
\]
is conilpotent, and the corresponding coderived category of bicomodules admits enough coflat
resolutions so that derived cotensor products are defined and associative.

\item Whenever linear duality is invoked, each graded/conformal piece of the object under consideration
is finite-dimensional.
\end{enumerate}
\end{definition}

\begin{remark}
This is the fortified regime in which hidden homological-algebra issues are kept on the surface.
The arguments below require no more than these hypotheses; they are chosen so that
every totalization, completion, normalization, and passage to coderived categories is legitimate.
\end{remark}

\begin{hypothesis}[Repository inputs]
\label{hyp:inputs}
We assume the repository has already established the following.
\begin{enumerate}[label=(I)]
\item \textbf{Associative chiral bar--cobar.}
For every augmented pro-nilpotent strongly admissible $E_1$-chiral algebra $A$ there is a
quasi-isomorphism
\[
\Cobar \Barch(A)\xrightarrow{\sim} A,
\]
and, on the appropriate coalgebra side, the corresponding bar--cobar adjunction is an equivalence.

\item \textbf{One-sided module/comodule duality.}
For every such $A$ there is an equivalence between complete left $A$-modules and conilpotent left
$\Barch(A)$-comodules, with the usual coderived/contraderived refinement in the curved or completed
regime.

\item \textbf{Enveloping algebra formalism.}
The chiral enveloping algebra
\[
A^e := A\chotimes A^{\op}
\]
is defined, and $A$-bimodules are equivalently left $A^e$-modules.
\end{enumerate}
\end{hypothesis}

\begin{definition}[Universal twisting cochain]
Let $A$ be strongly admissible and set $C:=\Barch(A)$. The bar--cobar equivalence produces a
canonical twisting cochain
\[
\tau:C\to A
\]
characterized by the Maurer--Cartan equation
\[
d\tau+\tau\star\tau=0.
\]
We call $\tau$ the \emph{universal twisting cochain}.
\end{definition}

\begin{definition}[Opposite and diagonal objects]
Let $A$ be an associative chiral algebra and let $C$ be a coaugmented dg coalgebra.
\begin{enumerate}[label=(\alph*)]
\item $A^{\op}$ is the same underlying object with multiplication reversed.
\item $C^{\cop}$ is the same underlying object with coproduct flipped.
\item The \emph{coalgebraic enveloping object} is
\[
C^e:=C\widehat\otimes C^{\cop}.
\]
\item The \emph{diagonal bicomodule} $C_\Delta$ is $C$ equipped with left and right coactions both
equal to $\Delta_C$.
\end{enumerate}
\end{definition}

\begin{lemma}[Bicomodules as comodules over the enveloping coalgebra; \ClaimStatusProvedHere]
\label{lem:bicom-e}
Let $C$ be a coaugmented conilpotent dg coalgebra. Then left $C^e$-comodules are canonically the
same thing as $C$-bicomodules.
\end{lemma}

\begin{proof}
Given a left $C^e$-coaction
\[
\rho^e:M\to C\otimes C^{\cop}\otimes M,
\]
compose with the counit on the second factor and on the first factor to obtain
\[
\lambda=(\id_C\otimes \eps\otimes \id_M)\rho^e:M\to C\otimes M,
\]
\[
\rho=(\eps\otimes \id_{C^{\cop}}\otimes \id_M)\rho^e:M\to M\otimes C.
\]
These are a left and a right $C$-coaction, and coassociativity of $\rho^e$ is equivalent to the
usual bicomodule compatibility condition.

Conversely, given commuting left and right coactions $\lambda$ and $\rho$, define
\[
\rho^e
=
(\id_C\otimes \tau_{C,M}\otimes \id_M)\circ(\lambda\otimes \id_M)\circ \rho
:
M\to C\otimes C^{\cop}\otimes M,
\]
where $\tau_{C,M}$ swaps the middle two factors. The bicomodule axioms are precisely the
coassociativity and counitality conditions for $\rho^e$.
\end{proof}

\begin{definition}[Chiral Hochschild and coHochschild theories]
Let $A$ be a strongly admissible associative chiral algebra and $C=\Barch(A)$.
For an $A$-bimodule $M$ and a $C$-bicomodule $N$, define
\[
\HH^{\mathrm{ch}}_\bullet(A,M):=\Tor_\bullet^{A^e}(A,M),
\qquad
\HH_{\mathrm{ch}}^\bullet(A,M):=\RHom_{A^e}(A,M),
\]
\[
\coHoch^{\mathrm{ch}}_\bullet(C,N):=\Cotor_\bullet^{C^e}(C_\Delta,N),
\qquad
\coHoch_{\mathrm{ch}}^\bullet(C,N):=\Coext_{C^e}^\bullet(C_\Delta,N).
\]
When $M=A$ and $N=C_\Delta$ we omit coefficients.
\end{definition}

\section{Ordered bar coalgebras and the chiral shuffle theorem}
\label{sec:shuffle}

In the associative chiral regime the basic bar object is ordered. This is the first point at which
the theory diverges sharply from the commutative shadow: \emph{every shuffle survives as a distinct
summand.}

\begin{definition}[The simplicial ordered bar object]
Let $A$ be an augmented associative chiral algebra. Define the simplicial object
$B^\Delta_\bullet(A)$ by
\[
B^\Delta_n(A):=A^{\chotimes n}\qquad (n\geq 0),
\]
with faces given by augmentation at the two ends and multiplication in adjacent slots, and
degeneracies given by unit insertion. Its normalized chains
\[
N(B^\Delta_\bullet(A))
\]
are canonically identified with the reduced ordered chiral bar complex $\Barch(A)$.
\end{definition}

In suspended notation we write a bar word of length $n$ as
\[
[a_1|\cdots|a_n].
\]

\begin{definition}[EZ-admissible pair]
A pair $(A,B)$ of strongly admissible associative chiral algebras is called \emph{EZ-admissible} if
the normalization functor on bisimplicial objects built from $A$ and $B$ commutes, filtration-wise,
with the completed chiral tensor products and totalizations appearing in the proof of the
Eilenberg--Zilber theorem.
\end{definition}

\begin{theorem}[Ordered chiral shuffle theorem; \ClaimStatusProvedHere]
\label{thm:shuffle}
Let $(A,B)$ be an EZ-admissible pair of strongly admissible associative chiral algebras.
There are natural filtered chain maps
\[
\mathsf{Sh}^{\mathrm{ch}}_{A,B}:
\Barch(A)\widehat\otimes \Barch(B)\longrightarrow \Barch(A\chotimes B),
\]
\[
\mathsf{AW}^{\mathrm{ch}}_{A,B}:
\Barch(A\chotimes B)\longrightarrow \Barch(A)\widehat\otimes \Barch(B),
\]
such that
\[
\mathsf{AW}^{\mathrm{ch}}_{A,B}\circ \mathsf{Sh}^{\mathrm{ch}}_{A,B}=\id,
\qquad
\mathsf{Sh}^{\mathrm{ch}}_{A,B}\circ \mathsf{AW}^{\mathrm{ch}}_{A,B}\simeq \id.
\]
Hence there is a filtered quasi-isomorphism of conilpotent dg coalgebras
\[
\Barch(A)\widehat\otimes \Barch(B)\xrightarrow{\sim}\Barch(A\chotimes B).
\]
\end{theorem}

\begin{proof}
Set
\[
K_{\bullet,\bullet}:=B^\Delta_\bullet(A)\chotimes B^\Delta_\bullet(B).
\]
This is a bisimplicial object in the ambient abelian or stable category of $D$-modules.

Its diagonal identifies canonically with the simplicial bar object of the tensor-product algebra:
\[
\operatorname{diag}(K)_n
=
A^{\chotimes n}\chotimes B^{\chotimes n}
\cong
(A\chotimes B)^{\chotimes n}
=
B^\Delta_n(A\chotimes B).
\]
This identification is compatible with faces and degeneracies because multiplication and augmentation
on $A\chotimes B$ are componentwise.

Now apply the classical Eilenberg--Zilber theorem filtration-wise in each exact graded piece.
EZ-admissibility ensures that normalization commutes with the completed tensor constructions used
here, so one obtains a filtered quasi-isomorphism
\[
N(\operatorname{diag}K)\xrightarrow{\sim} \Tot\bigl(NK_{\bullet,\bullet}\bigr),
\]
with inverse given by Alexander--Whitney. Since
\[
N(B^\Delta_\bullet(A))\cong \Barch(A),
\qquad
N(B^\Delta_\bullet(B))\cong \Barch(B),
\qquad
N(B^\Delta_\bullet(A\chotimes B))\cong \Barch(A\chotimes B),
\]
this yields the claimed quasi-isomorphism.

The formulas are the usual normalized Eilenberg--MacLane formulas, now interpreted in the ordered
chiral setting. Explicitly, for bar words
\[
[a_1|\cdots|a_p]\in {\Barch}_p(A),\qquad [b_1|\cdots|b_q]\in {\Barch}_q(B),
\]
set $c_i=a_i\otimes 1$ for $1\leq i\leq p$ and $c_{p+j}=1\otimes b_j$ for $1\leq j\leq q$, and define
\[
\mathsf{Sh}^{\mathrm{ch}}_{p,q}
([a_1|\cdots|a_p]\otimes[b_1|\cdots|b_q])
=
\sum_{\sigma\in \operatorname{Sh}(p,q)}
\varepsilon(\sigma)\,
[c_{\sigma^{-1}(1)}|\cdots|c_{\sigma^{-1}(p+q)}].
\]
Similarly,
\[
\mathsf{AW}^{\mathrm{ch}}([x_1|\cdots|x_n])
=
\sum_{p=0}^n
[\pi_A(x_1)|\cdots|\pi_A(x_p)]\otimes[\pi_B(x_{p+1})|\cdots|\pi_B(x_n)],
\]
where
\[
\pi_A=\id_A\otimes \eps_B,\qquad \pi_B=\eps_A\otimes \id_B.
\]
The chain-map property uses only the simplicial identities.

Finally, the shuffle map is compatible with deconcatenation: cutting a shuffled word is equivalent to
cutting the source words and then shuffling the left halves and the right halves. Therefore
$\mathsf{Sh}^{\mathrm{ch}}_{A,B}$ and $\mathsf{AW}^{\mathrm{ch}}_{A,B}$ are coalgebra maps.
\end{proof}

\begin{remark}
The theorem is genuinely associative. In a commutative or $E_\infty$ setting many of the shuffle
summands are identified after the symmetric-group quotient. Here they remain distinct. The ordered
configuration space remembers their full combinatorics.
\end{remark}

\section{Opposite-duality, anti-involutions, and enveloping duality}
\label{sec:opposite}

The second foundational phenomenon is order reversal. It transforms opposite multiplication on the
algebra side into opposite deconcatenation on the bar side.

\begin{theorem}[Opposite-duality for ordered bar coalgebras; \ClaimStatusProvedHere]
\label{thm:opposite}
Let $A$ be a strongly admissible associative chiral algebra. There is a canonical isomorphism of
coaugmented conilpotent dg coalgebras
\[
\mathsf{R}_A:\Barch(A^{\op})\xrightarrow{\sim}\Barch(A)^{\cop}.
\]
If $A$ is finite-type Koszul so that linear duality returns an associative chiral algebra, then
\[
(A^{\op})^!\simeq (A^!)^{\op}.
\]
\end{theorem}

\begin{proof}
On suspended bar words define
\[
\mathsf{R}_A[s^{-1} a_1|\cdots|s^{-1} a_n]
=
(-1)^{\sum_{i<j}(|a_i|+1)(|a_j|+1)}
[s^{-1} a_n|\cdots|s^{-1} a_1].
\]
Geometrically, this is induced by the order-reversal involution on the ordered configuration space.

The bar differential is the alternating sum of adjacent collision maps. The $i$-th face for
$A^{\op}$ multiplies the $i$-th and $(i+1)$-st slots in the opposite order. Reversal carries that
face to the $(n-i)$-th face for $A$ with exactly the sign written above. Hence $\mathsf{R}_A$ is a
chain map.

The coproduct on $\Barch(A)$ is deconcatenation. Reversal sends a cut after the $p$-th letter to a
cut before the $(n-p)$-th letter. Therefore
\[
\Delta\circ\mathsf{R}_A
=
\tau\circ(\mathsf{R}_A\otimes\mathsf{R}_A)\circ\Delta,
\]
where $\tau$ flips the tensor factors, which is exactly the statement that the target is the
opposite coalgebra. The claim about linear duals follows because linear duality exchanges opposite
coproduct with opposite product.
\end{proof}

\begin{corollary}[Anti-involutions survive duality; \ClaimStatusProvedHere]
\label{cor:anti}
Any anti-involution $\iota:A\to A^{\op}$ induces an anti-involution
\[
\iota^!:A^!\to (A^!)^{\op}.
\]
If $\iota^2=\id$, then $(\iota^!)^2=\id$.
\end{corollary}

\begin{proof}
Apply Theorem~\ref{thm:opposite} functorially.
\end{proof}

\begin{lemma}[Closure of admissibility under opposite and enveloping constructions; \ClaimStatusProvedHere]
\label{lem:closure}
If $A$ is strongly admissible, then so are $A^{\op}$ and $A^e=A\chotimes A^{\op}$.
\end{lemma}

\begin{proof}
All hypotheses in the definition of strong admissibility are stable under replacing multiplication by
its opposite and under taking completed tensor products of filtered objects.
\end{proof}

\begin{corollary}[Enveloping duality; \ClaimStatusProvedHere]
\label{cor:enveloping}
Let $A$ be strongly admissible, and assume $(A,A^{\op})$ is EZ-admissible. Set $C=\Barch(A)$.
Then there is a canonical quasi-isomorphism of conilpotent dg coalgebras
\[
\Barch(A^e)\xrightarrow{\sim} C^e=C\widehat\otimes C^{\cop}.
\]
If $A$ is finite-type Koszul, then
\[
(A^e)^!\simeq (A^!)^e.
\]
\end{corollary}

\begin{proof}
By the ordered shuffle theorem,
\[
\Barch(A^e)=\Barch(A\chotimes A^{\op})\xrightarrow{\sim}\Barch(A)\widehat\otimes \Barch(A^{\op}).
\]
Now apply Theorem~\ref{thm:opposite} to the second factor:
\[
\Barch(A^{\op})\xrightarrow{\sim}\Barch(A)^{\cop}.
\]
This proves the coalgebra statement. The algebra statement follows after finite-type linear duality.
\end{proof}

\section{The two-sided twisted bicomodule and the tangent bar construction}
\label{sec:tangent}

The construction of the correct coalgebra-side avatar of an $A$-bimodule. There are two a priori
different constructions:
the first is obtained by differentiating the ordered bar construction at a square-zero extension;
the second is obtained by twisting a two-sided tensor product with the universal twisting cochain.
A central point of this chapter is that these two constructions are the same.

\begin{definition}[Two-sided twisted bicomodule]
\label{def:Kbi}
Let $A$ be strongly admissible, let $C=\Barch(A)$, and let $M$ be an $A$-bimodule.
Define the \emph{two-sided twisted bicomodule}
\[
\KK_A^{\mathrm{bi}}(M):=C\widehat\otimes M\widehat\otimes C
\]
with differential
\[
d_{\KK}
=
d_C\otimes 1\otimes 1 + 1\otimes d_M\otimes 1 + 1\otimes 1\otimes d_C + d_L+d_R,
\]
where
\[
d_L(c\otimes m\otimes d)
=
\sum (-1)^{|c_{(1)}|}\, c_{(1)}\otimes \tau(c_{(2)})\cdot m\otimes d,
\]
\[
d_R(c\otimes m\otimes d)
=
\sum (-1)^{|c|+|m|}\, c\otimes m\cdot \tau(d_{(1)})\otimes d_{(2)},
\]
with the usual Koszul signs determined by the grading conventions.
The left and right $C$-coactions are induced by comultiplication on the two outer factors.
\end{definition}

\begin{lemma}[\ClaimStatusProvedHere]
\label{lem:Kbi-dg}
$\KK_A^{\mathrm{bi}}(M)$ is a dg $C$-bicomodule.
\end{lemma}

\begin{proof}
The untwisted part of the differential squares to zero. The quadratic terms in $d_L$ and in $d_R$
vanish by the Maurer--Cartan equation
\[
d\tau+\tau\star\tau=0.
\]
The mixed commutator $d_Ld_R+d_Rd_L$ vanishes because the left and right $A$-actions on the
bimodule $M$ commute. Compatibility with the bicomodule structure is immediate from coassociativity
of the coproduct on the outer $C$-factors.
\end{proof}

\begin{definition}[Square-zero extension and one-defect bar complex]
Let $M$ be an $A$-bimodule and form the square-zero extension
\[
A_\eps:=A\oplus \eps M,\qquad \eps^2=0.
\]
Inside the reduced ordered bar complex $\Barch(A_\eps)$, let $T_A^{\mathrm{bar}}(M)$ be the
subcomplex spanned by bar words containing exactly one $\eps M$-slot.
We call this the \emph{one-defect ordered bar complex}.
\end{definition}

\begin{proposition}[\ClaimStatusProvedHere]
\label{prop:one-defect}
Modulo $\eps^2$, there is a canonical decomposition
\[
\Barch(A_\eps)\cong \Barch(A)\oplus \eps\, T_A^{\mathrm{bar}}(M).
\]
The complex $T_A^{\mathrm{bar}}(M)$ is naturally a $C$-bicomodule.
\end{proposition}

\begin{proof}
A reduced bar word in $\Barch(A_\eps)$ contains some number of $\eps M$-slots.
Since $\eps^2=0$, any word with at least two such slots vanishes modulo $\eps^2$.
Therefore only the zero-defect and one-defect pieces survive.

The bar differential merges adjacent slots. Multiplying two ordinary $A$-slots remains in $A$.
Multiplying an $A$-slot with the unique $M$-slot remains in $M$ by the bimodule structure.
There is no nonzero product of two defect slots. Hence the one-defect sector is a subcomplex.
Deconcatenation cuts a one-defect word into a defect-free part and a one-defect part, or vice versa,
so the one-defect sector is naturally a bicomodule over the ordinary bar coalgebra $C=\Barch(A)$.
\end{proof}

\begin{theorem}[Tangent identification; \ClaimStatusProvedHere]
\label{thm:tangent=K}
There is a canonical isomorphism of dg $C$-bicomodules
\[
T_A^{\mathrm{bar}}(M)\xrightarrow{\sim}\KK_A^{\mathrm{bi}}(M).
\]
Equivalently, the one-defect ordered bar complex of a square-zero extension is exactly the two-sided
twisted bicomodule of the defect bimodule.
\end{theorem}

\begin{proof}
As a graded object,
\[
C=\bigoplus_{p\geq 0} (\susp^{-1}\bar A)^{\otimes p},
\]
hence
\[
C\widehat\otimes M\widehat\otimes C
\cong
\bigoplus_{p,q\geq 0}
(\susp^{-1}\bar A)^{\otimes p}\otimes M\otimes(\susp^{-1}\bar A)^{\otimes q}.
\]
This is in tautological bijection with one-defect bar words
\[
[a_1|\cdots|a_p|m|a_{p+1}|\cdots|a_{p+q}]
\]
in which the unique defect sits between a left bar segment of length $p$ and a right bar segment of
length $q$.

Under this identification, the ordinary bar differential on the left segment and the right segment
matches the internal bar differential coming from the two outer copies of $C$, while the two bar
faces adjacent to the defect slot match exactly the two twisting terms $d_L$ and $d_R$.
Therefore the differentials agree. Deconcatenation also agrees on the nose, because cutting a
one-defect word is the same as cutting the left and right bar segments around the unique defect.
\end{proof}

\begin{corollary}[Infinitesimal dual coalgebra; \ClaimStatusProvedHere]
\label{cor:infdual}
Modulo $\eps^2$, the bar coalgebra of the square-zero extension is
\[
\Barch(A_\eps)\cong C\oplus \eps\,\KK_A^{\mathrm{bi}}(M).
\]
If $A$ is finite-type Koszul, then after linear duality
\[
(A_\eps)^!\cong A^!\oplus \eps\, M^!
\qquad (\mathrm{mod}\ \eps^2),
\]
where
\[
M^!:=\Omega_C^{\mathrm{bi}}\bigl(\KK_A^{\mathrm{bi}}(M)\bigr).
\]
\end{corollary}

\begin{proof}
Combine Proposition~\ref{prop:one-defect} and Theorem~\ref{thm:tangent=K}. The statement after
dualization is the corresponding first-order cobar reconstruction.
\end{proof}

\begin{proposition}[Infinitesimal annular variation; \ClaimStatusProvedHere]
\label{prop:infann}
Modulo $\eps^2$ one has
\[
\HH^{\mathrm{ch}}_\bullet(A_\eps)
\cong
\HH^{\mathrm{ch}}_\bullet(A)\oplus \eps\, \HH^{\mathrm{ch}}_\bullet(A,M).
\]
On the coalgebra side this becomes
\[
\coHoch^{\mathrm{ch}}_\bullet(\Barch(A_\eps))
\cong
\coHoch^{\mathrm{ch}}_\bullet(C)\oplus \eps\,
\coHoch^{\mathrm{ch}}_\bullet(C,\KK_A^{\mathrm{bi}}(M)).
\]
\end{proposition}

\begin{proof}
The Hochschild complex of the square-zero extension decomposes, modulo $\eps^2$, into the sector with
no defect and the sector with exactly one defect; the second computes Hochschild homology with
coefficients in $M$. The coalgebra statement follows from
Corollary~\ref{cor:infdual} and the coHochschild identification proved later in
Theorem~\ref{thm:HH-coHH-homology}.
\end{proof}

\section{PBW-complete two-sided Koszul duality}
\label{sec:bimod-bicomod}

The repository already establishes one-sided module/comodule duality. Bootstrapping it to the
two-sided theory by passing from $A$ to its enveloping algebra $A^e$ and then using
Corollary~\ref{cor:enveloping}.

\begin{definition}[PBW-complete and coderived bicomodule categories]
Let $A$ be strongly admissible and set $C=\Barch(A)$.
\begin{enumerate}[label=(\alph*)]
\item $\Dpbw_{\mathrm{bi}}(A)$ denotes the derived category of PBW-complete $A$-bimodules obtained
from complete semifree $A^e$-resolutions.

\item $\Dco_{\mathrm{bi}}(C)$ denotes the coderived category of conilpotent $C$-bicomodules.
Equivalently, by Lemma~\ref{lem:bicom-e}, it is the coderived category of left $C^e$-comodules.
\end{enumerate}
\end{definition}

\begin{definition}[Two-sided twisted cobar]
Let $N$ be a $C$-bicomodule. Define
\[
\Omega_C^{\mathrm{bi}}(N):=A\widehat\otimes N\widehat\otimes A
\]
with the differential dual to the one in Definition~\ref{def:Kbi}.
\end{definition}

\begin{theorem}[PBW-complete bimodule/bicomodule equivalence; \ClaimStatusProvedHere]
\label{thm:bimod-bicomod}
Let $A$ be strongly admissible, and assume $(A,A^{\op})$ is EZ-admissible.
Set $C=\Barch(A)$.
Then there is an exact equivalence
\[
\KK_A^{\mathrm{bi}}:\Dpbw_{\mathrm{bi}}(A)\xrightarrow{\sim}\Dco_{\mathrm{bi}}(C)
\]
with inverse $\Omega_C^{\mathrm{bi}}$. In the curved regime the same statement holds after replacing
the ordinary derived category on the algebra side by the corresponding completed
contra\-derived/coderived version.
\end{theorem}

\begin{proof}
By the enveloping algebra formalism, $A$-bimodules are left $A^e$-modules.
By Corollary~\ref{cor:enveloping},
\[
\Barch(A^e)\xrightarrow{\sim} C^e.
\]
Apply the one-sided module/comodule Koszul duality from Hypothesis~\ref{hyp:inputs}(II) to the
associative chiral algebra $A^e$. This gives an equivalence between complete $A^e$-modules and
conilpotent $C^e$-comodules. Since left $C^e$-comodules are the same as $C$-bicomodules by
Lemma~\ref{lem:bicom-e}, we obtain the asserted equivalence.

The functor is precisely the two-sided twisted bar construction
\[
M\longmapsto C\widehat\otimes M\widehat\otimes C,
\]
because that is the one-sided bar functor for the algebra $A^e$ rewritten using the enveloping
identification. The inverse is the corresponding two-sided cobar functor.
\end{proof}

\begin{theorem}[Diagonal correspondence; \ClaimStatusProvedHere]
\label{thm:diagonal}
Under the equivalence of Theorem~\ref{thm:bimod-bicomod}, the diagonal bimodule ${}_A A_A$ is sent
to the diagonal bicomodule $C_\Delta$. Equivalently,
\[
\KK_A^{\mathrm{bi}}(A)\simeq C_\Delta
\]
in $\Dco_{\mathrm{bi}}(C)$.
\end{theorem}

\begin{proof}
By definition,
\[
\KK_A^{\mathrm{bi}}(A)=C\widehat\otimes A\widehat\otimes C
\]
with the two-sided twisting differential.

First contract the left half:
the standard bar--cobar counit resolution yields a quasi-isomorphism
\[
C\widehat\otimes_\tau A\xrightarrow{\sim} \mathbf 1
\]
compatible with the right $C$-coaction. Hence
\[
\KK_A^{\mathrm{bi}}(A)\simeq C
\]
as a right $C$-comodule.

Similarly, contracting the right half using
\[
A\widehat\otimes_\tau C\xrightarrow{\sim}\mathbf 1
\]
gives
\[
\KK_A^{\mathrm{bi}}(A)\simeq C
\]
as a left $C$-comodule.

To identify the opposite coactions, filter $\KK_A^{\mathrm{bi}}(A)$ by total coradical degree in the
two outer copies of $C$. On the associated graded the twisting terms disappear, and the bicomodule
structure is literally the left and right coproducts on the two outer factors. Under either of the
two contractions above, the associated graded collapses to a single copy of $C$, and the surviving
left and right coactions are exactly the two legs of $\Delta_C$. Because the filtration is complete
and bounded below in the strongly admissible regime, this associated-graded identification lifts to
the filtered complex. Therefore the resulting bicomodule is $C_\Delta$.
\end{proof}

\begin{corollary}[The diagonal is the unit for composition; \ClaimStatusProvedHere]
\label{cor:unit}
For every $X\in \Dco_{\mathrm{bi}}(C)$ there are canonical isomorphisms
\[
C_\Delta \square_C^{\mathbf R} X \simeq X \simeq X \square_C^{\mathbf R} C_\Delta.
\]
\end{corollary}

\begin{proof}
Transport the tautologies
\[
A\otimes_A^{\mathbf L} M\simeq M\simeq M\otimes_A^{\mathbf L} A
\]
for $A$-bimodules along Theorem~\ref{thm:bimod-bicomod} and use
Theorem~\ref{thm:diagonal}.
\end{proof}

\begin{corollary}[Tensor--cotensor gluing; \ClaimStatusProvedHere]
\label{cor:tensor-cotensor}
For all $M,N\in \Dpbw_{\mathrm{bi}}(A)$ there is a canonical isomorphism
\[
\KK_A^{\mathrm{bi}}(M\otimes_A^{\mathbf L} N)
\simeq
\KK_A^{\mathrm{bi}}(M)\square_C^{\mathbf R}\KK_A^{\mathrm{bi}}(N)
\]
in $\Dco_{\mathrm{bi}}(C)$.
\end{corollary}

\begin{proof}
Transport the derived tensor product across the equivalence of
Theorem~\ref{thm:bimod-bicomod}.
\end{proof}

\begin{remark}
Corollary~\ref{cor:tensor-cotensor} is the coalgebraic gluing law for intervals.
The diagonal bicomodule $C_\Delta$ is therefore the universal boundary condition for interval
composition.
\end{remark}

\section{Hochschild--coHochschild duality}
\label{sec:HH-coHH}

The identification of annular traces and centers on the algebra side with coHochschild theories of the bar
coalgebra.

\begin{theorem}[Associative chiral Hochschild/coHochschild homology; \ClaimStatusProvedHere]
\label{thm:HH-coHH-homology}
Let $A$ be strongly admissible, let $(A,A^{\op})$ be EZ-admissible, and set $C=\Barch(A)$.
Then for every complete $A$-bimodule $M$ there is a canonical isomorphism
\[
\HH^{\mathrm{ch}}_\bullet(A,M)\simeq
\coHoch^{\mathrm{ch}}_\bullet(C,\KK_A^{\mathrm{bi}}(M)).
\]
In particular,
\[
\HH^{\mathrm{ch}}_\bullet(A)\simeq \coHoch^{\mathrm{ch}}_\bullet(C)
\simeq
C_\Delta \square_{C^e}^{\mathbf R} C_\Delta.
\]
\end{theorem}

\begin{proof}
By definition,
\[
\HH^{\mathrm{ch}}_\bullet(A,M)=\Tor_\bullet^{A^e}(A,M).
\]
Apply the equivalence of Theorem~\ref{thm:bimod-bicomod} to the pair $(A,M)$.
The diagonal bimodule $A$ goes to $C_\Delta$ by Theorem~\ref{thm:diagonal}, and $M$ goes to
$\KK_A^{\mathrm{bi}}(M)$ by definition. Derived tensor over $A^e$ is transported to derived cotensor
over $C^e$. Therefore
\[
\Tor_\bullet^{A^e}(A,M)\simeq \Cotor_\bullet^{C^e}(C_\Delta,\KK_A^{\mathrm{bi}}(M)),
\]
which is the stated formula. Taking $M=A$ yields the second assertion.
\end{proof}

\begin{theorem}[Associative chiral Hochschild/coHochschild cohomology; \ClaimStatusProvedHere]
\label{thm:HH-coHH-cohomology}
Under the same hypotheses, for every complete $A$-bimodule $M$ there is a canonical isomorphism
\[
\HH_{\mathrm{ch}}^\bullet(A,M)\simeq
\coHoch_{\mathrm{ch}}^\bullet(C,\KK_A^{\mathrm{bi}}(M)).
\]
In particular,
\[
\HH_{\mathrm{ch}}^\bullet(A)\simeq
\coHoch_{\mathrm{ch}}^\bullet(C)
=
\Coext_{C^e}^\bullet(C_\Delta,C_\Delta).
\]
\end{theorem}

\begin{proof}
Derived Hom is preserved by the equivalence of Theorem~\ref{thm:bimod-bicomod}. Therefore
\[
\RHom_{A^e}(A,M)\simeq \RHom_{C^e}(C_\Delta,\KK_A^{\mathrm{bi}}(M)),
\]
which is the desired coHochschild cohomology complex. The specialization to $M=A$ follows from the
diagonal correspondence.
\end{proof}

\begin{corollary}[The annulus as self-cotrace; \ClaimStatusProvedHere]
\label{cor:annulus}
The annular closed object attached to the universal boundary condition $C_\Delta$ is the coHochschild
object
\[
\mathcal H_C^{\mathrm{cl}}
:=
C_\Delta \square_{C^e}^{\mathbf R} C_\Delta
\simeq
\coHoch^{\mathrm{ch}}_\bullet(C)
\simeq
\HH^{\mathrm{ch}}_\bullet(A).
\]
\end{corollary}

\begin{proof}
This is the special case $M=A$ of Theorem~\ref{thm:HH-coHH-homology}.
\end{proof}

\begin{corollary}[Cap action; \ClaimStatusProvedHere]
\label{cor:cap}
There is a canonical action
\[
\coHoch^{\mathrm{ch}}_\bullet(C)\square_C^{\mathbf R} C_\Delta \longrightarrow C_\Delta
\]
obtained by transporting the Hochschild cap action
\[
\HH^{\mathrm{ch}}_\bullet(A)\otimes A\longrightarrow A
\]
across the equivalence of Theorem~\ref{thm:bimod-bicomod}.
\end{corollary}

\begin{proof}
Apply Theorem~\ref{thm:bimod-bicomod} to the classical cap action viewed as a map of bimodules.
\end{proof}

\begin{remark}
The cohomology theorem is the center theorem. The homology theorem is the annulus theorem. Together
they say that the algebraic center and algebraic trace of $A$ can be read directly from the
coalgebraic diagonal $C_\Delta$.
\end{remark}

\section{The diagonal bicomodule as universal boundary condition}
\label{sec:open-trace}

There are two operations one wants on the boundary side.

First, interval composition, already realized by derived cotensor and controlled by
Corollary~\ref{cor:tensor-cotensor}. Second, the ordered pair-of-pants operation, which is not
cotensor but rather the transport of the external tensor product of bimodules. This second
operation is inherently ordered; in general we do \emph{not} claim it is symmetric without extra
structure.

\begin{definition}[Ordered fusion product]
Let $X,Y\in \Dco_{\mathrm{bi}}(C)$. Define
\[
X\star Y
:=
\KK_A^{\mathrm{bi}}\Bigl(
\Omega_C^{\mathrm{bi}}(X)\widehat\otimes \Omega_C^{\mathrm{bi}}(Y)
\Bigr),
\]
where $\widehat\otimes$ is the completed tensor product of bimodules with outer $A$-actions:
the left $A$-action is on the first factor and the right $A$-action is on the second factor.
We call $\star$ the \emph{ordered fusion product}.
\end{definition}

\begin{remark}
The adjective ``ordered'' is essential. Unless further symmetry data are supplied, $\star$ is only
a monoidal product reflecting the ordering of incoming open boundaries. It is the algebraic shadow
of an ordered pair-of-pants, not yet of the fully symmetric open cobordism category.
\end{remark}

\begin{theorem}[Ordered pair-of-pants algebra; \ClaimStatusProvedHere]
\label{thm:pair-of-pants}
The diagonal bicomodule $C_\Delta$ carries a canonical associative algebra structure with respect to
the ordered fusion product:
\[
\mu_P:C_\Delta\star C_\Delta\longrightarrow C_\Delta,
\qquad
\eta_P:\mathbf 1_\star\longrightarrow C_\Delta,
\]
where $\mathbf 1_\star:=\KK_A^{\mathrm{bi}}(k)$ is the image of the trivial bimodule.
\end{theorem}

\begin{proof}
By Theorem~\ref{thm:diagonal}, $C_\Delta\simeq \KK_A^{\mathrm{bi}}(A)$. The multiplication and unit
of the algebra $A$,
\[
\mu_A:A\widehat\otimes A\to A,\qquad \eta_A:k\to A,
\]
transport across $\KK_A^{\mathrm{bi}}$ to maps
\[
\mu_P:C_\Delta\star C_\Delta\to C_\Delta,\qquad \eta_P:\mathbf 1_\star\to C_\Delta.
\]
Associativity and unitality follow immediately from associativity and unitality of $A$.
\end{proof}

\begin{remark}[Bar chain model on the stratified pair-of-pants]%
\label{rem:stratified-pants}%
The pair-of-pants $P=\mathbb{P}^1\setminus\{0,1,\infty\}$ is the three-punctured sphere
equipped with three boundary circles $\partial_{\mathrm{in}}^1$, $\partial_{\mathrm{in}}^2$,
$\partial_{\mathrm{out}}$.
The ordered bar complex on $P$ is built from configurations of points on $P$ valued in the
ordered Fulton--MacPherson compactification
$\overline{\mathrm{FM}}_n^{\mathrm{ord}}(P)$, whose boundary strata record collisions
at each of the three punctures.
The pair-of-pants multiplication
\[
\mu_P\colon C_\Delta\star C_\Delta\longrightarrow C_\Delta
\]
of Theorem~\ref{thm:pair-of-pants} is then the chain-level avatar of the chiral product:
on homology it descends to the product on $\HH_\bullet^{\mathrm{ch}}(\cA)$, while at
cochain level it retains the full $A_\infty$ data encoded by the collision strata.
Concretely, $P$ is the multiplication cobordism: the two input circles carry the
annular bar complexes $B_\bullet^{\mathrm{ann}}(\cA)$, and $\mu_P$ composes them
through the pair-of-pants into the annular bar complex on the output circle.
The associativity of $\mu_P$ is the statement that the two decompositions of
the genus-zero four-punctured sphere into pairs of pants give chain-homotopic compositions.
\end{remark}

\begin{definition}[Ordered open trace formalism]
An \emph{ordered genus-zero open trace formalism} consists of:
\begin{enumerate}[label=(\roman*)]
\item a composition product $\square_C^{\mathbf R}$ with unit object $C_\Delta$,
\item an ordered fusion product $\star$,
\item an associative algebra structure $(C_\Delta,\mu_P,\eta_P)$ for $\star$,
\item a closed object $\mathcal H_C^{\mathrm{cl}}$ obtained by annular closure,
\item an annulus action of $\mathcal H_C^{\mathrm{cl}}$ on $C_\Delta$,
\end{enumerate}
subject to the sewing relations induced from the algebra side.
\end{definition}

\begin{theorem}[Ordered genus-zero open trace formalism; \ClaimStatusProvedHere]
\label{thm:ordered-open}
The data
\[
\bigl(\Dco_{\mathrm{bi}}(C),\square_C^{\mathbf R},C_\Delta,\star,\mu_P,\eta_P,\mathcal H_C^{\mathrm{cl}}\bigr)
\]
constitute an ordered genus-zero open trace formalism. More precisely:
\begin{enumerate}[label=(\alph*)]
\item $C_\Delta$ is the unit for interval composition $\square_C^{\mathbf R}$.
\item $(C_\Delta,\mu_P,\eta_P)$ is an associative algebra object for ordered fusion.
\item Annular closure is
\[
\mathcal H_C^{\mathrm{cl}}\simeq \coHoch^{\mathrm{ch}}_\bullet(C).
\]
\item The annulus acts on the boundary object by the action of
Corollary~\ref{cor:cap}.
\item All ordered genus-zero sewing identities are transported from the corresponding identities for
$A$-bimodules.
\end{enumerate}
\end{theorem}

\begin{proof}
Statement (a) is Corollary~\ref{cor:unit}. Statement (b) is
Theorem~\ref{thm:pair-of-pants}. Statement (c) is
Corollary~\ref{cor:annulus}. Statement (d) is Corollary~\ref{cor:cap}.
Finally, every ordered genus-zero sewing relation on the coalgebra side is the image under the
equivalence $\KK_A^{\mathrm{bi}}$ of the corresponding relation on the algebra side.
\end{proof}

\begin{remark}[What is proved and what is not]
The theorem gives the algebraic ordered open package. To upgrade it to a fully symmetric
one-brane open TCFT one needs additional structure that symmetrizes ordered fusion, for example a
boundary symmetry datum or a direct geometric realization by bordered configuration spaces.
This is formulated below as part of the conjectural frontier.
\end{remark}

\subsection{Calabi--Yau transport}

The ordered open trace formalism acquires a Frobenius structure under a Calabi--Yau hypothesis on
$A$.

\begin{definition}[Complete chiral Calabi--Yau structure]
A strongly admissible associative chiral algebra $A$ is called \emph{$d$-Calabi--Yau} if:
\begin{enumerate}[label=(\alph*)]
\item $A$ is perfect as an $A^e$-module, and
\item there is an isomorphism in $\Dpbw_{\mathrm{bi}}(A)$
\[
\varphi_A:A\xrightarrow{\sim}\RHom_{A^e}(A,A^e)[d].
\]
\end{enumerate}
Equivalently, $A$ carries a nondegenerate $d$-shifted derived trace pairing.
\end{definition}

\begin{theorem}[Shifted ordered Frobenius structure; \ClaimStatusProvedHere]
\label{thm:CY}
Assume $A$ is $d$-Calabi--Yau. Then the underlying complex of $C_\Delta$ carries a canonical
$d$-shifted ordered Frobenius structure. More precisely:
\begin{enumerate}[label=(\roman*)]
\item the ordered pair-of-pants multiplication $\mu_P$ and unit $\eta_P$ are as above;
\item there is a trace
\[
\eps_P:C_\Delta\to k[-d]
\]
transported from the Calabi--Yau trace on $A$;
\item there is a coproduct
\[
\Delta_P:C_\Delta\to C_\Delta\star C_\Delta[d]
\]
adjoint to $\mu_P$ under $\eps_P$;
\item the ordered Frobenius identities
\[
(\mu_P\star \id)\circ(\id\star \Delta_P)
=
\Delta_P\circ \mu_P
=
(\id\star\mu_P)\circ(\Delta_P\star \id)
\]
hold.
\end{enumerate}
\end{theorem}

\begin{proof}
On the algebra side, the Calabi--Yau structure identifies $A$ with its bimodule dual shifted by
$d$. Therefore multiplication
\[
\mu_A:A\widehat\otimes A\to A
\]
admits a unique adjoint coproduct
\[
\Delta_A:A\to A\widehat\otimes A[d]
\]
with respect to the trace pairing $\eps_A:A\to k[-d]$; explicitly one may define $\Delta_A$ as the
adjoint of $\mu_A$ under $\varphi_A$. The ordered Frobenius identities are formal consequences of
associativity of $\mu_A$ and the nondegeneracy of the pairing.

Transport the entire package $(\mu_A,\eta_A,\eps_A,\Delta_A)$ across the equivalence
$\KK_A^{\mathrm{bi}}$ and use Theorem~\ref{thm:diagonal}. The resulting structure maps are
$\mu_P,\eta_P,\eps_P,\Delta_P$, and all identities are preserved because the equivalence preserves
composition.
\end{proof}

\begin{corollary}[Cardy operator on the coalgebra side; \ClaimStatusProvedHere]
\label{cor:cardy}
Under the hypotheses of Theorem~\ref{thm:CY}, the ordered annulus endomorphism of the universal
boundary object is
\[
\mathsf{Ann}_{C_\Delta}:=\mu_P\circ \Delta_P\in \operatorname{End}(C_\Delta)[d].
\]
Its annular closure factors through the closed object
\[
\mathcal H_C^{\mathrm{cl}}\simeq \coHoch^{\mathrm{ch}}_\bullet(C).
\]
\end{corollary}

\begin{proof}
Transport the usual annulus endomorphism $\mu_A\circ \Delta_A$ of the Calabi--Yau algebra $A$ across
$\KK_A^{\mathrm{bi}}$ and use Corollary~\ref{cor:annulus}.
\end{proof}

\section{A master theorem and its conceptual form}
\label{sec:master}

For repository integration it is useful to package the preceding results into a single statement.

\begin{theorem}[Master theorem; \ClaimStatusProvedHere]
\label{thm:master}
Let $A$ be a strongly admissible augmented associative chiral algebra on $X$, assume
$(A,A^{\op})$ is EZ-admissible, and set
\[
C:=\Barch(A).
\]
Then:
\begin{enumerate}[label=(\arabic*)]
\item \textbf{Ordered shuffle.}
There is a quasi-isomorphism of conilpotent dg coalgebras
\[
\Barch(A\chotimes B)\simeq \Barch(A)\widehat\otimes \Barch(B)
\]
for every EZ-admissible partner $B$.

\item \textbf{Opposite-duality.}
There is a canonical isomorphism
\[
\Barch(A^{\op})\simeq \Barch(A)^{\cop}.
\]

\item \textbf{Enveloping duality.}
There is a canonical quasi-isomorphism
\[
\Barch(A^e)\simeq C^e.
\]

\item \textbf{Two-sided bimodule/bicomodule duality.}
There is an exact equivalence
\[
\Dpbw_{\mathrm{bi}}(A)\simeq \Dco_{\mathrm{bi}}(C).
\]

\item \textbf{Diagonal correspondence.}
Under this equivalence one has
\[
{}_A A_A \longleftrightarrow C_\Delta.
\]

\item \textbf{Composition.}
The diagonal bicomodule is the unit for derived cotensor:
\[
C_\Delta\square_C^{\mathbf R} X\simeq X\simeq X\square_C^{\mathbf R}C_\Delta.
\]

\item \textbf{Hochschild/coHochschild homology.}
For every complete $A$-bimodule $M$,
\[
\HH^{\mathrm{ch}}_\bullet(A,M)
\simeq
\coHoch^{\mathrm{ch}}_\bullet(C,\KK_A^{\mathrm{bi}}(M)).
\]

\item \textbf{Hochschild/coHochschild cohomology.}
For every complete $A$-bimodule $M$,
\[
\HH_{\mathrm{ch}}^\bullet(A,M)
\simeq
\coHoch_{\mathrm{ch}}^\bullet(C,\KK_A^{\mathrm{bi}}(M)).
\]

\item \textbf{Tangent theory.}
For a square-zero extension $A_\eps=A\oplus \eps M$,
\[
\Barch(A_\eps)\cong C\oplus \eps\,\KK_A^{\mathrm{bi}}(M)\qquad (\mathrm{mod}\ \eps^2).
\]

\item \textbf{Ordered open trace formalism.}
The coalgebraic boundary object $C_\Delta$ carries ordered pair-of-pants, annulus, and composition
operations transported from the algebra side, and if $A$ is $d$-Calabi--Yau then the underlying
complex of $C_\Delta$ becomes a $d$-shifted ordered Frobenius object.
\end{enumerate}
\end{theorem}

\begin{proof}
This is a summary of Theorems~\ref{thm:shuffle}, \ref{thm:opposite},
\ref{thm:bimod-bicomod}, \ref{thm:diagonal}, \ref{thm:HH-coHH-homology},
\ref{thm:HH-coHH-cohomology}, \ref{thm:tangent=K}, and \ref{thm:ordered-open}, together with
Corollaries~\ref{cor:enveloping}, \ref{cor:unit}, and \ref{cor:infdual}, and
Theorem~\ref{thm:CY}.
\end{proof}

\begin{remark}[The Chriss--Ginzburg architecture]
\label{rem:chriss-ginzburg-architecture}
The master theorem describes the \emph{algebra side} of associative chiral
Koszul duality: the diagonal bicomodule $C_\Delta$, its composition
algebra, and the Hochschild--coHochschild dictionary. The explicit
constructions of Appendix~\ref{sec:ordered-bar-explicit} supply the
\emph{spectral side}: the $R$-matrix as descent datum
(Construction~\ref{constr:r-matrix-covering}), the ordered-bar
realization of the open-colour dual with its Kac--Moody Yangian
specialization
(Construction~\ref{constr:dg-shifted-yangian-from-bar}), and the
Yang--Baxter equation as $d^2=0$ on the degree-$3$ complex
(Proposition~\ref{prop:ybe-from-d-squared}).

The unifying principle is the Chriss--Ginzburg paradigm: every algebraic
structure is a Maurer--Cartan element in a convolution dg~Lie algebra.
The $R$-matrix $R(z)$ is the MC element in the Gerstenhaber complex
$C^\bullet_{\mathrm{GS}}(\cA,\cA)$ that governs $\Sigma_n$-equivariant
descent (Remark~\ref{rem:r-matrix-mc}); the Yang--Baxter equation is the
MC equation; the classical $r$-matrix $r_0$ is its linearisation
(Corollary~\ref{cor:classical-ybe}); and the Yangian $Y_\hbar(\mathfrak g)$
is the algebra of symmetries of this MC datum. The passage from the master
theorem's algebraic core to the spectral data of ordered
configuration spaces is the passage from the abstract MC moduli to its
explicit coordinatisation.
\end{remark}

\section{The \texorpdfstring{$E_1$}{E1} five main theorems at genus zero}
\label{sec:e1-five-theorems-genus0}
\index{five main theorems!E1@$E_1$ variant!genus zero|textbf}

Volume~I establishes five main theorems for $E_\infty$-chiral algebras: bar-cobar adjunction, Koszul inversion, complementarity, leading coefficient, and the Hochschild ring. Each has an $E_1$-chiral counterpart, obtained by replacing symmetric coinvariants with the ordered bar complex and its braid-group equivariance. The five $E_1$ theorems at genus zero are formal consequences of the master theorem (Theorem~\ref{thm:master}) together with the explicit constructions of Appendix~\ref{sec:ordered-bar-explicit}; they are collected here to make the parallel with Vol~I explicit.

\begin{theorem}[Theorem~$\mathrm{A}^{E_1}$; \ClaimStatusProvedHere]
\label{thm:e1-theorem-A}
\index{five main theorems!E1@$E_1$ variant!Theorem A}
$\Cobar \dashv \Barch\colon
\mathsf{Alg}^{E_1}_{\mathrm{ch}} \rightleftarrows
\mathsf{Coalg}^{E_1}_{\mathrm{ch}}$.
Koszul exchange acts by $R^{-1}(z)$
\textup{(}Theorem~\textup{\ref{thm:opposite}}).
\end{theorem}

\begin{proof}
Ordered bar--cobar equivalence (\S\ref{sec:setup}) plus
opposite-duality (Theorem~\ref{thm:opposite}).
\end{proof}

\begin{theorem}[Theorem~$\mathrm{B}^{E_1}$; \ClaimStatusProvedHere]
\label{thm:e1-theorem-B}
\index{five main theorems!E1@$E_1$ variant!Theorem B}
On the $E_1$ Koszul locus,
$\Cobar(\Barch(\cA)) \xrightarrow{\sim} \cA$.
The $E_1$ Koszul dual is
$\cA^{!,E_1} := H^*(\Barch(\cA))^\vee$.
For $\cA = \hat\fg_k$:
$\cA^{!,E_1} = Y(\fg)$.
\end{theorem}

\begin{proof}
PBW-completeness (\S\ref{sec:setup}) and the dg-shifted Yangian
identification
(Construction~\ref{constr:dg-shifted-yangian-from-bar}).
\end{proof}

\begin{theorem}[Theorem~$\mathrm{C}^{E_1}$; \ClaimStatusProvedHere]
\label{thm:e1-theorem-C}
\index{five main theorems!E1@$E_1$ variant!Theorem C}
$\HH^{\mathrm{ch}}_\bullet(\cA)
\simeq
C_\Delta \square_{C^e}^{\mathbf{R}} C_\Delta$,
where $C = \Barch(\cA)$.
If $\cA$ is $d$-Calabi--Yau, then $C_\Delta$ is
$d$-shifted ordered Frobenius
\textup{(}Theorem~\textup{\ref{thm:CY}}).
\end{theorem}

\begin{proof}
Theorems~\ref{thm:HH-coHH-homology}, \ref{thm:ordered-open},
\ref{thm:CY}.
\end{proof}

\begin{theorem}[Theorem~$\mathrm{D}^{E_1}$; \ClaimStatusProvedHere]
\label{thm:e1-theorem-D}
\index{five main theorems!E1@$E_1$ variant!Theorem D}
\index{R-matrix!as E1 modular characteristic@as $E_1$ modular characteristic}
The classical $r$-matrix
$r(z) = \operatorname{Res}^{\mathrm{coll}}_{0,2}(\Theta_\cA^{E_1})
\in \operatorname{End}(V \otimes V)(\!(z^{-1})\!)$
satisfies the CYBE
\textup{(}Proposition~\textup{\ref{prop:ybe-from-d-squared}}).
Under averaging,
$\operatorname{av}(r(z)) = \kappa(\cA)$.
\end{theorem}

\begin{proof}
Construction~\ref{constr:r-matrix-monodromy},
Proposition~\ref{prop:r-matrix-descent},
Proposition~\ref{prop:ybe-from-d-squared}.
\end{proof}

\begin{theorem}[Theorem~$\mathrm{H}^{E_1}$; \ClaimStatusProvedHere]
\label{thm:e1-theorem-H}
\index{five main theorems!E1@$E_1$ variant!Theorem H}
For every complete $\cA$-bimodule $M$\textup{:}
$\HH^{\mathrm{ch}}_\bullet(\cA, M)
\simeq
\coHoch^{\mathrm{ch}}_\bullet(C, \KK_\cA^{\mathrm{bi}}(M))$,\quad
$\HH_{\mathrm{ch}}^\bullet(\cA, M)
\simeq
\coHoch_{\mathrm{ch}}^\bullet(C, \KK_\cA^{\mathrm{bi}}(M))$.
\end{theorem}

\begin{proof}
Theorems~\ref{thm:HH-coHH-homology}
and~\ref{thm:HH-coHH-cohomology}.
\end{proof}

\begin{remark}[$E_1$ vs $E_\infty$]
\label{rem:e1-einfty-five-comparison}
\index{five main theorems!E1 vs E-infty comparison@$E_1$ vs $E_\infty$ comparison}
\begin{center}
\renewcommand{\arraystretch}{1.15}
\begin{tabular}{lcc}
\toprule
& \emph{$E_\infty$} & \emph{$E_1$} \\
\midrule
\textbf{A}: Exchange
 & $\Com^! = \Lie$
 & $\Ass^! = \Ass \otimes \mathrm{sgn}$ \\
\textbf{B}: Dual
 & $\cA^!$ (Verdier)
 & $Y(\fg)$ ($R$-matrix) \\
\textbf{C}: Complementarity
 & $Q_g(\cA) + Q_g(\cA^!)$
 & $C_\Delta \square^{\mathbf{R}} C_\Delta$ \\
\textbf{D}: Characteristic
 & $\kappa(\cA)$ (scalar)
 & $r(z)$ ($R$-matrix) \\
\textbf{H}: Hochschild
 & $\ChirHoch$
 & HH $=$ coHH \\
\bottomrule
\end{tabular}
\end{center}
The $E_\infty$ theorem is the $\Sigma_n$-coinvariant of the $E_1$ theorem.
\end{remark}

\begin{remark}[The Drinfeld double programme]
\label{rem:drinfeld-double-programme}
The boundary algebras $\cA$ (Dirichlet) and $\cA^!$ (Neumann) are linked by the universal twisting morphism $\tau \in \cA^! \otimes \cA$ and by the Verdier intertwining $\mathbb{D}_{\Ran}(\barB(\cA)) \simeq \barB(\cA^!)$ (Theorem~A of Volume~I). The full Drinfeld double $U = \cA \bowtie \cA^!$, whose module category should be the line-operator category $\cC_{\mathrm{line}}$, is not yet constructed at the bar-complex level. The ingredients are present: the algebra structure (boundary operator composition), the coalgebra structure (deconcatenation coproduct on $B^{\mathrm{ord}}(\cA)$), the Hopf pairing (from complementarity / the cyclic structure), and the opposite-duality (Theorem~\ref{thm:opposite}, the $E_1$ analogue of the antipode). The assembly of these into a single Hopf algebra $U$ with $\cC_{\mathrm{line}} \simeq U\text{-mod}$ is the \emph{Drinfeld double programme}, connecting the bar-cobar framework to Dimofte's geometric Tannakian reconstruction \cite{Dimofte25}.
\end{remark}

\begin{remark}[Dimofte's $K$-matrix and shadow depth]
\label{rem:dimofte-k-matrix}
In the PIRSA lecture series \cite{Dimofte25}, Dimofte isolates a single chiral operator $K_\cA(z)$ which controls the deviation of the boundary coproduct from the primitive (Hopf) coproduct. Its explicit form is
\begin{equation}
\label{eq:dimofte-k-matrix}
K_\cA(z) \;=\; z^{\alpha_0}\,\exp\!\Biggl(\sum_{n>0} \frac{\alpha_{-n}}{-n}\, z^{n}\Biggr)\,\exp\!\Biggl(\sum_{n>0} \frac{\alpha_{n}}{-n}\, z^{-n}\Biggr),
\end{equation}
where $\{\alpha_n\}_{n\in\Z}$ are the modes of a free boson identified with the $\mathfrak{u}(1)^r$ Cartan factor of the boundary algebra $\cA$, and $\alpha_0$ is the zero-mode charge. The operator $K_\cA(z)$ is itself a vertex operator: it lives inside $\cA$, not over it.
The interpretation is diagnostic. For class~$\mathbf{G}$ (Heisenberg, free fermion, lattice vertex algebra) the boundary coproduct is already primitive and $K_\cA(z) = 1$; the boundary Hopf algebra is the Hopf primitive and no deviation is needed. For classes $\mathbf{L}$ and $\mathbf{C}$ (affine Kac--Moody at generic level, $\beta\gamma$) the deviation $K_\cA(z) \neq 1$ is polynomial in the mode expansion: finitely many negative powers of $z$ contribute and the corresponding Hopf corrections terminate. For class~$\mathbf{M}$ (Virasoro, $\mathcal{W}_N$) the deviation has an infinite expansion and the corrections never terminate. Dimofte reads this as a ``BPS decomposition'' of the boundary Hilbert space along the $\mathfrak{u}(1)^r$ direction: the free-boson factor is the generic stratum and the higher $\alpha_n$ modes record BPS corrections.
The sanity check is immediate: at $K_\cA(z) = 1$, equation~\eqref{eq:dimofte-k-matrix} reduces to $z^0 \cdot 1 \cdot 1 = 1$ (all non-zero modes $\alpha_n$ ($n\neq 0$) act trivially and $\alpha_0 = 0$), and the twisted coproduct returns to the primitive coproduct $\Delta(a) = a \otimes 1 + 1 \otimes a$ on the boundary generators, recovering the class~$\mathbf{G}$ behaviour.
The conceptual punchline matches the shadow-depth classification of Volume~I: the partition $\{\mathbf{G},\mathbf{L},\mathbf{C},\mathbf{M}\}$ by shadow depth $d_{\mathrm{alg}} \in \{0,1,2,\infty\}$ coincides with the partition by ``$K$-matrix complexity''.
\begin{center}
\renewcommand{\arraystretch}{1.2}
\begin{tabular}{lll}
\toprule
Class & $K_\cA(z)$ & $d_{\mathrm{alg}}$ \\
\midrule
$\mathbf{G}$ (Heisenberg, fermion, lattice) & $1$ & $0$ \\
$\mathbf{L}$ (affine Kac--Moody) & polynomial, simple pole & $1$ \\
$\mathbf{C}$ ($\beta\gamma$) & polynomial, double pole & $2$ \\
$\mathbf{M}$ (Virasoro, $\mathcal{W}_N$) & infinite expansion & $\infty$ \\
\bottomrule
\end{tabular}
\end{center}
Equivalently: $K_\cA(z) = 1$ iff the boundary coproduct is primitive iff the shadow tower terminates at depth zero. This is the $K$-matrix shadow of the pole-order hierarchy summarised in Table~\ref{tab:two-colour-koszul-duals} and in the shadow-depth discussion of Volume~I. It is not an $r$-matrix statement: the $K$-matrix modifies the coproduct, not the product, and the level-$k$ vanishing check does not apply directly; the corresponding $r$-matrix check is the one already recorded in the Volume~II affine constructions (classical $r$-matrix $k\,\Omega/z$ vanishes at $k=0$, in which case $\cA = \cH_0$ collapses to the trivial Heisenberg and $K_\cA(z) = 1$, consistent with class~$\mathbf{G}$).
\end{remark}

\begin{remark}[Double bosonization and the HT universal algebra]
\label{rem:dimofte-double-bosonization}
Majid's double bosonization construction \cite{Majid95} produces a quasi-triangular Hopf algebra from a braided Hopf algebra. Given a braided Hopf algebra $H$ in a braided monoidal category $\cC$, and its categorical dual $H^{*\mathrm{op}}$, the double bosonization is
\begin{equation}
\label{eq:double-bosonization}
R(H) \;=\; H^{*\mathrm{op}} \bowtie_{R^{21}} \cB \bowtie_R H,
\end{equation}
where $\cB$ is the transmutation of $\cC$ and $R$ is the $R$-matrix of the braiding. The double bosonization is universal in the following sense: any quasi-triangular Hopf algebra whose module category recovers $\cC$ as the category of braided $H$-modules factors through $R(H)$. When $\cC = \Vect$ and $H$ is ordinary, $R(H)$ is the classical Drinfeld double $D(H)$.
The conjectural picture at the $E_1$-chiral level is this. The ordered bar coalgebra $B^{\mathrm{ord}}(\cA) = T^c(s^{-1}\bar\cA)$ (augmentation ideal) carries the $R$-matrix $r(z)$ at degree two as the binary collision residue of the local OPE on $\bar\cA$. The full tower of higher collisions presents $B^{\mathrm{ord}}(\cA)$ as a braided Hopf object in a ``meromorphic $z$-dependent'' braided monoidal category: the category of line operators $\cC_{\mathrm{line}}$ discussed at length in this chapter. Apply the double bosonization~\eqref{eq:double-bosonization} to this braided Hopf datum. The output, $R(B^{\mathrm{ord}}(\cA))$, should be the HT universal algebra of the boundary theory: a quasi-triangular Hopf algebra whose module category is a rigidified enhancement of $\cC_{\mathrm{line}}$, and which extends the Drinfeld double programme (Remark~\ref{rem:drinfeld-double-programme}) from its ad-hoc assembly $\cA \bowtie \cA^!$ to a canonical universal construction.
One piece is missing: the correct definition of ``braided Hopf algebra in the meromorphic $z$-dependent tensor category''. Ordinary braided Hopf algebras live in a braided monoidal category with a scalar braiding; here the braiding is $r(z)$, which depends on a spectral parameter and carries poles. The spectral-braiding core programme of Chapter~\ref{chap:spectral-braiding} (Dimofte's item~6 of~\cite{Dimofte25}) supplies precisely this: a meromorphic braided category whose braiding is $r(z)$ with controlled pole behaviour. With that category in hand, the double bosonization~\eqref{eq:double-bosonization} becomes a well-defined functor on the $E_1$-chiral side, and $R(B^{\mathrm{ord}}(\cA))$ acquires a precise meaning.
Two caveats. First, this construction is stated as a conjecture; no part of it is proved here. Second, the Calabi--Yau parallel, in which the double bosonization acts on CY-braided objects and produces BPS algebras, is the natural home of this construction in Volume~III and is deferred to that volume.
\end{remark}

\section{Repository integration notes}
\label{sec:integration}

For direct integration into a research repository, the following replacements are recommended.

\begin{enumerate}[label=(\arabic*)]
\item Replace the abstract hypotheses in \S\ref{sec:setup} by the precise labels of the repository's
existing bar--cobar and module/comodule theorems.

\item Replace the shorthand notation $\widehat\otimes$ and $\chotimes$ by the repository's preferred
completed tensor notation and chiral monoidal symbol.

\item If the repository already has a sign appendix for the ordered bar complex, point the proofs of
Theorems~\ref{thm:opposite} and \ref{thm:tangent=K} to that appendix.

\item If the repository distinguishes coderived and contraderived categories more sharply, split
Theorem~\ref{thm:bimod-bicomod} accordingly.

\item The bordered configuration-space construction is now in Volume~I,
\S\ref*{V1-sec:bordered-fm}.
All five items of the original bordered conjecture are resolved:
items~(1)--(3) and~(5) by the Volume~I bordered FM
compactification, and item~(4) (annular closure computes
coHochschild homology) by
Theorem~\ref{thm:bordered-annular}.
\end{enumerate}

\section{Signs and conventions}

All formulas above are invariant under the usual choice of suspension conventions, provided one keeps
the same convention simultaneously in the ordered bar complex, the shuffle map, and the twisting
differential. For clarity we summarize the minimal sign content.

\begin{enumerate}[label=(\arabic*)]
\item Reversal in Theorem~\ref{thm:opposite} uses the standard suspended Koszul sign
\[
(-1)^{\sum_{i<j}(|a_i|+1)(|a_j|+1)}.
\]
The shift by~$1$ in each exponent reflects the bar grading: a bar word $[s^{-1}a_1|\cdots|s^{-1}a_n]$ has internal degree $|s^{-1}a_i| = |a_i| - 1$ (desuspension lowers cohomological degree by~$1$: $|s^{-1}| = -1$). Reversing the word permutes the $n$ desuspended elements past each other, and the Koszul sign is $(-1)^{(|a_i|-1)(|a_j|-1)}$ for transposing elements of degrees $|a_i|-1$ and $|a_j|-1$. Since $(|a|-1)\equiv(|a|+1)\pmod{2}$, this equals $(-1)^{(|a_i|+1)(|a_j|+1)}$ as written.

\item The left twisting term in Definition~\ref{def:Kbi} carries the sign $(-1)^{|c_{(1)}|}$, needed to move the
coderivation on $C$ past the first coalgebra leg before applying the left $A$-action via $\tau(c_{(2)})\cdot m$.

\item The right twisting term carries the additional sign $(-1)^{|c|+|m|}$, needed to move the twisting cochain across
the middle bimodule factor before applying the right $A$-action via $m\cdot\tau(d_{(1)})$.

\item In the strongly admissible regime the filtration is complete and bounded below, so all spectral
sequence and homotopy-transfer arguments used in the proofs converge. In particular, the PBW filtration on $A^e$ and the coradical filtration on $C^e$ are compatible under the bar--cobar equivalence, and the passage to associated gradeds commutes with the algebraic structures in question.
\end{enumerate}

\section{New results proved in this chapter}

For quick reference, here is the list of new results proved here rather than taken as repository
inputs.

\begin{enumerate}[label=(\arabic*)]
\item Ordered chiral shuffle theorem (Theorem~\ref{thm:shuffle}).
\item Opposite-duality and anti-involution transport (Theorem~\ref{thm:opposite},
Corollary~\ref{cor:anti}).
\item Enveloping duality (Corollary~\ref{cor:enveloping}).
\item Two-sided twisted bicomodule and its dg structure (Definition~\ref{def:Kbi},
Lemma~\ref{lem:Kbi-dg}).
\item Identification of the tangent ordered bar complex with the twisted bicomodule
(Theorem~\ref{thm:tangent=K}).
\item PBW-complete bimodule/bicomodule equivalence (Theorem~\ref{thm:bimod-bicomod}).
\item Diagonal correspondence (Theorem~\ref{thm:diagonal}).
\item Composition unit theorem (Corollary~\ref{cor:unit}).
\item Tensor--cotensor gluing (Corollary~\ref{cor:tensor-cotensor}).
\item Hochschild/coHochschild homology and cohomology equivalence
(Theorems~\ref{thm:HH-coHH-homology} and \ref{thm:HH-coHH-cohomology}).
\item Infinitesimal annular variation (Proposition~\ref{prop:infann}).
\item Ordered genus-zero open trace formalism (Theorem~\ref{thm:ordered-open}).
\item Calabi--Yau transport to an ordered Frobenius structure (Theorem~\ref{thm:CY}).
\item Covering-space frame for ordered-to-unordered descent
(Construction~\ref{constr:r-matrix-covering}).
\item $R$-matrix descent identification
(Proposition~\ref{prop:r-matrix-descent}).
\item Yang--Baxter equation from $d^2=0$ via dual topological-combinatorial
proof (Proposition~\ref{prop:ybe-from-d-squared}).
\item Classical Yang--Baxter equation as linearisation
(Corollary~\ref{cor:classical-ybe}).
\item Ordered AOS relations: $d^2=0$ from associativity alone, with explicit term
decomposition at degrees $3$--$5$ (Proposition~\ref{prop:ordered-aos}).
\item Euler-eta identity: master formula for the graded bar Euler characteristic
(Theorem~\ref{thm:euler-eta-identity}).
\item Bar cohomology concentration on the Koszul locus, with depth spectral sequence
(Theorem~\ref{thm:bar-cohomology-concentration}).
\end{enumerate}

\section{The ordered bar complex: explicit constructions}
\label{sec:ordered-bar-explicit}

The diagonal bicomodule $C_\Delta$, the Hochschild--coHochschild dictionary, and the $E_1$ five main theorems are the algebraic skeleton. The geometry that drives the skeleton lives on ordered configuration spaces: the bar differential is an integral over $\mathrm{Conf}_n^{\mathrm{ord}}(\C)$, its monodromy around the braid generator is the $R$-matrix, the Yang--Baxter equation is $d^2=0$, and for affine Kac--Moody input the resulting associative algebra with spectral parameter is the Yangian. We compute the ordered bar complex for the Heisenberg and $\widehat{\mathfrak{sl}}_2$ (where every formula can be checked by hand), extract the $R$-matrix and its Yang--Baxter equation, construct the dg-shifted Yangian, and close with the Euler-eta identity and bar cohomology concentration theorem.

\subsection{Ordered configuration spaces}
\label{subsec:ordered-config-spaces}

\begin{definition}[Ordered configuration space]
\label{def:ordered-config-space}
\index{configuration space!ordered|textbf}
For a smooth curve~$X$ and $n\ge 1$, the \emph{ordered
configuration space} is
\[
\mathrm{Conf}_n^{\mathrm{ord}}(X)
\;:=\;
\bigl\{(z_1,\dots,z_n)\in X^n
\;\big|\;
z_i\neq z_j\;\text{for all}\;i\neq j\bigr\}.
\]
The symmetric group~$\Sigma_n$ acts freely by permuting
coordinates. The unordered configuration space is the
quotient:
$\mathrm{Conf}_n(X)
=\mathrm{Conf}_n^{\mathrm{ord}}(X)/\Sigma_n$.
\end{definition}

\begin{proposition}[Topology of ordered configuration spaces;
\ClaimStatusProvedElsewhere]
\label{prop:ordered-config-topology}
\index{configuration space!ordered!topology}
\begin{enumerate}[label=\textup{(\roman*)}]
\item For $X=\mathbb C$:
 $\pi_1(\mathrm{Conf}_n^{\mathrm{ord}}(\mathbb C))
 =P_n$, the \emph{pure braid group} on~$n$ strands.
 The quotient
 $\pi_1(\mathrm{Conf}_n(\mathbb C))=B_n$
 is the full braid group.
\item $\mathrm{Conf}_n^{\mathrm{ord}}(\mathbb C)$ is a
 $K(\pi,1)$: its universal cover is contractible, so the
 homotopy type is determined by~$P_n$.
\item $H^*(\mathrm{Conf}_n^{\mathrm{ord}}(\mathbb C);\Bbbk)$
 is the \emph{Arnold algebra}: generated by classes
 $\omega_{ij}$ ($1\le i<j\le n$, degree~$1$) with relations
 $\omega_{ij}^2=0$ and the Arnold relation
 $\omega_{ij}\omega_{jk}+\omega_{jk}\omega_{ki}
 +\omega_{ki}\omega_{ij}=0$.
 For the \emph{ordered} space, no symmetric-group quotient
 is taken and the $\omega_{ij}$ for $i>j$ are independent
 generators: $\omega_{ji}=-\omega_{ij}+d(\cdots)$ at the
 cochain level.
\item For $X=\Sigma_g$ (genus~$g$):
 $\pi_1(\mathrm{Conf}_n^{\mathrm{ord}}(\Sigma_g))$
 is the \emph{surface pure braid group}, an extension
 of~$P_n$ by~$\pi_1(\Sigma_g)^n$. The additional
 generators are the loops winding around handles of the
 surface; they produce the period corrections in the
 ordered bar differential.
\end{enumerate}
\end{proposition}

\begin{remark}[Why ordered spaces for $E_1$-algebras]
\label{rem:why-ordered}
\index{E1 algebra@$E_1$-algebra!ordered configuration spaces}
An $E_\infty$-chiral algebra has
$\Sigma_n$-equivariant operations, so its bar complex
uses the unordered space $\mathrm{Conf}_n(X)$.
The pole-free BD-commutative subclass is a special case.
An $E_1$-algebra (associative chiral algebra) has
operations indexed by \emph{ordered} compositions, so its
bar complex uses $\mathrm{Conf}_n^{\mathrm{ord}}(X)$.
The ordering records the topological direction: in
$\mathbb R_+\times\mathbb C$, points along the
$\mathbb R_+$-axis are ordered by time, and the
$E_1$-structure is the time-ordered product. The
$R$-matrix $R(z)\in\operatorname{End}(V\otimes V)(\!(z)\!)$
is the data needed to pass from ordered to unordered:
it is the monodromy of the local system on
$\mathrm{Conf}_2^{\mathrm{ord}}(\mathbb C)$ around the
generator of $\pi_1=\mathbb Z$.
\end{remark}

\begin{construction}[Compactification: ordered
Fulton--MacPherson space; \ClaimStatusProvedElsewhere]
\label{constr:ordered-fm}
\index{Fulton--MacPherson!ordered compactification}
The ordered FM compactification
$\overline{\mathrm{FM}}_n^{\mathrm{ord}}(X)$ is the
iterated real oriented blowup of
$\mathrm{Conf}_n^{\mathrm{ord}}(X)$ along the ordered
diagonals $D_{i,i+1}=\{z_i=z_{i+1}\}$ ($i=1,\dots,n-1$).
Codimension-one boundary faces correspond to
\emph{consecutive} collisions $z_i\to z_{i+1}$
(not arbitrary subsets, as in the unordered case).
The boundary stratification is indexed by
\emph{planar trees}: trees whose leaves carry a planar
(= total) ordering inherited from the ordering of the
points.

For $n=2$:
$\overline{\mathrm{FM}}_2^{\mathrm{ord}}(\mathbb C)
\cong\overline{\mathbb C}$ (the Riemann sphere), with
boundary divisor $D_{1,2}=\{z_1=z_2\}$ (one point at
infinity).

For $n=3$:
$\overline{\mathrm{FM}}_3^{\mathrm{ord}}(\mathbb C)$
has \emph{two} codimension-one faces: $D_{1,2}$ (first two
points collide) and $D_{2,3}$ (last two collide). The
face $D_{1,3}$ (first and last collide without the middle
point) is codimension~$2$: it requires both $D_{1,2}$
and $D_{2,3}$ to degenerate simultaneously. This is the
associahedron $K_3$: a line segment with two endpoints.
\end{construction}

\subsection{The ordered bar differential}
\label{subsec:ordered-bar-definition}

\begin{definition}[Ordered chiral bar complex]
\label{def:ordered-chiral-bar}
\index{bar complex!ordered chiral|textbf}
Let $\cA$ be an $E_1$-chiral algebra on a curve~$X$. The
\emph{ordered chiral bar complex} is the differential graded
coalgebra
\[
\Barchord(\cA)
\;=\;
\bigoplus_{n\ge 1}
(s^{-1}\bar\cA)^{\otimes n}
\;\otimes\;
C_\bullet\!\bigl(\mathrm{Conf}_n^{\mathrm{ord}}(X)\bigr),
\]
where $\mathrm{Conf}_n^{\mathrm{ord}}(X)
=\{(z_1,\dots,z_n)\in X^n\mid z_i\neq z_j\;\forall\,i\neq j\}$
is the \emph{ordered} configuration space (no $\Sigma_n$-quotient).
The bar differential is
\[
d_{\Barchord}
\;=\;
d_{\mathrm{int}}+d_{\mathrm{res}},
\]
where $d_{\mathrm{int}}$ applies the internal differential
of~$\cA$ to each tensor factor and $d_{\mathrm{res}}$
extracts OPE residues at ordered collision divisors
$D_{i,i+1}=\{z_i=z_{i+1}\}$ using the \emph{ordered}
logarithmic forms $\eta_{i,i+1}=d\log(z_i-z_{i+1})$.

\end{definition}

\begin{construction}[The ordered bar differential at low degree;
\ClaimStatusProvedHere]
\label{constr:ordered-bar-low-degree}
\index{bar complex!ordered!low degree formulas}
Let $\{e_I\}$ be a homogeneous basis for $\bar\cA$ with
structure constants $e_I{}_{(n)}\,e_J=\sum_K c^K_{IJ;n}\,e_K$.

\emph{Degree~$1$}: $\Barchord_1=s^{-1}\bar\cA$.
The differential is $d(s^{-1}e_I)=0$ (no contractions
possible on a single element).

\emph{Degree~$2$}: $\Barchord_2
=\bigoplus_{I,J}s^{-1}e_I\otimes s^{-1}e_J\otimes
C_\bullet(\mathrm{Conf}_2^{\mathrm{ord}})$.
The differential contracts the unique \emph{ordered} pair
$(z_1,z_2)$ at the collision $z_1\to z_2$:
\begin{equation}\label{eq:ordered-bar-d2}
d(s^{-1}e_I\otimes s^{-1}e_J)
\;=\;
\sum_{K,n}\pm\,c^K_{IJ;n}\;s^{-1}e_K
\;\otimes\;
[\delta_{D_{1,2}}^{(n)}],
\end{equation}
where $[\delta_{D_{1,2}}^{(n)}]$ is the $n$-th residue class
at the collision divisor. For the simple-pole residue
($n=0$): this is $e_I{}_{(0)}\,e_J$, the Lie bracket.
For higher poles ($n\ge 1$): these are the higher products
$e_I{}_{(n)}\,e_J$.

The sign: $(-1)^{|e_I|}$ from permuting $s^{-1}$ past
$e_I$ in the desuspension (cohomological convention
$|d|=+1$).

\emph{Degree~$3$}: $\Barchord_3
=\bigoplus_{I,J,K}s^{-1}e_I\otimes s^{-1}e_J\otimes
s^{-1}e_K\otimes
C_\bullet(\mathrm{Conf}_3^{\mathrm{ord}})$.
Two codimension-one faces: $D_{1,2}$ (first pair collides)
and $D_{2,3}$ (second pair collides). The differential:
\begin{equation}\label{eq:ordered-bar-d3}
d(s^{-1}e_I\otimes s^{-1}e_J\otimes s^{-1}e_K)
\;=\;
\sum_L\pm\,c^L_{IJ;0}\;
s^{-1}e_L\otimes s^{-1}e_K
\;\otimes\;[\delta_{D_{1,2}}]
\;+\;
\sum_M\pm\,c^M_{JK;0}\;
s^{-1}e_I\otimes s^{-1}e_M
\;\otimes\;[\delta_{D_{2,3}}].
\end{equation}
$d^2=0$ at degree~$3$: the composition $d\circ d$ on
$s^{-1}e_I\otimes s^{-1}e_J\otimes s^{-1}e_K$ produces
\[
d^2\;\ni\;
\sum_{L,M,N}\pm\bigl(
c^N_{LK;0}\,c^L_{IJ;0}
-c^N_{IM;0}\,c^M_{JK;0}
\bigr)\,s^{-1}e_N
\;\otimes\;[\delta_{D_{1,2}}\cap\delta_{D_{2,3}}].
\]
This vanishes if and only if
$(e_I{}_{(0)}e_J){}_{(0)}e_K
=e_I{}_{(0)}(e_J{}_{(0)}e_K)$,
which is \emph{associativity of the $0$-th product}, the
$E_1$ axiom. (For commutative chiral algebras, the
$0$-th product satisfies the Jacobi identity instead,
and the Arnold relation provides the additional
cancellation. In the ordered complex, no Arnold relation
is needed: associativity alone suffices.)
\end{construction}

\begin{proposition}[Ordered AOS relations: $d^2=0$ from
associativity alone; \ClaimStatusProvedHere]
\label{prop:ordered-aos}
\index{AOS relations!ordered|textbf}
\index{bar complex!ordered!d squared zero@$d^2=0$}
The identity $d^2=0$ on
$\Barchord_n(\cA)$ for $n=3,4,5$ follows from
\emph{associativity of the $0$-th product} and
\emph{sign alternation on independent face pairs}. No Jacobi
identity, no Arnold relation, and no symmetric-group
equivariance are required.

At each degree~$n$, the terms in $d^2$ decompose into two
classes:
\begin{enumerate}[label=\textup{(\roman*)}]
\item \emph{Overlapping terms}: two consecutive collisions
 that share a common index. These cancel by associativity
 of the zeroth product $a_{(0)}(b_{(0)}c)
 =(a_{(0)}b)_{(0)}c$.
\item \emph{Independent terms}: two consecutive collisions
 at disjoint pairs $(i,i{+}1)$ and $(j,j{+}1)$ with
 $|i-j|\ge 2$. These cancel by sign alternation:
 the two orderings of independent face collapses produce
 identical algebraic contributions with opposite Koszul
 signs.
\end{enumerate}

Explicit verification at low degrees:

\emph{Degree~$3$}: Two faces $D_{1,2}$ and $D_{2,3}$.
These overlap at index~$2$. The unique codimension-two
stratum is $D_{1,2}\cap D_{2,3}=\{z_1{=}z_2{=}z_3\}$,
and $d^2[a|b|c]=\pm\bigl((a_{(0)}b)_{(0)}c
-a_{(0)}(b_{(0)}c)\bigr)=0$ by associativity.
Face count: $2$ faces of the associahedron~$K_3$
(a line segment).

\emph{Degree~$4$}: Three faces $D_{1,2}$, $D_{2,3}$,
$D_{3,4}$. The codimension-two strata are:
$D_{1,2}\cap D_{2,3}$ (overlapping, cancelled by
associativity), $D_{2,3}\cap D_{3,4}$ (overlapping,
cancelled by associativity), and $D_{1,2}\cap D_{3,4}$
(independent, cancelled by sign alternation).
Face count: $3$ consecutive collision divisors of
$\overline{\mathrm{FM}}_4^{\mathrm{ord}}$.

\emph{Degree~$5$}: Four faces $D_{1,2},\ldots,D_{4,5}$.
Six codimension-two strata: three overlapping pairs
$(D_{1,2}\cap D_{2,3}$, $D_{2,3}\cap D_{3,4}$,
$D_{3,4}\cap D_{4,5})$ and three independent pairs
$(D_{1,2}\cap D_{3,4}$, $D_{1,2}\cap D_{4,5}$,
$D_{2,3}\cap D_{4,5})$. All cancel by the two mechanisms.
Face count: $4$ faces of the associahedron~$K_5$
(the Stasheff polytope with $14$ vertices).

The geometric content: the ordered FM compactification
$\overline{\mathrm{FM}}_n^{\mathrm{ord}}(\mathbb C)$
has $n{-}1$ codimension-one faces (consecutive collisions
only), indexed by the facets of the associahedron~$K_n$.
The unordered compactification
$\overline{\mathrm{FM}}_n(\mathbb C)$ has
$\binom{n}{2}$ codimension-one faces (all pair collisions),
indexed by the facets of the permutohedron. The ratio
$(n{-}1)/\binom{n}{2}=2/n$ measures the combinatorial
simplification of the ordered theory.
\end{proposition}

\begin{proof}
By induction on degree, using the two cancellation mechanisms.

At degree~$n$, the operator $d^2$ on
$[a_1|\cdots|a_n]\in\Barchord_n$ is
\[
d^2 = \sum_{1\le i<j\le n-1}
\pm\,d_{D_{j,j+1}}\circ d_{D_{i,i+1}}.
\]
Decompose the sum according to whether $j=i+1$
(overlapping) or $j\ge i+2$ (independent).

\emph{Overlapping pairs} ($j=i+1$):
$d_{D_{i+1,i+2}}\circ d_{D_{i,i+1}}$ collapses
$a_i,a_{i+1}$ first, then the result with $a_{i+2}$,
producing $(a_i{}_{(0)}a_{i+1}){}_{(0)}a_{i+2}$.
The reverse order
$d_{D_{i,i+1}}\circ d_{D_{i+1,i+2}}$ produces
$a_i{}_{(0)}(a_{i+1}{}_{(0)}a_{i+2})$.
These appear with the same sign in $d^2$
(both are collapses at the common codimension-two
stratum $D_{i,i+1}\cap D_{i+1,i+2}$), so their
difference is
$\pm\bigl((a_i{}_{(0)}a_{i+1}){}_{(0)}a_{i+2}
-a_i{}_{(0)}(a_{i+1}{}_{(0)}a_{i+2})\bigr)=0$
by associativity.

\emph{Independent pairs} ($j\ge i+2$):
the two faces $D_{i,i+1}$ and $D_{j,j+1}$ are disjoint
in $\overline{\mathrm{FM}}_n^{\mathrm{ord}}$. The
algebraic contributions of $d_{D_{j,j+1}}\circ d_{D_{i,i+1}}$
and $d_{D_{i,i+1}}\circ d_{D_{j,j+1}}$ are identical
(the two collapses act on disjoint tensor slots), but
the Koszul signs differ by $(-1)$: this is the standard
sign arising from the orientation of the boundary of
a manifold with corners, where the two orderings of
codimension-one faces meeting at a codimension-two corner
carry opposite orientations.
\end{proof}

\begin{remark}[Depth decomposition of $d^2=0$]
\label{rem:depth-decomposition-d-squared}
\index{bar complex!ordered!depth decomposition of $d^2=0$}
Proposition~\ref{prop:ordered-aos} treats the zeroth-product
component $d_0$ of the ordered bar differential, the leading
term in the pole-order depth filtration (see the proof of
Theorem~\ref{thm:bar-cohomology-concentration}(iv) below for the
full filtration). The \emph{full} ordered differential
decomposes as $d=d_0+d_1+d_2+\cdots$, where $d_p$ is the
component involving $n$-th products $a_{(n)}b$ at pole depth
$n=p$. The identity $d^2=0$ then decomposes by total depth:
\begin{enumerate}[label=\textup{(\roman*)}]
\item \emph{Depth~$0$}: $d_0^2=0$. This is associativity of
 the zeroth product plus sign alternation on independent face
 pairs, precisely the content of the proposition.
\item \emph{Depth~$1$}: $d_0 d_1+d_1 d_0=0$. This is
 the compatibility between the Lie bracket $a_{(0)}b$ and
 the first higher product $a_{(1)}b$ (the central extension /
 Killing form term for affine algebras). It expresses
 sesquilinearity of the $\lambda$-bracket at linear order.
\item \emph{Depth~$2$}: $d_0 d_2+d_1^2+d_2 d_0=0$.
 For affine algebras (pole order~$\le 2$) this is the
 highest nontrivial condition; for the Virasoro algebra
 (pole order~$4$) it is an intermediate identity.
\item \emph{General depth~$p$}:
 $\sum_{i+j=p}d_i\,d_j=0$.
 These are the PVA axioms (sesquilinearity, Leibniz rules,
 Jacobi identity for the $\lambda$-bracket) expressed in
 the ordered bar language.
\end{enumerate}
The key point is that the \emph{combinatorial mechanism}
established by the proposition (overlapping faces cancel by
associativity, independent faces cancel by sign
alternation) persists at every depth level. The higher-depth
conditions $\sum_{i+j=p}d_i\,d_j=0$ use exactly the same
face-pair decomposition into overlapping and independent
strata; what changes is the \emph{algebraic identity}
invoked at the overlapping strata (higher-product analogues
of associativity, i.e.\ the PVA axioms). The spectral-parameter
generalization does not alter the fundamental distinction
between the ordered bar (associahedron faces, associativity)
and the symmetric bar (permutohedron faces, Jacobi + Arnold).
\end{remark}

\begin{remark}[Jacobi is not needed: the $\mathfrak{sl}_2$ counterexample]
\label{rem:jacobi-not-needed}
\index{bar complex!ordered!Jacobi not needed}
A persistent confusion in the literature is the belief
that $d^2=0$ on the bar complex requires the Jacobi identity.
This is true for the \emph{unordered} (Chevalley--Eilenberg)
bar complex but \textbf{false} for the ordered bar complex.

Consider $\mathfrak{sl}_2$ with generators $e,f,h$
and the ordered bar complex built from the Lie bracket
$a_{(0)}b=[a,b]$ with \emph{only consecutive
collapses}. At degree~$3$,
$d^2=0$ would require associativity of the bracket:
$(e_{(0)}f)_{(0)}h=e_{(0)}(f_{(0)}h)$. But the Lie bracket
is not associative:
\[
 d^2[e|f|h]
 =[[e,f],h]-[e,[f,h]]
 =[h,h]-[e,2f]
 =0-2h
 =-2h\ne 0.
\]
This is exactly the failure of associativity for the Lie
bracket, and it shows that $d^2=0$ on the ordered-consecutive
bar complex genuinely requires associativity, not the Jacobi
identity. For the full chiral algebra $V_k(\mathfrak{sl}_2)$,
the higher products $a_{(n)}b$ with $n\ge 1$ contribute
additional terms to the bar differential that restore $d^2=0$.

Now consider the non-associative ordered bar complex in
which \emph{all} pair collapses are included (not just
consecutive ones). At degree~$3$, $d^2[e|f|h]$ acquires
contributions from the non-consecutive collision $D_{1,3}$
(first and third slots collide), giving
$d^2[e|f|h]\ni\pm[e,h]\ne 0$. This is the
$\mathfrak{sl}_2$ counterexample: $12$ of the $27$
ordered triples have $d^2\ne 0$ in the non-consecutive
theory.

The resolution: the \emph{ordered} bar complex uses only
consecutive collisions (the associahedron boundary), and
$d^2=0$ follows from associativity. The \emph{unordered}
bar complex uses all pair collisions (the permutohedron
boundary), and $d^2=0$ requires both the Jacobi identity
and the Arnold relation. For $E_\infty$-chiral algebras
both are available; for $E_1$-chiral algebras only
associativity is available, and only the ordered complex
is well-defined.
\end{remark}

\begin{remark}[Affine $\mathfrak{sl}_2$: ordered bar with
all-pair collapses has $d^2\ne 0$]
\label{rem:sl2-d-squared-nonzero}
\index{Kac--Moody!ordered bar!d squared nonzero@$d^2\ne 0$}
The distinction between ordered-consecutive and
ordered-all-pair is not academic. For
$V_k(\mathfrak{sl}_2)$ with generators $\{e,f,h\}$, the
ordered bar complex with \emph{all} pair collapses
(including non-consecutive) has $d^2\ne 0$ on $12$ of the
$27$ basis triples:
\[
\renewcommand{\arraystretch}{1.1}
\begin{array}{ll}
d^2\ne 0\colon & [e|f|h],\;[e|h|f],\;[f|e|h],\;
 [f|h|e],\;[h|e|f],\;[h|f|e], \\
 & [e|f|e],\;[e|f|f],\;[f|e|e],\;
 [f|e|f],\;[e|h|e],\;[f|h|f]. \\[4pt]
d^2=0\colon & [e|e|e],\;[f|f|f],\;[h|h|h], \\
 & [e|e|f],\;[e|e|h],\;[f|f|e],\;
 [f|f|h],\;[h|h|e],\;[h|h|f], \\
 & [e|h|h],\;[f|h|h],\;[h|e|e],\;
 [h|f|f],\;[h|e|h],\;[h|f|h].
\end{array}
\]
This is expected: $V_k(\mathfrak{sl}_2)$ is
$E_\infty$-chiral (it is a local vertex algebra), so its
operations are defined on unordered configurations. The
ordered-all-pair differential conflates the ordered and
unordered structures without the $R$-matrix twist
needed for descent
(Proposition~\ref{prop:r-matrix-descent}), and the
resulting inconsistency manifests as $d^2\ne 0$.

The \emph{Yangian} $Y_\hbar(\mathfrak{sl}_2)$, by contrast,
is genuinely $E_1$-chiral: its operations are defined on
ordered configurations with only consecutive collisions.
The ordered-consecutive differential has $d^2=0$ by
Proposition~\ref{prop:ordered-aos}, and the resulting
Koszul dual is $U(\mathfrak{sl}_2[t])$
(Proposition~\ref{prop:open-colour-double-bar}).
The ordered bar complex ``belongs to'' the Yangian, not
the affine algebra.
\end{remark}

\begin{remark}[Ordered vs.\ unordered]
\label{rem:ordered-vs-unordered}
\index{bar complex!ordered vs unordered}
The unordered bar complex
$\Barch(\cA)=\Barchord(\cA)_{\Sigma_n}$
is the $\Sigma_n$-coinvariants of the ordered complex.
The comparison map
$\Barchord(\cA)\to\Barch(\cA)$
is the quotient by the symmetric group action. For
$E_\infty$-chiral algebras in the pole-free
BD-commutative subclass, the two are equivalent. For local
vertex algebras with OPE poles, the comparison is still
$E_\infty$-chiral but requires the OPE-derived
$R$-matrix twist. For genuinely $E_1$-chiral
algebras (Yangians, quantum groups), the ordered complex
carries strictly more information: the $R$-matrix is the
data needed to descend from ordered to unordered.

Geometrically, the descent datum is non-trivial because
$\mathrm{Conf}_n^{\mathrm{ord}}(\mathbb R)$ is discrete:
its connected components are the $n!$~chambers cut out by
the total order, while
$\mathrm{Conf}_n(\mathbb C)$ is connected, with fundamental
group the braid group~$B_n$. Passing from the
$E_1$-direction~$\mathbb R$ to the holomorphic
plane~$\mathbb C$ destroys the total order, and the
$\Sigma_n$-equivariant structure that was tautological on
chambers must be \emph{constructed} on the covering
$\mathrm{Conf}_n^{\mathrm{ord}}(\mathbb C)\to
\mathrm{Conf}_n(\mathbb C)$, where it becomes monodromy
data, a representation of~$B_n$. The next subsection
makes this precise: the $R$-matrix is the descent datum
for this covering, extracted from the flat connection that
the bar differential defines on ordered configuration
space.
\end{remark}

\subsection{The Heisenberg ordered bar complex}
\label{subsec:ordered-bar-heisenberg}

\begin{computation}[Ordered bar complex for Heisenberg;
\ClaimStatusProvedHere]
\label{comp:ordered-bar-heisenberg}
\index{Heisenberg algebra!ordered bar complex}
For $\cH_k$ on $X=\mathbb A^1$: the ordered configuration
space $\mathrm{Conf}_n^{\mathrm{ord}}(\mathbb A^1)$ is
contractible (it deformation-retracts onto
$\{z_1<z_2<\cdots<z_n\}\cong\mathbb R^n_+$).
The ordered bar complex reduces to
\[
\Barchord(\cH_k)
\;\simeq\;
\bigoplus_{n\ge 1}(s^{-1}\alpha)^{\otimes n}
\]
with differential $d(s^{-1}\alpha\otimes s^{-1}\alpha)
=\pm s^{-1}(\alpha_{(1)}\alpha)=\pm s^{-1}(k)$
(the residue of the double pole, after the $d\log$ kernel
absorbs one power from the double-pole OPE).
The bar cohomology concentrates immediately:
$H^0=\Bbbk$, $H^1=\Bbbk\langle s^{-1}\alpha\rangle$,
$H^n=0$ for $n\ge 2$.
(The weight-depth identity $w + d = n - 1$ of
Theorem~\ref{thm:gap-migration}(i) is trivially
satisfied at the unique nontrivial degree: $w = 0$,
$d = 1$, $n = 2$.)

The Koszul dual is $(\cH_k)^!=\mathrm{Sym}^{\mathrm{ch}}(V^*)$
(the chiral symmetric algebra on the dual; Vol~I, \S\ref*{V1-sec:heisenberg-koszul}). The $R$-matrix is
the scalar $R(z)=z^k$: nontrivial as a function of~$z$ (the double-pole OPE produces the collision residue $r(z) = k/z$ after $d\log$ absorption), but scalar-valued because the generator space is one-dimensional. The ordered and unordered bar complexes are therefore equivalent up to the $R$-twisted sign character of~$\Sigma_n$.
\end{computation}

\subsection{The ordered bar complex for $\widehat{\mathfrak{sl}}_2$}
\label{subsec:ordered-bar-sl2}

\begin{computation}[Ordered bar complex for
$\widehat{\mathfrak{sl}}_2$ at low degree;
\ClaimStatusProvedHere]
\label{comp:ordered-bar-sl2}
\index{Kac--Moody!ordered bar complex}
For $\widehat{\mathfrak{sl}}_2$ at level~$k$, the generators
are the three currents $e(z)$, $f(z)$, $h(z)$ of conformal
weight~$1$. The ordered bar complex at tensor degree~$2$
has $3^2=9$ generators:
\[
\Barchord_2
\;=\;
\operatorname{span}\{
s^{-1}e\otimes s^{-1}e,\;
s^{-1}e\otimes s^{-1}f,\;
s^{-1}e\otimes s^{-1}h,\;
\dots\}
\;\otimes\;
C_\bullet(\mathrm{Conf}_2^{\mathrm{ord}}(\mathbb A^1)).
\]
Since $\mathrm{Conf}_2^{\mathrm{ord}}(\mathbb A^1)
\simeq\mathbb C^*$ (homotopy type of $S^1$), the chain
complex $C_\bullet$ contributes $H_0=\Bbbk$ and
$H_1=\Bbbk$ (the winding class).

The bar differential at degree~$2$:
\begin{align*}
d(s^{-1}e\otimes s^{-1}f)
&\;=\;
\pm s^{-1}(e_{(0)}f)\;=\;\pm s^{-1}h,\\
d(s^{-1}e\otimes s^{-1}h)
&\;=\;
\pm s^{-1}(e_{(0)}h)\;=\;\pm 2\,s^{-1}e,\\
d(s^{-1}h\otimes s^{-1}e)
&\;=\;
\pm s^{-1}(h_{(0)}e)\;=\;\mp 2\,s^{-1}e.
\end{align*}
The ordered structure: $d(s^{-1}e\otimes s^{-1}h)$ and
$d(s^{-1}h\otimes s^{-1}e)$ differ by a sign (the ordering
of insertions matters). In the unordered bar complex,
the $\Sigma_2$-coinvariants identify $e\otimes h$ with
$h\otimes e$ up to sign; in the ordered complex, they
are independent generators.

Bar cohomology at degree~$1$:
$H^1(\Barchord_2)$ is the Koszul dual
coalgebra. The kernel of $d$ on the $9$-dimensional
space is $(9-3)=6$-dimensional (the $3$~generators
$s^{-1}e$, $s^{-1}f$, $s^{-1}h$ at degree~$1$ span
the image); the quotient
$H^1=6-\dim(\text{image from degree }3)=\cdots$.
The full computation gives
$\dim H^1(\Barchord)=3=\dim\mathfrak{sl}_2$
(the Koszul dual has the same dimension as the Lie
algebra, consistent with
$(\widehat{\mathfrak{sl}}_2)^!
=\widehat{\mathfrak{sl}}_2{}_{-k-4}$).

\emph{The $R$-matrix from the ordered bar complex.}
The $R$-matrix $R(z)\in\operatorname{End}(V\otimes V)(\!(z)\!)$
for $V=\mathfrak{sl}_2$ is the operator that intertwines
the two orderings:
\[
R(z)\colon
\Barchord_2(e_i\otimes e_j;\,z_1<z_2)
\;\xrightarrow{\;\sim\;}
\Barchord_2(e_j\otimes e_i;\,z_2<z_1).
\]
For $\widehat{\mathfrak{sl}}_2$ at level~$k$:
$R(z)=1+\hbar\,k\,\Omega/z+O(z^{-2})$ with
$\hbar=1/(k+2)$, where
$\Omega=e\otimes f+f\otimes e+\tfrac12 h\otimes h$
is the Casimir. (At $k=0$ the level-stripped
coefficient $\hbar\to 1/2$ remains nonzero because
the $\widehat{\mathfrak{sl}}_2$ Sugawara shift
$k+2$ survives; the strict classical $r$-matrix
on the underlying affine current algebra is
$k\,\Omega/z$, which vanishes at $k=0$.) The Yang--Baxter equation
$R_{12}(z_1{-}z_2)\,R_{13}(z_1{-}z_3)\,
R_{23}(z_2{-}z_3)
=R_{23}(z_2{-}z_3)\,R_{13}(z_1{-}z_3)\,
R_{12}(z_1{-}z_2)$
is equivalent to $d^2=0$ on
$\Barchord_3$: the three orderings
$(z_1<z_2<z_3)$, $(z_1<z_3<z_2)$, $(z_2<z_1<z_3)$
produce the three factors, and their consistency is
the codimension-two cancellation on the ordered
configuration space.
\end{computation}

\begin{computation}[$E_1$ vs $E_\infty$ dimension surplus;
\ClaimStatusProvedHere]
\label{comp:e1-einfty-dimension-surplus}
\index{bar complex!ordered vs symmetric!dimension surplus|textbf}
The ordered bar complex $\Barchord(\cA)_n$
has strictly larger dimension than the symmetric bar complex
$\Barch(\cA)_n$ whenever the algebra has more than one
generator. The surplus
$\Delta_n := \dim\Barchord_n
- \dim\Barch_n$ carries the pure braid group
representation: the Kontsevich integral identifies
$\Delta_n$ with a space of Vassiliev invariants at
degree~$n$. For a single-generator algebra the
surplus is zero in raw dimension but nontrivial in
spectral-parameter dependence (the directed spectral
flow that $\Sigma_n$-coinvariants destroy).
\[
\renewcommand{\arraystretch}{1.15}
\begin{array}{lcrrrl}
\text{Family}
 & n
 & \dim\Barchord_n
 & \dim\Barch_n
 & \Delta_n
 & \text{Content of surplus} \\
\hline
\widehat{\mathfrak{sl}}_2
 & 2 & 9 & 6 & 3
 & \mathrm{adj}(\mathfrak{sl}_2) \\
\widehat{\mathfrak{sl}}_2
 & 3 & 27 & 10 & 17
 & \text{braid representation } B_3\to\operatorname{GL}(V^{\otimes 3}) \\
\widehat{\mathfrak{sl}}_2
 & 4 & 81 & 15 & 66
 & \text{higher Vassiliev data} \\
\mathrm{Vir}_c
 & 2\text{--}4 & 1 & 1 & 0
 & \text{directed spectral flow}
\end{array}
\]
The $\widehat{\mathfrak{sl}}_2$ dimensions: at tensor
degree~$n$ with all generators at weight~$1$,
$\dim\Barchord_n = d^n = 3^n$
(ordered sequences of~$d = 3$ generators), while
$\dim\Barch_n = \binom{d+n-1}{n} = \binom{n+2}{n}$
(unordered multisets, via the symmetric bar).
At $n = 2$: $9$ vs $6$, surplus $= 3 = \dim(\mathfrak{sl}_2)$,
the adjoint representation, spanned by the
antisymmetric tensors $e_I\otimes e_J - R(z)\,e_J\otimes e_I$
that the $\Sigma_2$-quotient kills.
At $n = 3$: $27$ vs $10$, surplus $= 17$, decomposing into
irreducible $B_3$-representations that carry the
degree-$3$ Vassiliev data. At $n = 4$: $81$ vs $15$,
surplus $= 66$.

For Virasoro ($d = 1$, single generator of weight~$2$):
$\dim\Barchord_n = 1 = \dim\Barch_n$
at each tensor degree (one ordered sequence and one
unordered multiset of length~$n$ from a single generator).
The surplus is \emph{not} in raw dimensions but in the
\emph{spectral asymmetry}: the ordered bar differential
$d_{\mathrm{res}}$ on
$\Barchord_n(\mathrm{Vir})$
retains the spectral-parameter dependence
$r(z) = (c/2)/z^3 + 2T/z$ (the collision residue of
the quartic-pole OPE, after $d\log$ absorption), while
the symmetric bar sees only the $\Sigma_n$-invariant
part. The directed spectral flow (the fact that
$r(z) \ne r(-z)$ due to the odd-power term
$2T/z$) is the surplus datum, invisible at the
level of dimensions but controlling the higher
$\Ainf$ operations $m_k$ ($k \ge 3$) that characterise
class~$\mathbf{M}$.
\end{computation}

\subsection{The $R$-matrix as ordered-to-unordered descent}
\label{subsec:r-matrix-monodromy}

The ordered bar complex $\Barchord(\cA)$ carries
strictly more information than the unordered complex $\Barch(\cA)$
whenever $\cA$ is genuinely $E_1$. The surplus has a precise name:
it is the \emph{descent datum} for the $\Sigma_n$-covering
$\mathrm{Conf}_n^{\mathrm{ord}}(X)\to\mathrm{Conf}_n(X)$.
The $R$-matrix is this datum.

\begin{construction}[The covering-space frame;
\ClaimStatusProvedHere]
\label{constr:r-matrix-covering}
\index{R-matrix!descent datum|textbf}
\index{descent!ordered to unordered bar complex}
The projection
\[
\pi_n\colon
\mathrm{Conf}_n^{\mathrm{ord}}(X)
\;\longrightarrow\;
\mathrm{Conf}_n(X)
\;=\;
\mathrm{Conf}_n^{\mathrm{ord}}(X)/\Sigma_n
\]
is a principal $\Sigma_n$-bundle, \'etale with fibre~$\Sigma_n$.
A local system~$\mathcal{L}$ on the base is equivalent to a
$\Sigma_n$-equivariant local system~$\pi_n^*\mathcal{L}$ on the total
space: the descent correspondence
\[
\mathrm{LocSys}\bigl(\mathrm{Conf}_n(X)\bigr)
\;\simeq\;
\mathrm{LocSys}^{\Sigma_n}
\bigl(\mathrm{Conf}_n^{\mathrm{ord}}(X)\bigr).
\]
The bar complex at tensor degree~$n$ defines a local system
$\mathcal{F}_n^{\mathrm{ord}}$ on $\mathrm{Conf}_n^{\mathrm{ord}}(X)$
whose stalk at $(z_1,\dots,z_n)$ is
$(s^{-1}\bar\cA)^{\otimes n}$, with flat connection induced by
$d_{\mathrm{res}}$. To descend $\mathcal{F}_n^{\mathrm{ord}}$ to a
local system on $\mathrm{Conf}_n(X)$, one must exhibit a
$\Sigma_n$-equivariant structure: isomorphisms
\[
R_\sigma\colon
\mathcal{F}_n^{\mathrm{ord}}\big|_{(z_1,\dots,z_n)}
\;\xrightarrow{\;\sim\;}
\mathcal{F}_n^{\mathrm{ord}}\big|_{(z_{\sigma(1)},\dots,z_{\sigma(n)})}
\]
for each $\sigma\in\Sigma_n$, satisfying the cocycle condition
$R_{\sigma\tau}=R_\tau\circ R_\sigma$ and compatible with the
flat connection.

For an $E_\infty$-chiral algebra, the equivariant structure
is canonical: $R_\sigma=\epsilon(\sigma)\cdot\tau_\sigma$, where
$\tau_\sigma$ permutes tensor factors and $\epsilon(\sigma)$ is
the Koszul sign. The descent recovers the ordinary
$\Sigma_n$-coinvariant bar complex.

For an $E_1$-chiral algebra, the operations respect only
the linear ordering. The equivariant structure must be
\emph{constructed} from the OPE data, and the construction
is the $R$-matrix.
\end{construction}

\begin{construction}[$R$-matrix as monodromy of the ordered
bar complex; \ClaimStatusProvedHere]
\label{constr:r-matrix-monodromy}
\index{R-matrix!from monodromy|textbf}
\index{monodromy!R-matrix}
The simplest descent instance is $n=2$. The ordered
configuration space
$\mathrm{Conf}_2^{\mathrm{ord}}(\mathbb C)
=\{(z_1,z_2)\in\mathbb C^2\mid z_1\neq z_2\}$
has fundamental group $\pi_1=\mathbb Z$, generated by the
loop $\gamma\colon t\mapsto
(z_2\,e^{2\pi it}+z_1(1-e^{2\pi it}),\,z_2)$
(the first point circles the second). The $\Sigma_2$-quotient
$\mathrm{Conf}_2(\mathbb C)$ has
$\pi_1(\mathrm{Conf}_2(\mathbb C))=\mathbb Z$, generated by
the image of~$\gamma$: the braid in which one point
winds around the other.

The residue differential $d_{\mathrm{res}}$ on the ordered
bar complex
${\Barch}_2^{\mathrm{ord}}=\bar\cA\otimes\bar\cA\otimes
C_\bullet(\mathrm{Conf}_2^{\mathrm{ord}})$
defines a flat connection $\nabla$ on the trivial bundle
over $\mathrm{Conf}_2^{\mathrm{ord}}$ with fibre
$\bar\cA\otimes\bar\cA$. Flatness is $d_{\mathrm{res}}^2=0$,
which holds because $d^2=0$ on ${\Barch}_2^{\mathrm{ord}}$.
The \emph{$R$-matrix} is the monodromy of $\nabla$ around
the generator~$\gamma$:
\begin{equation}\label{eq:r-matrix-monodromy}
R(z)
\;:=\;
\operatorname{Mon}_\gamma(\nabla)
\;=\;
\mathcal P\exp\!\Bigl(
\oint_\gamma\sum_{n\ge 0}r_n\,z^{-n-1}\,dz
\Bigr)
\;\in\;
\operatorname{End}(\bar\cA\otimes\bar\cA)[\![z^{-1}]\!],
\end{equation}
where $z=z_1-z_2$ is the relative coordinate,
$r_n=\sum_{I,J}c^K_{IJ;n}\,e_I\otimes e_J$ is the OPE matrix
at pole order~$n$, and $\mathcal P$ denotes path-ordering.
At leading order:
\begin{equation}\label{eq:r-matrix-leading}
R(z)\;=\;1+r_0/z+O(z^{-2}),
\qquad
r_0\;=\;\sum_{I,J}c^K_{IJ;0}\,e_I\otimes e_J,
\end{equation}
where $r_0$ is the \emph{classical $r$-matrix}: the matrix
of zeroth products (Lie brackets).

The $R$-matrix intertwines the two orderings:
\[
R(z)\colon
{\Barch}_2^{\mathrm{ord}}(a\otimes b;\,z_1,z_2)
\;\xrightarrow{\;\sim\;}
{\Barch}_2^{\mathrm{ord}}(b\otimes a;\,z_2,z_1).
\]
For the pole-free BD-commutative subclass of
$E_\infty$-chiral algebras, $R(z)=\tau$ (the flip): ordered
and unordered complexes are identified. For local
vertex algebras in the $E_\infty$ regime with OPE poles,
$R(z)\ne\tau$ but is derived from the local OPE.
For genuine $E_1$-algebras, $R(z)$ carries nontrivial spectral
dependence as independent input and is the fundamental new datum.
\end{construction}

\begin{proposition}[Descent identification;
\ClaimStatusProvedHere]
\label{prop:r-matrix-descent}
\index{descent!R-matrix!ordered to unordered}
Let $\cA$ be a strongly admissible $E_1$-chiral algebra with
ordered bar coalgebra $C=\Barchord(\cA)$ and
$R$-matrix~$R(z)$. Then the unordered bar complex is the
$R$-twisted $\Sigma_n$-descent:
\begin{equation}\label{eq:descent-identification}
\Barch(\cA)_n
\;\simeq\;
\bigl(\Barchord(\cA)_n\bigr)^{R\text{-}\Sigma_n},
\end{equation}
where $\Sigma_n$ acts by $\sigma_i\cdot(a_1\otimes\cdots\otimes a_n)
=(\tau_{i,i+1}\circ R_{i,i+1}(z_i{-}z_{i+1}))\,
(a_1\otimes\cdots\otimes a_n)$
on the adjacent transposition $(i,\,i{+}1)$. The
underlying $B_n$-representation is well-defined if and only
if the $R$-matrix satisfies the Yang--Baxter
equation~\textup{(Proposition~\ref{prop:ybe-from-d-squared})};
it factors through $\Sigma_n$ if and only if~$R$
satisfies \emph{strong unitarity}
$R_{12}(z)\,R_{21}(-z)=\id$.
For the pole-free BD-commutative subclass of
$E_\infty$-chiral algebras, $R(z)=1$ and the descent
reduces to ordinary $\Sigma_n$-coinvariants; for local
vertex algebras in the $E_\infty$ regime with OPE poles, the descent is
$R$-twisted (nontrivially, though $R(z)$ is derived from
the local OPE).
\end{proposition}

\begin{proof}
The covering $\pi_n$ presents $\mathrm{Conf}_n(X)$ as the
orbit space of the free $\Sigma_n$-action on
$\mathrm{Conf}_n^{\mathrm{ord}}(X)$. The monodromy of the
flat connection $\nabla=d+d_{\mathrm{res}}$ along the lift of a
generator $\sigma_i\in\pi_1(\mathrm{Conf}_n)$ decomposes into
the permutation $\tau_{i,i+1}$ of tensor slots and the
parallel transport $R_{i,i+1}(z_i-z_{i+1})$ of the fibre.
The composition
$\rho_R(\sigma_i)=\tau_{i,i+1}\circ R_{i,i+1}$ defines a
representation $\rho_R\colon B_n\to\operatorname{GL}(\bar\cA^{\otimes n})$
of the braid group: the braid relation
$\sigma_i\sigma_{i+1}\sigma_i
=\sigma_{i+1}\sigma_i\sigma_{i+1}$ is satisfied because it
reduces to the Yang--Baxter equation on~$R$
(Proposition~\ref{prop:ybe-from-d-squared}).

For $\rho_R$ to factor through the quotient
$\Sigma_n=B_n/\langle\sigma_i^2\rangle$, we need the
additional relation $\rho_R(\sigma_i)^2=\id$, i.e.\
strong unitarity $R_{12}(z)\,R_{21}(-z)=\id$. This holds
for $E_\infty$-chiral algebras (including those with OPE
poles) by the following algebraic argument. The OPE
residues $r_{IJ}^K$ of a local vertex algebra satisfy
$r_{12}=r_{21}$ (the Casimir element
$\Omega=\sum r_n\in(\bar\cA\otimes\bar\cA)[\![z^{-1}]\!]$
is symmetric under exchange of tensor factors). Path
reversal on $\mathrm{Conf}_2^{\mathrm{ord}}(\mathbb C)$ sends
$z\mapsto -z$ and exchanges the two sheets, so the parallel
transport along the reversed path is $R_{21}(-z)$. Since the
connection form is symmetric ($r_{12}=r_{21}$), path reversal
inverts the holonomy: $R_{21}(-z)=R_{12}(z)^{-1}$.
Hence $R_{12}(z)\,R_{21}(-z)=\id$.

The result is a genuine $\Sigma_n$-action $\rho_R$ on the
ordered bar complex. The descended local system is the
equaliser of~$\rho_R$ with the identity,
which is~\eqref{eq:descent-identification}.
\end{proof}

\begin{remark}[The $R$-matrix as Maurer--Cartan element]
\label{rem:r-matrix-mc}
\index{R-matrix!Maurer--Cartan interpretation}
The classical $r$-matrix
$r_0\in\operatorname{End}(\bar\cA\otimes\bar\cA)$ is a
Maurer--Cartan element in the Gerstenhaber--Schack dg~Lie
algebra $C^\bullet_{\mathrm{GS}}(\cA,\cA)$ of bilinear
operations, with the Gerstenhaber bracket as Lie bracket
and the bar differential as $d$. The Yang--Baxter equation
on~$R(z)$ is the exponentiated MC equation; the
classical Yang--Baxter equation on~$r_0$
(Corollary~\ref{cor:classical-ybe}) is the linearised MC
equation. The $R$-matrix is the descent datum because the
convolution dg~Lie algebra governing $\Sigma_n$-equivariant
structures on local systems is precisely the Gerstenhaber
complex, and MC elements in it parametrise the nontrivial
descents.
\end{remark}

\begin{proposition}[Yang--Baxter from $d^2=0$ on degree~$3$;
\ClaimStatusProvedHere]
\label{prop:ybe-from-d-squared}
\index{Yang--Baxter equation!from d squared zero@from $d^2=0$|textbf}
The identity $d^2=0$ on
$\Barchord_3(\cA)$ is equivalent to the
Yang--Baxter equation:
\begin{equation}\label{eq:ybe-from-bar}
R_{12}(z_{12})\,R_{13}(z_{13})\,R_{23}(z_{23})
\;=\;
R_{23}(z_{23})\,R_{13}(z_{13})\,R_{12}(z_{12}),
\end{equation}
where $z_{ij}=z_i-z_j$.
\end{proposition}

\begin{proof}
We give two proofs: one topological, one combinatorial.

\emph{Topological proof.}
The pure braid group
$P_3=\pi_1(\mathrm{Conf}_3^{\mathrm{ord}}(\mathbb C))$
has three standard generators
$\gamma_{12}$, $\gamma_{13}$, $\gamma_{23}$, where
$\gamma_{ij}$ is the loop in which $z_i$ winds around $z_j$
with all other points fixed.
The defining relation of~$P_3$ is
$\gamma_{12}\gamma_{13}\gamma_{23}
=\gamma_{23}\gamma_{13}\gamma_{12}$.
The flat connection $\nabla$ on the ordered bar complex
induces a monodromy representation
$\rho\colon P_3\to\operatorname{GL}(\bar\cA^{\otimes 3})$
sending $\gamma_{ij}\mapsto R_{ij}(z_{ij})$. The image
of the braid relation is~\eqref{eq:ybe-from-bar}.

\emph{Combinatorial proof via boundary faces.}
The ordered FM compactification
$\overline{\mathrm{FM}}_3^{\mathrm{ord}}(\mathbb C)$
has two codimension-one boundary faces:
$D_{1,2}=\{z_1=z_2\}$ (first pair collides) and
$D_{2,3}=\{z_2=z_3\}$ (second pair collides). The face
$D_{1,3}$ is not codimension-one in the \emph{ordered}
compactification: the first and third points can only
collide after one of the consecutive collisions, so
$D_{1,3}$ lives in codimension~$2$.
The unique codimension-two stratum is
$D_{1,2}\cap D_{2,3}=\{z_1=z_2=z_3\}$ (the triple collision).

The operator $d^2$ on ${\Barch}_3^{\mathrm{ord}}$ composes
the residue differential with itself. On a typical element
$s^{-1}a\otimes s^{-1}b\otimes s^{-1}c$, the first
application of $d$ produces two terms
(equation~\eqref{eq:ordered-bar-d3}), one from each face.
The second application produces terms supported on the
codimension-two stratum. There are exactly two pathways
to the triple collision:
\begin{enumerate}[label=(\roman*)]
\item $D_{1,2}$ then $D_{2,3}$: collapse $z_1\to z_2$
 (extracting $a_{(0)}b$), then the composite collides
 with $z_3$ (extracting $(a_{(0)}b)_{(0)}c$).
 The monodromy of the two-step path is
 $R_{12}(z_{12})\,R_{13}(z_{13})$.
\item $D_{2,3}$ then $D_{1,2}$: collapse $z_2\to z_3$
 first (extracting $b_{(0)}c$), then $z_1$ collides with
 the composite (extracting $a_{(0)}(b_{(0)}c)$).
 The monodromy is
 $R_{23}(z_{23})\,R_{13}(z_{13})$.
\end{enumerate}
The algebraic part of $d^2$ vanishes by associativity of
the zeroth product ($E_1$ axiom). The monodromy part
gives the Yang--Baxter equation: the two sides
of~\eqref{eq:ybe-from-bar} are the two pathways, and their
equality is $d^2=0$.
\end{proof}

\begin{corollary}[Classical Yang--Baxter equation;
\ClaimStatusProvedHere]
\label{cor:classical-ybe}
\index{Yang--Baxter equation!classical}
Writing $R(z)=1+r(z)+O(r^2)$ with
$r(z)=r_0/z+r_1/z^2+\cdots$, the linearisation of
the YBE~\eqref{eq:ybe-from-bar} is the classical
Yang--Baxter equation:
\begin{equation}\label{eq:cybe}
[r_{12}(z_{12}),\,r_{13}(z_{13})]
+[r_{12}(z_{12}),\,r_{23}(z_{23})]
+[r_{13}(z_{13}),\,r_{23}(z_{23})]
\;=\;0.
\end{equation}
At leading order ($r(z)=r_0/z$), the partial-fraction
identity
$\tfrac{1}{z_{12}z_{13}}+\tfrac{1}{z_{21}z_{23}}
+\tfrac{1}{z_{31}z_{32}}=0$
converts~\eqref{eq:cybe} into the Jacobi identity for the
Lie bracket $[a,b]=r_0(a\otimes b)-r_0(b\otimes a)$. The
classical $r$-matrix $r_0$ is thus a Maurer--Cartan element
in the convolution dg~Lie algebra
$\operatorname{Hom}(\bar\cA^{\otimes 2},\bar\cA)$ with the
Gerstenhaber bracket.
\end{corollary}

\begin{remark}[$\Sigma_n$-representations and the Eulerian weight
decomposition]
\label{rem:sn-eulerian-weight}
\index{Eulerian weight decomposition!on the bar complex|textbf}
\index{Sigma_n representation@$\Sigma_n$-representation!on bar complex}
\index{Harrison complex!Eulerian weight}
The ordered bar complex
$\Barchord_n(\cA) = (s^{-1}\bar\cA)^{\otimes n}$ carries a
$\Sigma_n$-action via Koszul signs: the transposition
$(i,\,i{+}1)$ acts on $s^{-1}a_i \otimes s^{-1}a_{i+1}$ by
\[
\sigma_i \cdot (s^{-1}a_i \otimes s^{-1}a_{i+1})
\;=\;
(-1)^{|s^{-1}a_i|\,|s^{-1}a_{i+1}|}\,
s^{-1}a_{i+1} \otimes s^{-1}a_i,
\]
where $|s^{-1}a| = |a| - 1$ is the desuspended degree.
For a single generator of conformal weight~$h$:
$|s^{-1}a| = h - 1$, and the sign
$(-1)^{(h-1)^2} = (-1)^{h-1}$ determines whether $\Sigma_n$
acts by the trivial or sign representation on the diagonal
$(s^{-1}a)^{\otimes n}$.

The Poincar\'e--Birkhoff--Witt decomposition
$\Sym^c(V) = \bigoplus_{j \geq 1} e_j \cdot \Sym^c(V)$
by the Eulerian idempotents $\{e_j\}$ in $\mathbb{Q}[\Sigma_n]$
organises this representation-theoretic structure.
At degree~$2$: $e_1 = \tfrac{1}{2}(\id - \sigma)$
(Harrison/Lie${}^c$ projection), $e_2 = \tfrac{1}{2}(\id + \sigma)$
(symmetric complement).
On the sign representation,
$e_1$ acts as~$1$ and $e_2$ as~$0$; on the trivial
representation, $e_1$ acts as~$0$ and $e_2$ as~$1$.

The modular characteristic $\kappa(\cA)$ lives at degree~$2$
and therefore has a definite Eulerian weight depending on the
desuspended parity:
\begin{center}
\renewcommand{\arraystretch}{1.2}
\begin{tabular}{lcccl}
\textbf{Family}
& $h$
& $|s^{-1}a|$
& \textbf{Eulerian wt of $\kappa$}
& \textbf{Harrison visibility} \\
\hline
Heisenberg ($J$)
& $1$ & $0$ (even)
& weight~$2$ (symmetric)
& invisible \\
Virasoro ($T$)
& $2$ & $1$ (odd)
& weight~$1$ (Harrison)
& visible \\
$\cW_3$ ($T, W$)
& $2, 3$ & $1, 2$ (mixed)
& decomposes across both
& partially visible
\end{tabular}
\end{center}
For the multi-generator case ($\cW_3$): the $TT$-block carries
the sign representation (both factors have odd desuspended
degree~$1$) and lives in Harrison weight~$1$; the
$WW$-block carries the trivial representation (even desuspended
degree~$2$) and lives in weight~$2$; the mixed $TW$/$WT$-block
exchanges with sign $(-1)^{1 \cdot 2} = -1$ and decomposes into
one Harrison and one symmetric component.

The averaging map $\mathrm{av}\colon \mathfrak{g}_\cA^{E_1} \to \gAmod$ of
the $\Eone$-primacy thesis
(Volume~I, Principle~\textup{\ref*{V1-princ:e1-primacy}})
is the $\Sigma_n$-coinvariant projection
$T^c(s^{-1}\bar\cA) \to \Sym^c(s^{-1}\bar\cA)$.
The Eulerian weight filtration refines this projection:
the Harrison (weight-$1$) component
$e_1 \cdot (s^{-1}\bar\cA)^{\otimes n}$ is the Lie-coalgebraic
part detected by the Francis--Gaitsgory bar $\barB^{\mathrm{FG}}$,
while the higher-weight components carry the surplus information
that the full $\barB^{\Sigma}$ retains over $\barB^{\mathrm{FG}}$.
At degree~$2$, $\kappa$ for the Virasoro algebra is entirely
Harrison-visible, while $\kappa$ for the Heisenberg is entirely
Harrison-invisible: this is the representation-theoretic
explanation for why the FG bar complex sees the full Virasoro
curvature but misses the Heisenberg level.

Computational verification:
\texttt{compute/lib/eulerian\_weight\_mc\_engine.py}
($67$~tests).
\end{remark}

\subsection{The ordered bar complex and its Kac--Moody Yangian specialization}
\label{subsec:dg-shifted-yangian-from-bar}

The $R$-matrix governs descent from ordered to unordered; the Yang--Baxter equation is $d^2=0$. The $E_\infty$-chiral Koszul dual $\cA^!_{\mathrm{ch}}$ is another chiral algebra on the same curve. The ordered ($E_1$) Koszul dual is an associative algebra with spectral parameter, living on the affine line of spectral parameters rather than on the curve. For affine Kac--Moody input, this algebra is the Yangian.

\begin{construction}[Open-colour dual from the ordered bar
complex; Kac--Moody Yangian specialization; \ClaimStatusProvedHere]
\label{constr:dg-shifted-yangian-from-bar}
\index{dg-shifted Yangian!from ordered bar complex|textbf}
The open-colour Koszul dual of an $E_1$-chiral algebra $\cA$ on the
affine line is the linear dual of the cohomology of the
ordered bar complex:
\[
\cA^!_{\mathrm{line}}
\;=\;
H^*\!\bigl(\Barchord(\cA)\bigr)^\vee.
\]
This is the \emph{small} (classical) open-colour Koszul dual: a
graded algebra whose multiplication is the dual of the
deconcatenation coproduct descended to cohomology. On the
chirally Koszul locus the small dual agrees with the
\emph{homotopy} open-colour Koszul dual
$\Omega^{\mathrm{ord}}(\Barchord(\cA))$,
which is a full $\Ainf$ algebra. Off the Koszul locus
(e.g.\ for $\mathrm{Vir}_c$), the homotopy version carries
genuinely higher $\Ainf$ operations $m_k$ ($k \ge 3$) that
the small dual does not encode.

For $\cA=\widehat{\mathfrak g}_k$ (affine Kac--Moody at
level~$k$), the open-colour Koszul dual is the finite-type Yangian:
$\cA^!_{\mathrm{line}}=Y_\hbar(\mathfrak g)$ at deformation parameter
$\hbar=1/(k+h^\vee)$. The input is the \emph{affine} Kac--Moody
algebra $\widehat{\fg}_k$; the output is the \emph{finite-type}
Yangian $Y(\fg)$. This identification is the
associative chiral analogue of the commutative identification
$\cA^!_{\mathrm{ch}}=\widehat{\mathfrak g}_{-k-2h^\vee}$
(Verdier-level duality): the ordered bar complex replaces
symmetric coinvariants with the braid-group--equivariant
structure, and the Yangian emerges as the algebra encoding
this equivariance.

\begin{remark}[Two-colour Koszul duality: why
$\cA^!_{\mathrm{ch}}$ and $\cA^!_{\mathrm{line}}$ differ]%
\label{rem:two-colour-architecture}%
The two-colour Koszul datum of~$\cA$
carries two bar-complex projections, each with its own
duality operation:
\begin{itemize}[nosep]
\item \emph{Closed colour.}
 The symmetric bar complex
 $\Barch(\cA) = \Barchord(\cA)_{\Sigma_n}$
 is a factorisation coalgebra on $\operatorname{Ran}(X)$.
 Verdier duality $\mathbb{D}_{\operatorname{Ran}}$
 produces a factorisation algebra; cobar recovers the
 chiral Koszul dual $\cA^!_{\mathrm{ch}}$, another
 chiral algebra on~$X$. For $\widehat{\fg}_k$:
 $\cA^!_{\mathrm{ch}} = \widehat{\fg}_{-k-2h^\vee}$.
 This is the Francis--Gaitsgory direction
 ($\mathrm{Com}^{\mathrm{ch}}\,!\,=
 \mathrm{Lie}^{\mathrm{ch}}$).
\item \emph{Open colour.}
 The ordered bar complex
 $\Barchord(\cA)$ on
 $\Conf_n^{\mathrm{ord}}(X)$ has no $\Sigma_n$-quotient.
 Linear duality $H^*(-)^\vee$ produces an associative
 algebra with spectral parameter:
 $\cA^!_{\mathrm{line}}$. For $\widehat{\fg}_k$:
 $\cA^!_{\mathrm{line}} = Y_\hbar(\fg)$.
 This is the $E_1$-chiral direction: Koszul duality
 for $\SCchtop$ exchanges the closed colour
 ($\mathsf{Com} \leftrightarrow \mathsf{Lie}$) while
 preserving the open colour ($\mathsf{Ass} = \mathsf{Ass}$)
 (Proposition~\ref{thm:SC-self-duality}), so the
 $E_1$-algebra structure survives on the dual.
\end{itemize}
The $R$-matrix is the cross-colour datum:
$\Barch(\cA)_n \simeq
(\Barchord(\cA)_n)^{R\text{-}\Sigma_n}$
(Proposition~\ref{prop:r-matrix-descent}). For
the pole-free BD-commutative subclass of
$E_\infty$-chiral algebras, $R(z) = \tau$ and both
projections carry equivalent information. For
$E_\infty$-chiral algebras with OPE poles
(all interesting vertex algebras: affine Kac--Moody,
Virasoro, Heisenberg),
$R(z)$ carries nontrivial spectral dependence
(derived from the local OPE via analytic
continuation), and $\cA^!_{\mathrm{line}}$ carries
strictly more structure than $\cA^!_{\mathrm{ch}}$:
the spectral parameter~$u$ is the fingerprint of the
chiral data retained by the ordered bar differential.
Theorem~\ref{thm:FG-shadow} gives the controlled
interpolation between the two projections via the
pole-order spectral sequence.

The two-colour duality is summarised for the standard families
in Table~\ref{tab:two-colour-koszul-duals}.
\end{remark}

\begin{table}[ht]
\centering
\renewcommand{\arraystretch}{1.2}
\caption{Two-colour Koszul duals for the standard families.
 The closed colour $\cA^!_{\mathrm{ch}}$ is the Verdier dual
 (chiral Koszul dual on $\operatorname{Ran}(X)$);
 the open colour $\cA^!_{\mathrm{line}}$ is the ordered
 Koszul dual (linear dual of ordered bar cohomology,
 an algebra with spectral parameter).
 The $R$-matrix column records the descent datum
 $\Barchord\to\Barch$; the class column
 is the shadow depth class
 ($\mathbf{G}$ = pole-free,
 $\mathbf{L}$ = finite shadow depth,
 $\mathbf{M}$ = infinite shadow depth).}
\label{tab:two-colour-koszul-duals}
\index{Koszul duality!two-colour comparison table|textbf}
\[
\renewcommand{\arraystretch}{1.25}
\begin{array}{lllll}
\cA
 & \cA^!_{\mathrm{ch}}
 & \cA^!_{\mathrm{line}}
 & R(z)
 & \text{Class} \\
\hline
\cH_k
 & \mathrm{Sym}^{\mathrm{ch}}(V^*)
 & \Bbbk[t]\text{ (polynomial)}
 & z^k\;\text{(scalar)}
 & \mathbf{L} \\
\widehat{\mathfrak{sl}}_2
 & \widehat{\mathfrak{sl}}_{2,\,-k-4}
 & Y_{1/(k+2)}(\mathfrak{sl}_2)
 & 1+\Omega/((k+2)z)+O(z^{-2})
 & \mathbf{L} \\
\widehat{\mathfrak{sl}}_3
 & \widehat{\mathfrak{sl}}_{3,\,-k-6}
 & Y_{1/(k+3)}(\mathfrak{sl}_3)
 & 1+\Omega/((k+3)z)+O(z^{-2})
 & \mathbf{L} \\
\mathrm{Vir}_c
 & \mathrm{Vir}_{26-c}
 & Y^{\Ainf}_\hbar(\mathrm{Vir})
 & 1+(c/2)/z^3+2T/z
 & \mathbf{M} \\
\mathcal{W}_{3,c}
 & \mathcal{W}_{3,\,100-c}
 & Y^{\Ainf}_\hbar(\mathcal{W}_3)
 & 1+O(z^{-3})
 & \mathbf{M} \\
\beta\gamma_\lambda
 & bc_\lambda
 & \Bbbk[t]\text{ (polynomial)}
 & z^{-1}\;\text{(scalar)}
 & \mathbf{L}
\end{array}
\]
\end{table}
\noindent
The affine families ($\cH_k$, $\widehat{\mathfrak{sl}}_2$,
$\widehat{\mathfrak{sl}}_3$) share the pattern: the
closed-colour dual $\cA^!_{\mathrm{ch}}$ remains in the
same family at the reflected level $k \mapsto -k - 2h^\vee$
(Feigin--Frenkel), while the open-colour dual
$\cA^!_{\mathrm{line}}$ exits the chiral category
entirely and becomes a \emph{finite-type Yangian},
an associative algebra with spectral parameter~$u$,
not a chiral algebra on a curve. The central extension
(the double-pole residue) is invisible to the open colour
(Proposition~\ref{prop:open-colour-double-bar}).

For Virasoro and $\mathcal{W}_3$, the closed-colour dual
stays in the same family at the reflected central charge
($c \mapsto 26 - c$, resp.\ $c \mapsto 100 - c$), but
the open-colour dual $Y^{\Ainf}_\hbar$ is a genuinely
$\Ainf$ object: the higher operations $m_k$ ($k \ge 3$)
do not vanish, reflecting the infinite shadow depth of
class~$\mathbf{M}$. The $R$-matrix has cubic leading
pole (from the quartic OPE pole after $d\log$ absorption),
in contrast to the simple leading pole for the affine
families.

For $\beta\gamma$, the closed-colour dual is the $bc$
system (statistics exchange: $\mathrm{Sym}^{\mathrm{ch}}
\leftrightarrow \bigwedge^{\mathrm{ch}}$), and the
open-colour dual reduces to the polynomial ring $\Bbbk[t]$
(the free-field case: trivial $\Ainf$ structure, $m_k = 0$
for $k \ge 3$).

\begin{proposition}[Open-colour double bar recovers current algebra;
\ClaimStatusProvedHere]
\label{prop:open-colour-double-bar}%
\index{double bar!open-colour!central extension invisible}%
Let $\fg$ be a finite-dimensional simple Lie algebra,
$k\notin\{-h^\vee\}$, $\hbar=1/(k{+}h^\vee)$, and let
$Y_\hbar(\fg)=\cA^!_{\mathrm{line}}$ be the Yangian obtained
from $\widehat{\fg}_k$ by open-colour Koszul duality
\textup{(Theorem~\ref{thm:duality-involution})}.
Then the open-colour double bar
$\Barchord(Y_\hbar(\fg))$ recovers the
\emph{current algebra} $U(\fg[t])$
\textup{(}the universal enveloping algebra of the polynomial
loop algebra, \textbf{without} central extension\textup{)}:
\[
 \Barchord\bigl(Y_\hbar(\fg)\bigr)
 \;\xrightarrow{\;\sim\;}
 U(\fg[t]).
\]
The central extension
$k\,n\,\delta_{n+m,0}\in[\widehat{\fg}_k,\widehat{\fg}_k]$
is invisible to the open-colour double bar.
Recovery of the full affine algebra $\widehat{\fg}_k$
\textup{(}including the central term\textup{)} requires the
\emph{closed-colour} direction: the Verdier\slash Ran duality of
the full $\SCchtop$ Koszul duality
\textup{(Theorem~\ref{thm:duality-involution})}.
\end{proposition}

\begin{proof}
The Yangian spectral OPE in the Drinfeld presentation has the
schematic form
\[
 [E_i(u),\,F_j(v)]
 \;=\;
 \frac{\delta_{ij}}{u-v}\,H_i(v)
 \;+\;\text{(regular)},
 \qquad
 [H_i(u),\,E_j(v)]
 \;=\;
 \frac{A_{ij}}{u-v}\,E_j(v)
 \;+\;\text{(regular)},
\]
with only \emph{simple} poles in $u-v$. By the pole-order
analysis of the bar differential
(Proposition~\ref{prop:r-matrix-descent}), the degree-$2$
component of $d_{\mathrm{bar}}$ extracts only the
\emph{zeroth product} (the residue at the simple pole):
\[
 d_{\mathrm{bar}}(s^{-1}E_i\otimes s^{-1}F_j)
 \;=\;
 \delta_{ij}\,s^{-1}H_i,
 \qquad
 d_{\mathrm{bar}}(s^{-1}H_i\otimes s^{-1}E_j)
 \;=\;
 A_{ij}\,s^{-1}E_j.
\]
The central term $k\,n\,\delta_{n+m,0}$ in the affine
commutation relations arises from a \emph{double} pole in the
affine OPE $J^a(z)J^b(w)\sim k\,\kappa^{ab}/(z{-}w)^2+\cdots$\,.
The Yangian spectral OPE has no double pole: the level enters
only as the deformation parameter $\hbar=1/(k{+}h^\vee)$,
which scales the \emph{position} of simple poles (e.g.\
$E_i(u)E_i(v)$ has a pole at $u-v=\hbar$, shifted off the
collision diagonal) but does not produce a pole of higher
order at $u=v$. Therefore the bar differential sees no
central extension.

The $E$--$E$ pole at $u-v=\hbar$ is shifted off the collision
locus $u=v$, so
$d_{\mathrm{res}}(s^{-1}E_i\otimes s^{-1}E_i)=0$, correctly
reproducing $[e_i,e_i]=0$ in $\fg[t]$. At degree~$3$,
$d^2=0$ reduces to the classical Yang--Baxter equation for
$r(u)=P/u$, which holds by Jacobi.

For the full recovery: the closed-colour bar complex retains the symmetric bar
$\Barch(\cA)_n\simeq
(\Barchord(\cA)_n)^{R\text{-}\Sigma_n}$
(Proposition~\ref{prop:r-matrix-descent}), whose Verdier
dual on $\operatorname{Ran}(X)$ carries the factorisation
algebra structure encoding the double-pole residue. The
double Koszul duality of
Theorem~\ref{thm:duality-involution} therefore recovers
$\widehat{\fg}_k$ with its central extension.
\end{proof}

\emph{Degree-$1$ generators: the Drinfeld currents.}
The Yangian generators in the Drinfeld presentation
$(E_i(u),\,F_i(u),\,H_i(u))$ are the degree-$1$ bar
cohomology classes:
\begin{align*}
E_i(u)&\;=\;\sum_{r\ge 0}E_i^{(r)}\,u^{-r-1}
\;:\;
[\,s^{-1}e_i\,]\in H^1(\Barchord),\\
F_i(u)&\;=\;\sum_{r\ge 0}F_i^{(r)}\,u^{-r-1}
\;:\;
[\,s^{-1}f_i\,]\in H^1(\Barchord),\\
H_i(u)&\;=\;1+\sum_{r\ge 0}H_i^{(r)}\,u^{-r-1}
\;:\;
[\,s^{-1}h_i\,]\in H^1(\Barchord),
\end{align*}
where $u=1/z$ is the spectral parameter (inverse of the
curve coordinate). The leading-order generators
$E_i^{(0)}$, $F_i^{(0)}$, $H_i^{(0)}$ correspond to the
zeroth modes $e_{i(0)}$, $f_{i(0)}$, $h_{i(0)}$ of the
currents, and generate a copy of $U(\mathfrak g)$ inside
$Y_\hbar(\mathfrak g)$. The higher modes $E_i^{(r)}$
($r\ge 1$) come from the $r$-th pole OPE coefficients
$e_{i(r)}$; they encode the deformation of
$U(\mathfrak g[t])$ into $Y_\hbar(\mathfrak g)$.

\emph{Degree-$2$ relations: the Yangian commutation
relations.}
The bar differential on
$s^{-1}h_i\otimes s^{-1}e_j$ maps to $s^{-1}(h_{i(0)}e_j)
=A_{ij}\,s^{-1}e_j$, so the cohomology relation
$[H_i,E_j]=A_{ij}\,E_j$ at leading order in~$u$ is
the Cartan matrix relation. Higher-order terms in the
$u$-expansion produce the Yangian deformation:
\[
[H_i(u),\,E_j(v)]
\;=\;
\frac{A_{ij}}{u-v}\,E_j(v)
\;+\;\text{(higher poles)},
\]
where the higher poles come from OPE coefficients
$h_{i(n)}e_j$ ($n\ge 1$). The Serre relations arise from
degree-$3$ bar cohomology via the $A_\infty$ structure.

\emph{The RTT presentation.}
The transfer matrix
$T(u)=\sum_I e_I\otimes t^I(u)
\in\operatorname{End}(V)\otimes Y_\hbar(\mathfrak g)[\![u^{-1}]\!]$
(where $V$ is the defining representation of~$\mathfrak g$
and $t^I(u)$ are the degree-$1$ bar cohomology generating
functions) satisfies
\begin{equation}\label{eq:rtt-from-bar}
R(u{-}v)\,T_1(u)\,T_2(v)
\;=\;
T_2(v)\,T_1(u)\,R(u{-}v),
\end{equation}
where $R(u-v)$ is the $R$-matrix from
Construction~\ref{constr:r-matrix-monodromy}. The RTT
relation~\eqref{eq:rtt-from-bar} is equivalent to $d^2=0$
on the three-fold ordered bar complex
${\Barch}_3^{\mathrm{ord}}$: the $R$-matrix intertwines the
orderings of the first two tensor factors, and the transfer
matrix encodes the insertion of a bar cohomology class in
the third slot. The consistency of these three insertions
is the Yang--Baxter
equation~\eqref{eq:ybe-from-bar}
(Computation~\ref{comp:ordered-bar-sl2},
Proposition~\ref{prop:ybe-from-d-squared}).

\emph{For $\widehat{\mathfrak{sl}}_2$ at level~$k$.}
The $R$-matrix is $R(z)=1+\hbar\,k\,\Omega/z+O(z^{-2})$ with
$\hbar=1/(k+2)$ and
$\Omega=e\otimes f+f\otimes e+\tfrac12 h\otimes h$
(the Casimir), and the Yangian $Y_\hbar(\mathfrak{sl}_2)$
is generated by $E(u)$, $F(u)$, $H(u)$ with
the same $\hbar$. The RTT relation
$R(u{-}v)\,T_1(u)\,T_2(v)=T_2(v)\,T_1(u)\,R(u{-}v)$
with $T(u)=\bigl(\begin{smallmatrix}
H(u) & F(u) \\ E(u) & -H(u)
\end{smallmatrix}\bigr)$
reproduces all Yangian relations from bar complex
associativity.

\emph{For $\widehat{\mathfrak{sl}}_3$ at level~$k$:
adjacent roots and the coefficient $\frac12$.}
In $\mathfrak{gl}_3$ matrix units:
$e_1=E_{12}$, $e_2=E_{23}$, $e_{12}=E_{13}$,
$f_1=E_{21}$, $f_2=E_{32}$, $f_{12}=E_{31}$.
The degree-$3$ ordered bar complex exhibits
\emph{triangle sectors}:
\[
T^L_{12}=E_{13}^{(1)}\,E_{21}^{(2)}\,E_{32}^{(3)},
\quad
T^L_{21}=E_{13}^{(1)}\,E_{32}^{(2)}\,E_{21}^{(3)};
\qquad
T^R_{12}=E_{12}^{(1)}\,E_{23}^{(2)}\,E_{31}^{(3)},
\quad
T^R_{21}=E_{23}^{(1)}\,E_{12}^{(2)}\,E_{31}^{(3)}.
\]
The triangle coefficient is
$\lambda^L=\lambda^R=c_1 c_2/(2c_{12})=1/2=\int_0^1(1-t)\,dt$
(the $n=2$ beta integral). The Serre relation follows from
root-space vanishing:
$2\alpha_1+\alpha_2\notin\Phi(\mathfrak{sl}_3)
\implies (\mathfrak{sl}_3)_{2\alpha_1+\alpha_2}=0
\implies [e_1,[e_1,e_2]]=0$.
\end{construction}

\subsection{The Euler-eta identity and bar cohomology concentration}
\label{subsec:euler-eta-identity}

The low-degree computations above display the ordered bar complex entry by entry. The generating function for bar dimensions across all tensor degrees and conformal weights has a closed form that connects the bar complex to Dedekind $\eta$-functions.

The ordered bar complex has a natural bigrading by
\emph{tensor degree}~$n$ (number of bar slots) and
\emph{conformal weight}~$w$ (total weight of the
insertions). The Euler characteristic in this bigrading
is computable in closed form for all standard families
and, for affine Kac--Moody algebras at generic
non-critical level, reveals a connection to Dedekind
$\eta$-functions and modular forms.

\begin{theorem}[Euler-eta identity; \ClaimStatusProvedHere]
\label{thm:euler-eta-identity}
\index{Euler-eta identity|textbf}
\index{bar complex!ordered!Euler characteristic}
Let $\cA$ be a strongly admissible $E_1$-chiral algebra with
ordered bar complex $\Barchord(\cA)$. Define
the \emph{graded bar Euler characteristic} by
\[
\chi(\Barchord;\,q)
\;:=\;
\sum_{n\ge 1}(-1)^n\dim \Barchord_{n,w}\;q^w,
\]
where $\Barchord_{n,w}$ is the weight-$w$,
tensor-degree-$n$ component, and the \emph{chiral character}
by $\operatorname{ch}(\cA;\,q):=\sum_{w\ge 0}\dim\cA_w\;q^w$.
Then
\begin{equation}\label{eq:euler-eta-master}
\chi(\Barchord;\,q)
\;=\;
-1+\frac{1}{\operatorname{ch}(\cA;\,q)}.
\end{equation}
\end{theorem}

\begin{proof}
The ordered bar complex at tensor degree~$n$ has underlying
graded vector space $(s^{-1}\bar\cA)^{\otimes n}$. The
generating function for the total graded dimension is
\[
\sum_{n\ge 0}
\dim(s^{-1}\bar\cA)^{\otimes n}\;t^n
\;=\;
\frac{1}{1-\operatorname{ch}(\bar\cA;\,q)\,t}
\;=\;
\frac{1}{1-(\operatorname{ch}(\cA;\,q)-1)\,t}.
\]
Setting $t=-1$ (the Euler characteristic) and subtracting
the $n=0$ term:
\[
\chi(\Barchord;\,q)
\;=\;
\frac{1}{1+(\operatorname{ch}(\cA;\,q)-1)}-1
\;=\;
\frac{1}{\operatorname{ch}(\cA;\,q)}-1.
\qedhere
\]
\end{proof}

\begin{corollary}[Specializations to standard families;
\ClaimStatusProvedHere]
\label{cor:euler-eta-specializations}
\index{Euler-eta identity!specializations}
\begin{enumerate}[label=\textup{(\roman*)}]
\item \emph{Affine Kac--Moody $V_k(\fg)$, $\fg$ simple of
 rank~$r$ and dimension~$d$, at generic non-critical level
 $k\ne -h^\vee$}:
 the vacuum character is
 $\operatorname{ch}(V_k(\fg);\,q)
 =\prod_{n\ge 1}(1-q^n)^{-d}$
 (the universal enveloping algebra character, valid in
 the vacuum sector at non-critical level), so
 \begin{equation}\label{eq:euler-eta-affine}
 \chi(\Barchord;\,q)
 \;=\;
 -1+\eta(q)^d,
 \end{equation}
 where $\eta(q)=q^{1/24}\prod_{n\ge 1}(1-q^n)$ is the
 Dedekind eta function (up to the standard $q^{1/24}$
 prefactor, which we suppress in the weight-graded
 convention $q=e^{2\pi i\tau}$ with
 $\cA_0=\Bbbk$). At the critical level $k=-h^\vee$,
 the vacuum character differs (the Feigin--Frenkel
 centre enlarges $V_{-h^\vee}(\fg)$) and the
 $\eta^d$ specialization does not apply; use the
 master formula~\eqref{eq:euler-eta-master} instead.
\item \emph{Heisenberg $\cH_k$} ($d=1$):
 $\chi=-1+\prod_{n\ge 1}(1-q^n)$. The coefficients
 are the \emph{Euler pentagonal numbers}: by
 Euler's pentagonal theorem,
 $\prod_{n\ge 1}(1-q^n)
 =\sum_{m\in\mathbb Z}(-1)^m q^{m(3m-1)/2}$,
 so $\chi_w=(-1)^m$ when $w=m(3m{-}1)/2$ for some
 $m\in\mathbb Z$, and $\chi_w=0$ otherwise.
\item \emph{Affine $\widehat{\mathfrak{sl}}_2$ at generic
 non-critical level $k\ne -2$} ($d=3$):
 $\chi=-1+\eta(q)^3$. By the Jacobi triple product,
 $\eta(q)^3
 =\sum_{n\ge 0}(-1)^n(2n{+}1)q^{n(n+1)/2}$,
 so the Euler characteristic is a signed sum of odd
 integers at triangular-number weights.
\item \emph{Virasoro $\mathrm{Vir}_c$}:
 $\operatorname{ch}(\mathrm{Vir}_c;\,q)
 =\prod_{n\ge 2}(1-q^n)^{-1}=(1{-}q)/\eta(q)$
 (up to $q^{1/24}$), so
 \begin{equation}\label{eq:euler-eta-virasoro}
 \chi(\Barchord;\,q)
 \;=\;
 -1+\frac{\eta(q)}{1-q}.
 \end{equation}
 The Virasoro Euler characteristic is the ratio of
 the pentagonal series to the geometric series: partial
 sums of the pentagonal theorem.
\end{enumerate}
\end{corollary}

\begin{proof}
Direct substitution of the known characters into
equation~\eqref{eq:euler-eta-master}. For~(ii), Euler's
pentagonal theorem is
$\prod_{n\ge 1}(1-q^n)
=\sum_{m=-\infty}^{\infty}(-1)^m q^{m(3m-1)/2}$.
For~(iii), the Jacobi triple product gives
$\prod_{n\ge 1}(1-q^n)^3
=\sum_{n=0}^{\infty}(-1)^n(2n+1)q^{n(n+1)/2}$.
For~(iv), $\operatorname{ch}(\mathrm{Vir}_c;\,q)
=\prod_{n\ge 2}(1-q^n)^{-1}$, so
$1/\operatorname{ch}=\prod_{n\ge 2}(1-q^n)
=\bigl[\prod_{n\ge 1}(1-q^n)\bigr]/(1-q)
=\eta(q)/(1-q)$.
\end{proof}

\begin{remark}[Scope of the $\eta^d$ specialization]
\label{rem:euler-eta-scope}
The \emph{universal} formula
$\chi=-1+1/\operatorname{ch}(\cA;\,q)$
(Theorem~\ref{thm:euler-eta-identity}) holds for all
strongly admissible chiral algebras. The
$\eta(q)^d$ specialization~\eqref{eq:euler-eta-affine}
is specific to \emph{affine Kac--Moody} vacuum modules
at \emph{generic non-critical level}: it uses the
character $\operatorname{ch}(V_k(\fg);\,q)
=\prod_{n\ge 1}(1-q^n)^{-d}$, which is the PBW
character of the vacuum sector. For~$\mathcal W$-algebras
obtained by Drinfeld--Sokolov reduction, the character
changes: the Virasoro algebra
($\mathcal W_2 = \mathrm{Vir}_c$) has
$\chi = -1+\eta(q)/(1{-}q)$
(Corollary~\ref{cor:euler-eta-specializations}(iv)),
and $\mathcal W_3$ at generic level gives
$\chi = -1+\prod_{n\ge 2}(1-q^n)\cdot
\prod_{m\ge 3}(1-q^m)$
(two generators at weights~$2$ and~$3$).
These are \emph{not} pure $\eta$-powers.
\end{remark}

\begin{computation}[Weight-by-weight Euler characteristics;
\ClaimStatusProvedHere]
\label{comp:euler-chi-table}
\index{Euler-eta identity!weight table}
The Euler characteristic $\chi_w$ at weight~$w$ for
$w=1,\ldots,20$ across the four standard families.
Let $p(w)$ denote the number of partitions of~$w$, and
$p_{\ge 2}(w)$ the number of partitions of~$w$ into
parts~$\ge 2$.
\[
\renewcommand{\arraystretch}{1.05}
\begin{array}{r|rrrr}
w & \cH_k & \widehat{\mathfrak{sl}}_2
 & \mathrm{Vir}_c & V_k(\mathfrak{sl}_3)
 \\
\hline
 1 & -1 & -3 & 0 & -8 \\
 2 & -1 & 0 & -1 & 20 \\
 3 & 0 & 5 & -1 & 0 \\
 4 & 0 & 0 & -1 & -70 \\
 5 & 1 & 0 & 0 & 64 \\
 6 & 0 & -7 & 0 & 56 \\
 7 & 1 & 0 & 1 & 0 \\
 8 & 0 & 0 & 1 & -125 \\
 9 & 0 & 0 & 1 & -160 \\
10 & 0 & 9 & 1 & 308 \\
11 & 0 & 0 & 1 & 0 \\
12 & -1 & 0 & 0 & 110 \\
13 & 0 & 0 & 0 & 0 \\
14 & 0 & 0 & 0 & -520 \\
15 & -1 & -11 & -1 & 0 \\
16 & 0 & 0 & -1 & 57 \\
17 & 0 & 0 & -1 & 560 \\
18 & 0 & 0 & -1 & 0 \\
19 & 0 & 0 & -1 & 0 \\
20 & 0 & 0 & -1 & 182
\end{array}
\]
\emph{Verification.}
Heisenberg: $1/\mathrm{ch}=\prod_{n\ge 1}(1-q^n)$;
by Euler's pentagonal theorem
$=\sum_{k\in\mathbf{Z}}(-1)^k q^{k(3k-1)/2}$,
so $\chi_1=-1$ ($k\!=\!1$), $\chi_2=-1$ ($k\!=\!{-}1$),
$\chi_5=+1$ ($k\!=\!{-}2$), $\chi_7=+1$ ($k\!=\!2$), etc.
$\widehat{\mathfrak{sl}}_2$:
$1/\mathrm{ch}=\prod(1-q^n)^3$;
by Jacobi's triple product
$=\sum_{m\ge 0}(-1)^m(2m+1)\,q^{m(m+1)/2}$,
so $\chi_1=-3$ ($m\!=\!1$), $\chi_3=+5$ ($m\!=\!2$),
$\chi_6=-7$ ($m\!=\!3$), $\chi_{10}=+9$ ($m\!=\!4$), etc.
Virasoro:
$1/\mathrm{ch}=\prod_{n\ge 2}(1-q^n)
=\prod_{n\ge 1}(1-q^n)/(1-q)$, i.e.\ the
cumulative partial sums of the Euler function.
$V_k(\mathfrak{sl}_3)$:
$1/\mathrm{ch}=\prod(1-q^n)^8$; coefficients grow
with the $8$-dimensional generator space.

The Heisenberg column is pentagonal: $\chi_w\ne 0$ only at
$w\in\{1,2,5,7,12,15,22,26,\ldots\}=\{m(3m{\pm}1)/2\}_{m\ge 1}$,
with alternating signs. The $\widehat{\mathfrak{sl}}_2$
column is Jacobi: nonzero only at triangular numbers
$w=n(n{+}1)/2$, with value $(-1)^n(2n{+}1)$ (the first
corrections are $3$-fold pentagonal interactions). The
Virasoro column exhibits plateaux of constant sign,
partial sums of the pentagonal series, reflecting the
single-generator structure with no weight-$1$ field.
The $\mathfrak{sl}_3$ column grows rapidly:
$\eta(q)^8$ has large coefficients from the $8$-dimensional
generator space.
\end{computation}

\begin{theorem}[Bar cohomology concentration;
\ClaimStatusProvedHere]
\label{thm:bar-cohomology-concentration}
\index{bar complex!ordered!cohomology concentration|textbf}
Let $\cA$ be a strongly admissible $E_1$-chiral algebra that
is \emph{chirally Koszul}: the bar--cobar equivalence
$\Cobar\Barch(\cA)\xrightarrow{\sim}\cA$ induces a
quasi-isomorphism on the bar complex whose cohomology
is concentrated in tensor degrees~$0$ and~$1$. Then:
\begin{enumerate}[label=\textup{(\roman*)}]
\item $H^0(\Barchord(\cA))=\Bbbk$ (the
 coaugmentation), $H^1(\Barchord(\cA))=V$
 (the strong generating space of~$\cA$), and
 $H^n(\Barchord(\cA))=0$ for all $n\ge 2$.
\item The Poincar\'e polynomial of the bar cohomology is
 \begin{equation}\label{eq:poincare-bar}
 P(t)\;=\;1+(\dim V)\cdot t.
 \end{equation}
\item Weight-refined: if the generators have conformal
 weights $s_1\le s_2\le\cdots\le s_N$ (where $N=\dim V$),
 then
 \begin{equation}\label{eq:poincare-bar-refined}
 P(t,q)\;=\;1+t\sum_{i=1}^N q^{s_i}.
 \end{equation}
\item The \emph{depth spectral sequence}
 $E_1^{p,*}=H_*^{\mathrm{Lie}}(\fg;\,\Bbbk)$
 (Lie algebra homology of the Lie algebra underlying the
 zeroth products) converges to $H^*(\Barchord)$.
 For affine Kac--Moody algebras, the spectral sequence
 collapses at $E_2$. For the Virasoro algebra, higher
 differentials are nonzero through $E_4$ (reflecting the
 quartic OPE pole order).
\end{enumerate}
\end{theorem}

\begin{proof}
Statement~(i): Koszulness means the bar cohomology
$H^*(\Barchord(\cA))$ is concentrated
in degrees~$0$ and~$1$. Degree~$0$ is the coaugmentation
$\Bbbk$ by convention. Degree~$1$ consists of the
indecomposable bar elements
$[s^{-1}a]\in\Barchord_1=s^{-1}\bar\cA$
that are not in the image of $d$ from degree~$2$. The bar
differential $d\colon\Barchord_2\to
\Barchord_1$ extracts the zeroth product
$a_{(0)}b$; its image is the space of decomposable
elements. The quotient is $V=\bar\cA/(\bar\cA\cdot\bar\cA)$,
the strong generating space.

Statement~(ii): immediate from~(i).

Statement~(iii): the conformal weight grading refines
the count, with each generator $v_i\in V$ contributing
$q^{s_i}$ to the weight-refined Poincar\'e polynomial.

Statement~(iv): filter $\Barchord$ by
the \emph{pole-order depth}: assign filtration degree~$p$
to a bar element whose differential involves OPE products
$a_{(n)}b$ with $n\le p$. The $E_0$ page uses only the
zeroth product (the Lie bracket), so
$E_1^{p,*}=H_*(\fg;\,\Bbbk)$ where $\fg$ is the Lie algebra
with bracket $[a,b]=a_{(0)}b$. The differentials
$d_r\colon E_r^{p,q}\to E_r^{p+r,q-r+1}$ are controlled by
the $r$-th pole OPE coefficients. For affine algebras
(pole order~$\le 2$), the spectral sequence has at most one
nontrivial differential and collapses at $E_2$. For the
Virasoro algebra (pole order~$4$), the central charge and
its descendants produce nonzero differentials through
$E_4$.
\end{proof}

\begin{computation}[Concentration data for standard families;
\ClaimStatusProvedHere]
\label{comp:concentration-data}
\index{bar complex!ordered!concentration data}
\[
\renewcommand{\arraystretch}{1.15}
\begin{array}{lcccl}
\cA & \dim V & \text{weights } s_i
 & P(t,q)
 & \text{$E_r$ collapse} \\
\hline
\cH_k & 1 & \{1\}
 & 1+tq & E_2 \\
\widehat{\mathfrak{sl}}_2 & 3 & \{1,1,1\}
 & 1+3tq & E_2 \\
\widehat{\mathfrak{sl}}_3 & 8 & \{1^8\}
 & 1+8tq & E_2 \\
V_k(\fg),\;\fg\text{ simple} & \dim\fg & \{1^{\dim\fg}\}
 & 1+(\dim\fg)\,tq & E_2 \\
\mathrm{Vir}_c & 1 & \{2\}
 & 1+tq^2 & E_4 \\
W_3 & 2 & \{2,3\}
 & 1+t(q^2+q^3) & E_4 \\
W_N & N{-}1 & \{2,3,\ldots,N\}
 & 1+t\sum_{j=2}^N q^j & E_4
\end{array}
\]
The affine family has all generators at weight~$1$
(currents), and the depth spectral sequence collapses at
$E_2$ because the only higher pole is the double pole
(the Killing form term $k\,\kappa^{ab}/(z{-}w)^2$), which
produces a single nontrivial differential $d_1$. The
$W$-algebra family has generators at weights
$2,3,\ldots,N$ (the Virasoro field and its higher-spin
partners), and the quartic pole in the Virasoro self-OPE
forces differentials through $E_4$.

Why does the spectral sequence terminate at $E_4$
rather than persisting indefinitely? The
weight-depth identity $w + d = n - 1$ and the
structural gap at depth $d = n$
(Theorem~\ref{thm:gap-migration}(i)--(ii))
force the differential $d_r$ for $r \ge 4$ to have
trivial source or trivial target at every bidegree.
The $\cW_N$ family has no weight-$1$ generator, so the
depth filtration runs out of room. For the Virasoro
case ($N = 2$, single generator at weight~$2$), the
quartic pole produces exactly three nontrivial
differentials $d_1, d_2, d_3$; the $E_4$ page is
concentrated in degree~$1$ and the bar cohomology is
$1$-dimensional, confirming the Poincar\'e polynomial
$P(t,q) = 1 + tq^2$.
\end{computation}

\begin{computation}[Dimension table for Heisenberg ordered bar
complex; \ClaimStatusProvedHere]
\label{comp:heisenberg-dim-table}
\index{Heisenberg algebra!ordered bar!dimension table}
The ordered bar complex $\Barchord(\cH_k)$ has
a single weight-$1$ generator $\alpha$. The space
$\Barchord_{n,w}$ consists of bar words
$[\alpha^{(m_1)}|\cdots|\alpha^{(m_n)}]$ with
$\sum m_i=w$ and each $m_i\ge 1$: thus
$\dim\Barchord_{n,w}=\binom{w-1}{n-1}$
(compositions of~$w$ into~$n$ positive parts).
\[
\renewcommand{\arraystretch}{1.05}
\begin{array}{r|rrrrrrrr}
n \backslash w & 1 & 2 & 3 & 4 & 5 & 6 & 7 & 8 \\
\hline
1 & 1 & 1 & 1 & 1 & 1 & 1 & 1 & 1 \\
2 & 0 & 1 & 2 & 3 & 4 & 5 & 6 & 7 \\
3 & 0 & 0 & 1 & 3 & 6 & 10 & 15 & 21 \\
4 & 0 & 0 & 0 & 1 & 4 & 10 & 20 & 35 \\
5 & 0 & 0 & 0 & 0 & 1 & 5 & 15 & 35 \\
6 & 0 & 0 & 0 & 0 & 0 & 1 & 6 & 21 \\
7 & 0 & 0 & 0 & 0 & 0 & 0 & 1 & 7 \\
8 & 0 & 0 & 0 & 0 & 0 & 0 & 0 & 1
\end{array}
\]
This is Pascal's triangle rotated: column sums are
$2^{w-1}$ (all compositions of~$w$), and the alternating
column sums give $\chi_w=\sum_n(-1)^n\binom{w-1}{n-1}
=(1-1)^{w-1}=\delta_{w,1}\cdot(-1)$, matching the
pentagonal theorem (only $w=1$ survives at the Euler
level, since the pentagonal exponents for the single-generator
Heisenberg start at $w=1$).

The closed-form generating function:
\begin{equation}\label{eq:heisenberg-bar-gf}
\sum_{n,w\ge 1}
\dim\Barchord_{n,w}\;t^n q^w
\;=\;
\frac{tq}{1-tq/(1-q)}
\;=\;
\frac{tq(1-q)}{1-q-tq}.
\end{equation}
\end{computation}

\begin{computation}[Generator-level dimensions for affine
families; \ClaimStatusProvedHere]
\label{comp:affine-generator-dims}
\index{bar complex!ordered!generator-level dimensions}
For $V_k(\fg)$ with $\dim\fg=d$, all generators have
weight~$1$, and the weight-$w$ space at tensor degree~$n$
counts the number of ways to distribute conformal
weight~$w$ among $n$ generators, each of weight~$\ge 1$:
\[
\dim\Barchord_{n,w}
\;=\;
d^n\cdot\binom{w-1}{n-1}
\qquad
(n\ge 1,\;w\ge n).
\]
The factor $d^n$ counts the $d^n$ ordered sequences of
generators; the binomial coefficient counts weight
distributions. The Euler characteristic at weight~$w$ is
\[
\chi_w
\;=\;
\sum_{n=1}^w (-1)^n d^n\binom{w-1}{n-1}
\;=\;
-d(1-d)^{w-1},
\]
which matches $[-1+\eta(q)^d]_w$ at leading weights
(exact agreement requires the full pentagonal expansion
of $\eta^d$).

\[
\renewcommand{\arraystretch}{1.05}
\begin{array}{l|rrrrrr}
(d,w) & w{=}1 & 2 & 3 & 4 & 5 & 6 \\
\hline
d{=}1\;\;(\cH_k) & 1 & 1 & 1 & 1 & 1 & 1 \\
d{=}3\;\;(\widehat{\mathfrak{sl}}_2) & 3 & 9 & 27 & 81 & 243 & 729 \\
d{=}8\;\;(\widehat{\mathfrak{sl}}_3) & 8 & 64 & 512 & 4096 & 32768 & 262144
\end{array}
\]
This table gives $\dim\Barchord_{1,w}=d$
(weight~$w$ descendants of the $d$ generators at each mode level)
for $w=1,\ldots,6$, i.e.\ the number of bar words of
length~$1$ at weight~$w$ is $d$ for all~$w$ (one descendant per
generator per mode). The full $(n,w)$ table is
$d^n\binom{w-1}{n-1}$.
\end{computation}

The generating functions describe the ordered bar complex at the level of dimensions. Returning to the structural content and verify the central claim (that the double bar of the Yangian recovers the current algebra) at the first nontrivial rank.

\begin{computation}[Open-colour double bar for
$\mathfrak{sl}_3$: explicit verification]%
\label{comp:double-bar-sl3}%
\index{double bar!sl3@$\mathfrak{sl}_3$!explicit verification}%
We verify Proposition~\ref{prop:open-colour-double-bar}
for $\fg=\mathfrak{sl}_3$ at degrees~$1$--$3$ by explicit
computation (all $43$ tests passing in the computational
engine
\texttt{compute/tests/test\_double\_bar\_sl3.py}).

\emph{Degree~$1$: generators.}\enspace
The Chevalley generators of $\mathfrak{sl}_3$ in
$\mathfrak{gl}_3$ matrix units are
\[
\begin{array}{c@{\;=\;}cc@{\;=\;}cc@{\;=\;}c}
e_1 & E_{12}, & e_2 & E_{23}, & e_{12} & E_{13}, \\
f_1 & E_{21}, & f_2 & E_{32}, & f_{12} & E_{31}, \\
h_1 & E_{11}-E_{22}, & h_2 & E_{22}-E_{33}. & &
\end{array}
\]
These $8=\dim(\mathfrak{sl}_3)$ generators span
${\bar B}^{\mathrm{ord}}_1$,
with $d=0$ on degree~$1$. At each mode level $r\ge 0$ they
correspond to $X^{(r)}_0,X^{(r)}_1,\ldots$ in
$Y_\hbar(\mathfrak{sl}_3)$, hence to the current modes
$X_{n}\in\mathfrak{sl}_3[t]$.

\emph{Degree~$2$: the Lie bracket, no central extension.}\enspace
The bar differential on degree-$2$ elements extracts the
\emph{collision residue}
$d_{\mathrm{bar}}(s^{-1}a\otimes s^{-1}b)=s^{-1}[a,b]$,
where $[a,b]$ is the $\mathfrak{sl}_3$ Lie bracket. The
Cartan matrix entries give
\[
d(s^{-1}h_i\otimes s^{-1}e_j)
\;=\;
A_{ij}\,s^{-1}e_j,
\qquad
A=\begin{pmatrix}2&-1\\-1&2\end{pmatrix}.
\]
The composite root brackets are
\[
d(s^{-1}e_1\otimes s^{-1}e_2)
=s^{-1}e_{12},
\qquad
d(s^{-1}e_{12}\otimes s^{-1}f_1)
=-s^{-1}e_2,
\qquad
d(s^{-1}e_{12}\otimes s^{-1}f_2)
=s^{-1}e_1,
\]
\[
d(s^{-1}f_{12}\otimes s^{-1}e_1)
=s^{-1}f_2,
\qquad
d(s^{-1}f_{12}\otimes s^{-1}e_2)
=-s^{-1}f_1,
\qquad
d(s^{-1}e_{12}\otimes s^{-1}f_{12})
=s^{-1}(h_1+h_2).
\]
All $64$ degree-$2$ differentials match the
$\mathfrak{sl}_3$ Lie bracket. \textbf{No central extension
appears:} the output is always an $\mathfrak{sl}_3$
generator (never a scalar Killing form term).

\emph{Degree~$3$: Jacobi identity = classical YBE.}\enspace
The degree-$3$ consistency
$d^2_{\mathrm{CE}}=0$ on the Chevalley--Eilenberg complex
is equivalent to the Jacobi identity
$[X,[Y,Z]]+[Y,[Z,X]]+[Z,[X,Y]]=0$.
All $8^3=512$ triples of $\mathfrak{sl}_3$ generators are
verified computationally. The Serre relations
\[
[e_1,[e_1,e_2]]=[e_1,e_{12}]=0,
\qquad
[e_2,[e_2,e_1]]=-[e_2,e_{12}]=0
\]
follow from root-space vanishing:
$2\alpha_1+\alpha_2$ and $\alpha_1+2\alpha_2$ are not
roots of $\mathfrak{sl}_3$. The triangle sectors
$T^L_{12}$, $T^L_{21}$ at degree~$3$ have coefficient
$\lambda=1/2=\int_0^1(1-t)\,dt$ (the $n=2$ beta integral).

\emph{Why the central extension is invisible.}\enspace
The Yangian spectral OPE
$[E_i(u),F_j(v)]=\delta_{ij}\,(u{-}v)^{-1}H_i(v)+\cdots$
has only \emph{simple} poles at $u=v$. The deformation
parameter $\hbar=1/(k{+}3)$ shifts some poles off the
collision diagonal ($E_i(u)E_i(v)$ has pole at $u{-}v=\hbar$,
not at $u=v$) but never produces a double pole. the $d\!\log$ kernel absorbs one pole power, so the collision
residue extracts \emph{only} the simple-pole part:
the Lie bracket. The central extension
$k\,n\,\delta_{n+m,0}\,\kappa^{ab}$ in the affine OPE
$J^a(z)\,J^b(w)\sim k\,\kappa^{ab}/(z{-}w)^2+\cdots$
arises from a \emph{double} pole, which the Yangian
spectral OPE does not have. Recovery of the central
term requires the closed-colour (Verdier/Ran) direction of
the full $\SCchtop$ Koszul duality.

\emph{Uniform argument for $\mathfrak{sl}_N$.}\enspace
The computation extends to all $N\ge 2$ by the same three
ingredients:
\begin{enumerate}[label=\textup{(\roman*)},nosep]
\item $Y_\hbar(\mathfrak{sl}_N)$ has only simple poles at
$u=v$ in the spectral OPE;
\item all root spaces of $A_{N-1}$ are one-dimensional
(root multiplicity $=1$);
\item the Jacobi identity holds for $\mathfrak{sl}_N$.
\end{enumerate}
Together these give
$\Barchord(Y_\hbar(\mathfrak{sl}_N))
\simeq U(\mathfrak{sl}_N[t])$ for all $N$.
More generally, the argument applies to any
finite-dimensional simple Lie algebra $\fg$ (all root
multiplicities $=1$ for types $A$--$G$). For
infinite-dimensional Kac--Moody algebras with root
multiplicities~$>1$, the argument requires modification:
the root-space one-dimensionality
(Theorem~\ref{thm:root-space-one-dim}) fails, and the
strictification mechanism needs new input.
\end{computation}

The preceding computation treats the Yangian as an \emph{output}: the open-colour Koszul dual of the affine algebra. Reversing the perspective and treat the Yangian as an \emph{input} (a genuinely $E_1$-chiral algebra) and compute its ordered bar complex directly. This is the double bar: two successive applications of open-colour Koszul duality. By Proposition~\ref{prop:open-colour-double-bar}, the double bar $\Barchord(Y_\hbar(\fg))$ recovers the current algebra $U(\fg[t])$, without central extension.

\begin{computation}[The $E_1$ ordered shadow of $Y_\hbar(\mathfrak{sl}_2)$:
the Yangian viewed as an $E_1$-chiral algebra;
\ClaimStatusProvedHere]%
\label{comp:yangian-e1-shadow}%
\index{Yangian!E1 ordered shadow@$E_1$ ordered shadow|textbf}%
\index{Yangian!sl2@$\mathfrak{sl}_2$!ordered bar complex}%
\index{double bar!Yangian!explicit computation}%
\index{RTT relations!from Yangian ordered bar}%
The Yangian $Y_\hbar(\mathfrak{sl}_2)$ is the open-colour
Koszul dual of $\widehat{\mathfrak{sl}}_{2,k}$
(Construction~\ref{constr:dg-shifted-yangian-from-bar}).
It is an \emph{associative} algebra with spectral parameter,
a genuinely $E_1$-chiral algebra in the sense of
\S\ref{sec:setup}. The computation of its ordered bar complex,
treating $Y_\hbar(\mathfrak{sl}_2)$ as the \emph{input}
chiral algebra (not the output of duality).

\medskip
\noindent\emph{(1) The ``OPE'' of the Yangian: spectral
commutation relations.}
The Yangian $Y_\hbar(\mathfrak{sl}_2)$ in the Drinfeld
presentation has generating series
$E(u)$, $F(u)$, $H(u)$ as
in~\eqref{eq:sl2-dual-generators}. The spectral
``OPE'' (the spectral commutation relations) is
the RTT relation
\begin{equation}\label{eq:yangian-rtt-input}
R(u{-}v)\,T_1(u)\,T_2(v)
\;=\;
T_2(v)\,T_1(u)\,R(u{-}v),
\end{equation}
where $R(u) = 1 + \hbar\,P/u$ is the Yang $R$-matrix
($P$ the permutation on $\mathbb{C}^2\otimes\mathbb{C}^2$,
$\hbar = 1/(k{+}2)$) and
$T(u) = \bigl(\begin{smallmatrix}
k_1(u) & F(u) \\ E(u) & k_2(u)
\end{smallmatrix}\bigr)$.
This is the \emph{input} algebraic structure. The
expanded FRT relations are:
\begin{align}
(u{-}v{-}\hbar)\,E(u)\,E(v)
&\;=\;
(u{-}v{+}\hbar)\,E(v)\,E(u),
\label{eq:yangian-EE}\\[3pt]
(u{-}v{+}\hbar)\,F(u)\,F(v)
&\;=\;
(u{-}v{-}\hbar)\,F(v)\,F(u),
\label{eq:yangian-FF}\\[3pt]
[E(u),\,F(v)]
&\;=\;
\frac{\hbar}{u{-}v}\,
\bigl(k_1(u)\,k_2(v) - k_1(v)\,k_2(u)\bigr),
\label{eq:yangian-EF}\\[3pt]
[H(u),\,E(v)]
&\;=\;
\frac{2\hbar}{u{-}v}\,E(v)
\;+\;\text{(regular)},
\label{eq:yangian-HE}\\[3pt]
[H(u),\,F(v)]
&\;=\;
\frac{-2\hbar}{u{-}v}\,F(v)
\;+\;\text{(regular)}.
\label{eq:yangian-HF}
\end{align}
The critical observation: all poles in the spectral
parameter $u{-}v$ are \emph{simple}. There is no
double pole. The central extension of the affine algebra
$\widehat{\mathfrak{sl}}_{2,k}$, which lives in the
double pole of the affine OPE
$J^a(z)J^b(w) \sim k\,\kappa^{ab}/(z{-}w)^2 + \cdots$,
has been absorbed into the deformation parameter
$\hbar = 1/(k{+}2)$. This is the mechanism by which
the open-colour Koszul duality converts the
\emph{level} (a central extension, a double-pole datum)
into the \emph{deformation parameter} (a simple-pole
datum).

\medskip
\noindent\emph{(2) The $m_2$ of the Yangian: the spectral
commutator.}
The binary product $m_2$ in the Yangian is
\emph{not} the OPE residue but the ordinary associative
product of the algebra. The spectral commutator
(the collision residue at $u = v$) extracts the Lie
bracket from the spectral OPE. For the matrix generators
$T_{ij}(u)$ of the RTT presentation, the spectral
commutator is:
\begin{equation}\label{eq:yangian-m2}
[T_{ij}(u),\,T_{kl}(v)]
\;=\;
\frac{\hbar}{u{-}v}\,
\bigl(T_{kj}(v)\,T_{il}(u) - T_{kj}(u)\,T_{il}(v)\bigr).
\end{equation}
The mode-expanded form: writing
$T_{ij}(u) = \delta_{ij} + \sum_{r\ge 0}
t_{ij}^{(r)}\,u^{-r-1}$, the leading relation is
\begin{equation}\label{eq:yangian-mode-comm}
[t_{ij}^{(r)},\,t_{kl}^{(s)}]
\;=\;
\hbar\sum_{p=0}^{\min(r,s)-1}
\bigl(t_{kj}^{(p)}\,t_{il}^{(r+s-1-p)}
- t_{kj}^{(r+s-1-p)}\,t_{il}^{(p)}\bigr).
\end{equation}
At leading mode level ($r = s = 0$):
$[t_{ij}^{(0)},\,t_{kl}^{(0)}] = 0$ (the sum is empty).
This is the \emph{abelian} leading-order relation,
consistent with the fact that the Yangian is a
deformation of $U(\mathfrak{sl}_2[t])$, and the
degree-zero piece $U(\mathfrak{sl}_2)$ sits inside
as a Hopf subalgebra whose generators
$e = t_{21}^{(0)}$, $f = t_{12}^{(0)}$,
$h = t_{11}^{(0)} - t_{22}^{(0)}$ satisfy the
$\mathfrak{sl}_2$ relations.
The first nontrivial commutator appears at $r = 0$,
$s = 1$ (or vice versa):
\[
[e,\,E^{(1)}]
= [t_{21}^{(0)},\,t_{21}^{(1)}]
= \hbar\bigl(t_{21}^{(0)}\,t_{21}^{(0)}
 - t_{21}^{(0)}\,t_{21}^{(0)}\bigr) = 0
\]
(same-root commutator vanishes). The cross-root
commutator:
\[
[e,\,F^{(1)}]
= [t_{21}^{(0)},\,t_{12}^{(1)}]
= \hbar\,(t_{11}^{(0)}\,t_{22}^{(0)} - t_{12}^{(0)}\,t_{21}^{(0)})
= \hbar\,\bigl(1 - \text{lower}\bigr).
\]

\medskip
\noindent\emph{(3) The collision residue of the Yangian.}
The collision residue of the Yangian spectral OPE
at $u = v$ is computed by taking
$\operatorname{Res}_{u=v}$ of the
RTT relation~\eqref{eq:yangian-rtt-input}. The
$R$-matrix $R(u{-}v) = 1 + \hbar\,P/(u{-}v)$ has
a simple pole at $u = v$ with residue $\hbar\,P$.
The RTT relation gives:
\begin{align}
\operatorname{Res}_{u=v}\bigl[
R(u{-}v)\,T_1(u)\,T_2(v)\bigr]
&\;=\;
\hbar\,P\,T_1(v)\,T_2(v),
\notag\\
\operatorname{Res}_{u=v}\bigl[
T_2(v)\,T_1(u)\,R(u{-}v)\bigr]
&\;=\;
\hbar\,T_2(v)\,T_1(v)\,P.
\notag
\end{align}
The collision residue is therefore
\begin{equation}\label{eq:yangian-collision-residue}
r^{\mathrm{coll}}_Y
\;=\;
\hbar\,\bigl(P\,T_1(v)\,T_2(v)
- T_2(v)\,T_1(v)\,P\bigr)
\;=\;
\hbar\,\bigl[P,\,T_1(v)\,T_2(v)\bigr].
\end{equation}
This is the \emph{permutation commutator} of the
monodromy matrix. The key structural point: the
collision residue is \emph{linear} in $\hbar$ (it comes
from a simple pole), and it involves only the
permutation~$P$, not the full $R$-matrix. This is
the Yangian analogue of the fact that the affine OPE
at level~$k$ has collision residue $r^{\mathrm{coll}}(z) =
k\,\Omega/z$ (a simple pole with the Casimir, scaled by
the level; at $k=0$ the affine current algebra is abelian
and the residue vanishes).

At the Lie-algebra level
(setting $T_1(v)\,T_2(v) \to 1\otimes 1$ at leading
order), the collision residue reduces to
$r^{\mathrm{coll}}_Y = \hbar\,[P,\,1] = 0$:
the leading-order collision is trivial. The first
nontrivial contribution comes at mode level~$1$:
\begin{equation}\label{eq:yangian-collision-mode1}
r^{\mathrm{coll}}_{Y,1}
\;=\;
\hbar\,\bigl[P,\,
\sum_{i,j}e_{ij}\otimes t_{ij}^{(0)}
+ t_{ij}^{(0)}\otimes e_{ij}
\bigr]
\;=\;
\hbar\,\sum_{i,j}
\bigl(e_{ji}\otimes t_{ij}^{(0)}
- t_{ij}^{(0)}\otimes e_{ji}\bigr),
\end{equation}
which is the $\mathfrak{sl}_2$ Lie bracket encoded
in the first-order deformation.

\medskip
\noindent\emph{(4) The ``Yangian of the Yangian'':
$\cA^!_{\mathrm{line}}(Y_\hbar(\mathfrak{sl}_2))$.}
This is the central structural question. The Yangian
$Y_\hbar(\mathfrak{sl}_2)$ is the open-colour Koszul
dual of $\widehat{\mathfrak{sl}}_{2,k}$:
\[
\widehat{\mathfrak{sl}}_{2,k}
\;\xrightarrow{\;\text{open-colour KD}\;}
Y_\hbar(\mathfrak{sl}_2).
\]
What is the open-colour Koszul dual of the Yangian
itself? By
Proposition~\ref{prop:open-colour-double-bar}, the
double bar
$\Barchord(Y_\hbar(\mathfrak{sl}_2))
\simeq U(\mathfrak{sl}_2[t])$.
The Koszul dual of the Yangian is therefore:
\begin{equation}\label{eq:yangian-koszul-dual}
\cA^!_{\mathrm{line}}\bigl(Y_\hbar(\mathfrak{sl}_2)\bigr)
\;=\;
H^*\!\bigl(\Barchord(Y_\hbar(\mathfrak{sl}_2))
\bigr)^\vee
\;=\;
U(\mathfrak{sl}_2[t])^\vee.
\end{equation}
This is \emph{not} $Y_\hbar(\mathfrak{sl}_2)$ itself.
The Yangian is not self-dual under open-colour Koszul
duality. The duality chain is:
\[
\underbrace{\widehat{\mathfrak{sl}}_{2,k}}_{
\text{affine, }E_\infty}
\;\xrightarrow{\;\text{KD}\;}
\underbrace{Y_\hbar(\mathfrak{sl}_2)}_{
\text{Yangian, }E_1}
\;\xrightarrow{\;\text{KD}\;}
\underbrace{U(\mathfrak{sl}_2[t])^\vee}_{
\text{current coalgebra dual, }E_1}
\;\xrightarrow{\;\text{KD}\;}
\;\cdots
\]
The three-step chain does not close: open-colour
Koszul duality is an \emph{involution} on the
$E_\infty$ side (for Kac--Moody:
$\widehat{\fg}_k \mapsto
\widehat{\fg}_{-k-2h^\vee} \mapsto
\widehat{\fg}_k$), but it is \emph{not} an involution
on the $E_1$ side. The reason:
\begin{itemize}[nosep]
\item The affine algebra $\widehat{\fg}_k$ has a
 double pole (the central extension) which the
 open-colour bar complex ``consumes'' to produce the
 deformation parameter $\hbar$.
\item The Yangian $Y_\hbar(\fg)$ has only simple poles
 (the RTT singularities), so its bar complex sees only
 the Lie bracket; the double-bar is
 $U(\fg[t])$, the \emph{undeformed} current algebra.
\item The deformation parameter $\hbar$ is invisible
 to the double bar: it controls the \emph{position}
 of shifted poles ($u{-}v = \hbar$), not the collision
 residue at $u = v$.
\end{itemize}
Therefore: the Yangian is not self-dual. Instead, the
self-duality lives at the level of the full
$\SCchtop$-algebra, where the two-colour Koszul duality
(Theorem~\ref{thm:duality-involution}) acts as an
involution on the \emph{pair}
$(\cA^!_{\mathrm{ch}},\,\cA^!_{\mathrm{line}})$.
The open colour alone does not close.

However, there is a weaker self-duality at $\hbar = 0$
(the classical limit): in the classical limit
$Y_0(\mathfrak{sl}_2) = U(\mathfrak{sl}_2[t])$, and
\[
\cA^!_{\mathrm{line}}(U(\mathfrak{sl}_2[t]))
\;=\;
U(\mathfrak{sl}_2[t])^\vee
\;\simeq\;
U(\mathfrak{sl}_2[t])
\]
by the self-duality of the PBW filtration on the
universal enveloping algebra. The classical current
algebra $U(\fg[t])$ \emph{is} self-dual under
open-colour Koszul duality, as it should be: it is a
quadratic algebra (generated by $\fg[t]$ with the Lie
bracket as the only relation), and quadratic Koszul
duality is an involution.

\medskip
\noindent\emph{(5) The double bar
$\Barchord(Y_\hbar(\mathfrak{sl}_2))$
at degrees $1$, $2$, $3$: explicit computation.}

\emph{Degree~$1$: generators.}\enspace
The Yangian $Y_\hbar(\mathfrak{sl}_2)$ in the Drinfeld
presentation has generators $E^{(r)}$, $F^{(r)}$,
$H^{(r)}$ for $r \ge 0$. The bar complex at tensor
degree~$1$ is spanned by the desuspended generators:
\[
\Barchord_1(Y_\hbar)
\;=\;
\operatorname{span}\bigl\{
s^{-1}E^{(r)},\;s^{-1}F^{(r)},\;s^{-1}H^{(r)}
\;\big|\;r\ge 0\bigr\}.
\]
These are the current modes
$e_r := e\,t^r$, $f_r := f\,t^r$, $h_r := h\,t^r$ of
$\mathfrak{sl}_2[t]$. (The bar complex at degree~$1$
has $d = 0$: there is no differential on single-slot
elements.) At each mode level~$r$, there are $3$
generators ($\dim\mathfrak{sl}_2 = 3$), so the
degree-$1$ space is $3$-dimensional at each mode level.

\emph{Degree~$2$: the Yangian bar differential extracts
the Lie bracket of $\mathfrak{sl}_2[t]$.}\enspace
The bar differential at degree~$2$ extracts the
collision residue of the Yangian spectral OPE. By
equations~\eqref{eq:yangian-HE}--\eqref{eq:yangian-HF},
the poles are all simple, so the $d\!\log$ kernel
absorbs one power and the collision residue is the
zeroth-order Lie bracket. Explicitly:
\begin{align*}
d(s^{-1}E^{(r)}\otimes s^{-1}F^{(s)})
&\;=\;
\sum_{p=0}^{\min(r,s)-1}
\hbar\,s^{-1}\bigl(
H^{(r+s-1-p)} - H^{(p)}\bigr)
\\
&\xrightarrow{\;\hbar\to 0\;}
s^{-1}\delta_{r+s,1}\,H^{(0)}
\;=\;s^{-1}[e_r,f_s],
\\[4pt]
d(s^{-1}H^{(r)}\otimes s^{-1}E^{(s)})
&\;=\;
2\hbar\,s^{-1}E^{(r+s)}
\;+\;\text{(higher)}
\\
&\xrightarrow{\;\hbar\to 0\;}
2\,s^{-1}E^{(r+s)}
\;=\;s^{-1}[h_r,e_s],
\\[4pt]
d(s^{-1}H^{(r)}\otimes s^{-1}F^{(s)})
&\;=\;
-2\hbar\,s^{-1}F^{(r+s)}
\;+\;\text{(higher)}
\\
&\xrightarrow{\;\hbar\to 0\;}
-2\,s^{-1}F^{(r+s)}
\;=\;s^{-1}[h_r,f_s].
\end{align*}
The collision residue at $u = v$ of the Yangian
$E$--$E$ relation~\eqref{eq:yangian-EE} is:
\[
d(s^{-1}E^{(r)}\otimes s^{-1}E^{(s)})
\;=\;0,
\]
because the $E(u)E(v)$ relation
$(u{-}v{-}\hbar)\,E(u)\,E(v)
= (u{-}v{+}\hbar)\,E(v)\,E(u)$
has the pole at $u{-}v = \hbar$ (shifted off the
collision diagonal $u = v$), so
$\operatorname{Res}_{u=v}[\cdots] = 0$. This
correctly reproduces $[e_r,e_s] = 0$ in
$\mathfrak{sl}_2[t]$. Similarly
$d(s^{-1}F^{(r)}\otimes s^{-1}F^{(s)}) = 0$.

The Cartan self-commutator:
$d(s^{-1}H^{(r)}\otimes s^{-1}H^{(s)}) = 0$
(the Cartan subalgebra is abelian).

\textbf{Summary at degree $2$:} the bar differential
on $\Barchord_2(Y_\hbar)$ reproduces exactly
the Lie bracket of the polynomial current algebra
$\mathfrak{sl}_2[t]$:
\[
d_{\mathrm{bar}}
\;=\;
[-,-]_{\mathfrak{sl}_2[t]}.
\]
The level $k$ (hence $\hbar = 1/(k{+}2)$) enters only
through the mode-mixing terms in the $[E,F]$
commutator, which contribute to higher-order
corrections. At leading order in modes (fixed
$r{+}s$), the bar differential is the
$\mathfrak{sl}_2[t]$ bracket \textbf{without central
extension}, confirming
Proposition~\ref{prop:open-colour-double-bar}
at degree~$2$.

\emph{Degree~$3$: Jacobi identity = classical YBE for
$\mathfrak{sl}_2[t]$.}\enspace
The degree-$3$ bar complex
$\Barchord_3(Y_\hbar)$ has elements
$s^{-1}X^{(r)}\otimes s^{-1}Y^{(s)}\otimes s^{-1}Z^{(t)}$
with $X,Y,Z \in \{E,F,H\}$. The identity $d^2 = 0$
at degree~$3$ is equivalent to the Jacobi identity for
$\mathfrak{sl}_2[t]$:
\[
[X_r,[Y_s,Z_t]] + [Y_s,[Z_t,X_r]] + [Z_t,[X_r,Y_s]]
\;=\;0.
\]
For the Serre-type triples:
\[
d^2(s^{-1}E^{(r)}\otimes s^{-1}E^{(s)}\otimes
s^{-1}F^{(t)})
\;=\;0
\]
because $d(s^{-1}E^{(r)}\otimes s^{-1}E^{(s)}) = 0$
($E$--$E$ collision vanishes) and
$d(s^{-1}E^{(s)}\otimes s^{-1}F^{(t)})
\propto s^{-1}H^{(s+t-\cdots)}$, followed by
$d(s^{-1}E^{(r)}\otimes s^{-1}H^{(\cdots)})
\propto s^{-1}E^{(\cdots)}$, which matches the
Jacobi identity
$[e_r,[e_s,f_t]] + [e_s,[f_t,e_r]] = 0$.
All $3^3 \cdot (\text{mode triples}) = 27 \cdot
\text{(modes)}$ checks reduce to the Jacobi identity
for $\mathfrak{sl}_2[t]$.

The \emph{classical Yang--Baxter equation} for the
Yangian $r$-matrix
$r_Y(z) = \hbar\,\Omega/z$ is:
\[
[r_{Y,12}(z_{12}),\,r_{Y,13}(z_{13})]
+ [r_{Y,12}(z_{12}),\,r_{Y,23}(z_{23})]
+ [r_{Y,13}(z_{13}),\,r_{Y,23}(z_{23})]
\;=\;0,
\]
which reduces (since $r_Y(z) = \hbar\,\Omega/z$ with
$\Omega$ the $\mathfrak{sl}_2$ Casimir) to the
classical YBE for $\Omega$, proved by the Jacobi
identity. This is the same classical $r$-matrix
as the one obtained from
$\widehat{\mathfrak{sl}}_{2,k}$
(Computation~\ref{comp:ordered-bar-sl2}), confirming
that the double bar sees only the undeformed classical
structure.

\textbf{The three-degree summary:}
\[
\renewcommand{\arraystretch}{1.15}
\begin{array}{c|l|l}
\text{Degree}
 & \Barchord_n(Y_\hbar(\mathfrak{sl}_2))
 & \text{Identification} \\
\hline
1 & \{s^{-1}X^{(r)} : X \in \{E,F,H\},\;r\ge 0\}
 & \mathfrak{sl}_2[t] \\
2 & d = [-,-]_{\mathfrak{sl}_2[t]},\;\text{no central ext.}
 & \text{Lie bracket of } \mathfrak{sl}_2[t] \\
3 & d^2 = 0 \iff \text{Jacobi for } \mathfrak{sl}_2[t]
 & \text{CYBE for } r_Y(z)=\hbar\Omega/z
\end{array}
\]
This confirms
$\Barchord(Y_\hbar(\mathfrak{sl}_2))
\simeq U(\mathfrak{sl}_2[t])$
(Proposition~\ref{prop:open-colour-double-bar})
at the first three degrees. The computation extends
to all degrees by the same mechanism: the Yangian
spectral OPE has only simple poles at $u = v$, so
the bar complex at each degree extracts only the Lie
bracket of the current algebra, and $d^2 = 0$
follows from the Jacobi identity.

\medskip
\noindent\emph{(6) Why the deformation parameter
$\hbar$ is invisible.}
The double bar sees only collision residues at the
\emph{collision diagonal} $u = v$. The Yangian
deformation parameter $\hbar$ controls the position
of \emph{shifted} poles (e.g.\ $u{-}v = \pm\hbar$
in the $E$--$E$ relation~\eqref{eq:yangian-EE}) but
does not affect the residue at $u = v$. The shifted
poles are invisible to the bar differential because
they live away from the collision locus: in the
ordered configuration space
$\mathrm{Conf}_n^{\mathrm{ord}}(\mathbb{C})$, the
bar differential integrates the $d\!\log$ form only
at the collision boundary, not at the shifted locus.
This is the open-colour analogue of the statement
(Proposition~\ref{prop:open-colour-double-bar}) that
the central extension is invisible: the level
$k$ (which produces $\hbar$) has been consumed in the
first Koszul duality step, and the second step cannot
recover it.

\medskip
\noindent\emph{(7) The $R$-matrix of the Yangian,
viewed as input.}
The Yangian is itself an $E_1$-chiral algebra, so
it has its own ordered bar complex with its own
$R$-matrix $R_Y(z)$. By
Construction~\ref{constr:r-matrix-monodromy},
$R_Y(z)$ is the monodromy of the flat connection
defined by the bar differential of the Yangian.
Since the bar differential of the Yangian extracts
only the $\mathfrak{sl}_2[t]$ Lie bracket (parts~(2)
and~(5) above), the $R$-matrix is:
\begin{equation}\label{eq:yangian-r-matrix-own}
R_Y(z)
\;=\;
1 + \Omega_{\mathfrak{sl}_2}/z + O(z^{-2}),
\end{equation}
where $\Omega_{\mathfrak{sl}_2}$ is the Casimir of
$\mathfrak{sl}_2$ (not the Yangian Casimir). This
is the \emph{same} leading-order $R$-matrix as for
$\widehat{\mathfrak{sl}}_{2,k}$
(Computation~\ref{comp:ordered-bar-sl2}). The
deformation parameter $\hbar$ of the Yangian does
not appear in its own $R$-matrix at leading order;
it enters only through the shifted-pole structure
that the bar differential does not see.

This is a precise illustration of the loss phenomenon:
applying open-colour Koszul duality twice to
$\widehat{\mathfrak{sl}}_{2,k}$ gives
$U(\mathfrak{sl}_2[t])$, not
$\widehat{\mathfrak{sl}}_{2,k}$. The central
extension (the datum that distinguishes the affine
algebra from the current algebra) has been lost in
the first duality step and cannot be recovered by
the second. Recovery requires the closed-colour
direction (Verdier/Ran duality on the full
chiral bar complex $\barBch(\cA)$).
\end{computation}

The preceding computation establishes $Y_\hbar(\mathfrak{sl}_2)$
as the prototypical $E_1$-chiral quantum group at rank~$1$.
The next example pushes the construction to rank~$2$,
where the root system has two simple roots, the $R$-matrix
acts on $\mathbb{C}^3\otimes\mathbb{C}^3$, and the Gauss
decomposition of the monodromy matrix has a richer structure.
The computation verifies that the $E_1$-chiral quantum group
framework of Definition~\ref{def:e1-chiral-quantum-group}
accommodates higher-rank Lie algebras without modification.

\begin{example}[Yangian $Y_\hbar(\mathfrak{sl}_3)$ as an
$E_1$-chiral quantum group;
\ClaimStatusProvedHere]%
\label{ex:yangian-sl3-e1-chiral-qg}%
\index{Yangian!sl3@$\mathfrak{sl}_3$!E1-chiral quantum group@$E_1$-chiral quantum group|textbf}%
\index{E1-chiral quantum group@$E_1$-chiral quantum group!sl3@$\mathfrak{sl}_3$ Yangian|textbf}%
\index{RTT relations!sl3@$\mathfrak{sl}_3$|textbf}%
\index{Gauss decomposition!sl3@$\mathfrak{sl}_3$ Yangian}%
\index{R-matrix!sl3@$\mathfrak{sl}_3$ Yang|textbf}%
The Yangian $Y_\hbar(\mathfrak{sl}_3)$ is the open-colour
Koszul dual of the affine algebra
$\widehat{\mathfrak{sl}}_{3,k}$
(Construction~\ref{constr:dg-shifted-yangian-from-bar}),
with $\hbar = 1/(k{+}3)$ (dual Coxeter number $h^\vee = 3$).
The $E_1$-chiral quantum group structure on
$Y_\hbar(\mathfrak{sl}_3)$ is verified in five steps,
following the pattern of
Computation~\ref{comp:yangian-e1-shadow} for $\mathfrak{sl}_2$
and the double-bar verification of
Computation~\ref{comp:double-bar-sl3}.

\medskip
\noindent\emph{(1) The RTT presentation: $9\times 9$ Yang
$R$-matrix.}
The fundamental $R$-matrix of $Y_\hbar(\mathfrak{sl}_3)$ acts
on
$\mathbb{C}^3\otimes\mathbb{C}^3$ and takes the Yang form
\begin{equation}\label{eq:sl3-yang-R}
R(u)
\;=\;
1 + \frac{\hbar}{u}\,P,
\end{equation}
where $P$ is the $9\times 9$ permutation matrix
$P(v_i\otimes v_j) = v_j\otimes v_i$ on
$\mathbb{C}^3\otimes\mathbb{C}^3$. In components:
$P^{kl}_{ij} = \delta^k_j\,\delta^l_i$.
The $R$-matrix satisfies the quantum Yang--Baxter equation
\begin{equation}\label{eq:sl3-qybe}
R_{12}(u{-}v)\,R_{13}(u{-}w)\,R_{23}(v{-}w)
\;=\;
R_{23}(v{-}w)\,R_{13}(u{-}w)\,R_{12}(u{-}v)
\end{equation}
in $\operatorname{End}((\mathbb{C}^3)^{\otimes 3})$,
verified by direct expansion using $P_{12}P_{13} = P_{23}P_{12}$
(the braid relation for transpositions in $S_3$).

The RTT relation is
\begin{equation}\label{eq:sl3-rtt}
R(u{-}v)\,T_1(u)\,T_2(v)
\;=\;
T_2(v)\,T_1(u)\,R(u{-}v),
\end{equation}
where $T(u) = (T_{ij}(u))_{1\le i,j\le 3}$ is a $3\times 3$
matrix of generating series. In components:
\begin{equation}\label{eq:sl3-rtt-components}
[T_{ij}(u),\,T_{kl}(v)]
\;=\;
\frac{\hbar}{u{-}v}\,
\bigl(T_{kj}(v)\,T_{il}(u) - T_{kj}(u)\,T_{il}(v)\bigr).
\end{equation}
This is the same structural equation as the $\mathfrak{sl}_2$
case~\eqref{eq:yangian-m2}, but now with indices running
over $\{1,2,3\}$: nine generating series $T_{ij}(u)$
instead of four.

\medskip
\noindent\emph{(2) Gauss decomposition and Drinfeld generators.}
The $3\times 3$ monodromy matrix admits a Gauss decomposition
\begin{equation}\label{eq:sl3-gauss}
T(u)
\;=\;
\begin{pmatrix}
1 & 0 & 0 \\
F_1(u) & 1 & 0 \\
F_{21}(u) & F_2(u) & 1
\end{pmatrix}
\begin{pmatrix}
K_1(u) & 0 & 0 \\
0 & K_2(u) & 0 \\
0 & 0 & K_3(u)
\end{pmatrix}
\begin{pmatrix}
1 & E_1(u) & E_{12}(u) \\
0 & 1 & E_2(u) \\
0 & 0 & 1
\end{pmatrix},
\end{equation}
where:
\begin{itemize}[nosep]
\item $E_i(u)$ ($i=1,2$) are the simple positive root
  currents, corresponding to simple roots $\alpha_1$ and
  $\alpha_2$ of $A_2$;
\item $E_{12}(u)$ is the composite root current for
  $\alpha_1+\alpha_2$, determined by
  \begin{equation}\label{eq:sl3-composite-E}
  E_{12}(u) = E_1(u)\,E_2(u) - E_2(u{-}\hbar)\,E_1(u);
  \end{equation}
\item $F_i(u)$ ($i=1,2$) are the simple negative root
  currents; $F_{21}(u)$ is the composite;
\item $K_i(u)$ ($i=1,2,3$) are the Cartan currents, with
  the constraint $K_1(u)\,K_2(u-\hbar)\,K_3(u-2\hbar) = 1$ (equivalently,
  the quantum determinant is central).
\end{itemize}
The independent Drinfeld generators are
$E_1(u)$, $E_2(u)$, $F_1(u)$, $F_2(u)$, $K_1(u)$,
$K_2(u)$ (six generating series), with $K_3(u)$
determined by the quantum determinant condition and
$E_{12}(u)$, $F_{21}(u)$ determined
by~\eqref{eq:sl3-composite-E} and its transpose.

Expanding $T_{ij}(u) = \delta_{ij} + \sum_{r\ge 0}
t_{ij}^{(r)}\,u^{-r-1}$, the Gauss decomposition
gives $8 = \dim(\mathfrak{sl}_3)$ independent generators
at each mode level (matching
Computation~\ref{comp:double-bar-sl3}).

\medskip
\noindent\emph{(3) Spectral commutation relations from the
Gauss decomposition.}
Substituting the Gauss decomposition~\eqref{eq:sl3-gauss}
into the RTT relation~\eqref{eq:sl3-rtt} and extracting
the $(i,j)$ components yields the Drinfeld
relations for $Y_\hbar(\mathfrak{sl}_3)$:
\begin{align}
(u{-}v{-}\hbar)\,E_i(u)\,E_i(v)
&\;=\;
(u{-}v{+}\hbar)\,E_i(v)\,E_i(u),
\qquad i = 1,2,
\label{eq:sl3-EiEi}\\[3pt]
(u{-}v{+}\hbar)\,F_i(u)\,F_i(v)
&\;=\;
(u{-}v{-}\hbar)\,F_i(v)\,F_i(u),
\qquad i = 1,2,
\label{eq:sl3-FiFi}\\[3pt]
[E_i(u),\,F_j(v)]
&\;=\;
\frac{\delta_{ij}\,\hbar}{u{-}v}\,
\bigl(K_i(u)\,K_{i+1}(v)^{-1}
- K_i(v)\,K_{i+1}(u)^{-1}\bigr),
\label{eq:sl3-EF}\\[3pt]
K_i(u)\,E_j(v)\,K_i(u)^{-1}
&\;=\;
\frac{u-v+\hbar\,A_{ij}/2}
{u-v-\hbar\,A_{ij}/2}\;E_j(v),
\label{eq:sl3-KE}\\[3pt]
K_i(u)\,F_j(v)\,K_i(u)^{-1}
&\;=\;
\frac{u-v-\hbar\,A_{ij}/2}
{u-v+\hbar\,A_{ij}/2}\;F_j(v),
\label{eq:sl3-KF}
\end{align}
where $A = \bigl(\begin{smallmatrix}2&-1\\-1&2\end{smallmatrix}\bigr)$ is the $A_2$ Cartan matrix.
The \emph{cross-root} relations ($i\neq j$, so $A_{ij}=-1$)
are:
\begin{align}
(u{-}v{+}\hbar)\,E_1(u)\,E_2(v)
&\;=\;
(u{-}v{-}\hbar)\,E_2(v)\,E_1(u)
\;+\;\text{(correction terms)},
\label{eq:sl3-E1E2}\\[3pt]
(u{-}v{+}\hbar)\,F_2(u)\,F_1(v)
&\;=\;
(u{-}v{-}\hbar)\,F_1(v)\,F_2(u)
\;+\;\text{(correction terms)}.
\label{eq:sl3-F2F1}
\end{align}
The Serre relations in the Yangian are the cubic spectral
identities:
\begin{equation}\label{eq:sl3-serre}
\operatorname{Sym}_{u_1,u_2}
\bigl[E_i(u_1),\,[E_i(u_2),\,E_j(v)]\bigr]
\;=\;0
\qquad(i\neq j),
\end{equation}
where $\operatorname{Sym}_{u_1,u_2}$ symmetrises in the
two spectral parameters. These are the $A_2$ quantum Serre
relations, the spectral-parameter lifts of
$[e_i,[e_i,e_j]] = 0$ verified in
Computation~\ref{comp:double-bar-sl3}.

The structural point, as for $\mathfrak{sl}_2$: all poles
in $u{-}v$ are simple. The double poles of the affine OPE
$J^a(z)\,J^b(w) \sim k\,\kappa^{ab}/(z{-}w)^2 + \cdots$
have been absorbed into $\hbar = 1/(k{+}3)$.

\medskip
\noindent\emph{(4) The spectral coproduct.}
The Yangian $Y_\hbar(\mathfrak{sl}_3)$ carries a Hopf algebra
structure with spectral coproduct
$\Delta : Y_\hbar(\mathfrak{sl}_3) \to
Y_\hbar(\mathfrak{sl}_3)^{\widehat\otimes\, 2}$ given in the
RTT presentation by
\begin{equation}\label{eq:sl3-coproduct-rtt}
\Delta(T_{ij}(u))
\;=\;
\sum_{a=1}^{3}\,T_{ia}(u)\otimes T_{aj}(u).
\end{equation}
This is the \emph{matrix coproduct}: the monodromy matrix
$T(u)$ is group-like in the matrix sense
($\Delta(T) = T\,\dot\otimes\,T$, where $\dot\otimes$ denotes
the matrix tensor product with entries multiplied in the
algebra tensor product). In terms of the Drinfeld generators
obtained from the Gauss decomposition~\eqref{eq:sl3-gauss},
the coproduct reads:
\begin{align}
\Delta(K_i(u))
&\;=\;
K_i(u)\otimes K_i(u)
\;+\;\text{(lower triangular corrections)},
\label{eq:sl3-coprod-K}\\[3pt]
\Delta(E_i(u))
&\;=\;
E_i(u)\otimes 1
\;+\;
K_i(u)\,K_{i+1}(u)^{-1}\otimes E_i(u)
\;+\;\text{(mixed terms)},
\label{eq:sl3-coprod-E}\\[3pt]
\Delta(F_i(u))
&\;=\;
1\otimes F_i(u)
\;+\;
F_i(u)\otimes K_i(u)^{-1}\,K_{i+1}(u)
\;+\;\text{(mixed terms)}.
\label{eq:sl3-coprod-F}
\end{align}
The ``mixed terms'' involve products of root currents
from distinct Gauss factors; they vanish at leading
mode level, where the coproduct reduces to the
standard $U(\mathfrak{sl}_3)$ coproduct
$\Delta(e_i) = e_i\otimes 1 + 1\otimes e_i$,
$\Delta(f_i) = f_i\otimes 1 + 1\otimes f_i$,
$\Delta(h_i) = h_i\otimes 1 + 1\otimes h_i$.

The spectral coproduct is coassociative:
$(\Delta\otimes\operatorname{id})\circ\Delta
= (\operatorname{id}\otimes\Delta)\circ\Delta$,
which in the RTT presentation follows from associativity of
matrix multiplication:
\[
\sum_{a,b}T_{ia}\otimes T_{ab}\otimes T_{bj}
\;=\;
\sum_{a,b}T_{ia}\otimes T_{ab}\otimes T_{bj}.
\]
The counit is $\varepsilon(T_{ij}(u)) = \delta_{ij}$,
and the antipode is $S(T(u)) = T(u)^{-1}$
(the matrix inverse, which exists as a formal power series
in $u^{-1}$).

This coproduct is the chiral coproduct
$\Delta: \cA\to\cA\widehat\otimes\cA$ of
Definition~\ref{def:e1-chiral-quantum-group}.
It lives on $Y_\hbar(\mathfrak{sl}_3)$ itself (not on the
bar complex $\Barchord(Y_\hbar)$, whose deconcatenation
coproduct is an independent structural datum).

\medskip
\noindent\emph{(5) The $r$-matrix and quasi-triangularity.}
The classical $r$-matrix of $Y_\hbar(\mathfrak{sl}_3)$ is
\begin{equation}\label{eq:sl3-classical-r}
r(z)
\;=\;
\frac{\hbar\,\Omega_{\mathfrak{sl}_3}}{z},
\end{equation}
where $\Omega_{\mathfrak{sl}_3}$ is the quadratic Casimir
of $\mathfrak{sl}_3$:
\begin{equation}\label{eq:sl3-casimir}
\Omega_{\mathfrak{sl}_3}
\;=\;
\sum_{a=1}^{8}I^a\otimes I_a
\;=\;
\sum_{i=1}^{2}
\bigl(e_i\otimes f_i + f_i\otimes e_i\bigr)
+ \bigl(e_{12}\otimes f_{12} + f_{12}\otimes e_{12}\bigr)
+ \sum_{i,j=1}^{2}(A^{-1})_{ij}\,h_i\otimes h_j,
\end{equation}
with $A^{-1} = \frac{1}{3}\bigl(\begin{smallmatrix}2&1\\1&2\end{smallmatrix}\bigr)$ the inverse Cartan matrix. The Casimir is normalised so that
$\kappa(\Omega) = \dim(\mathfrak{sl}_3) = 8$
(trace in the adjoint representation).

The classical $r$-matrix~\eqref{eq:sl3-classical-r}
satisfies the classical Yang--Baxter equation:
\begin{equation}\label{eq:sl3-cybe}
[r_{12}(z_{12}),\,r_{13}(z_{13})]
+ [r_{12}(z_{12}),\,r_{23}(z_{23})]
+ [r_{13}(z_{13}),\,r_{23}(z_{23})]
\;=\;0,
\end{equation}
which reduces to the $\mathfrak{sl}_3$ Jacobi identity
(the same mechanism as for $\mathfrak{sl}_2$ in
Computation~\ref{comp:yangian-e1-shadow}, part~(5)).

The quantum $R$-matrix is the Yang
matrix~\eqref{eq:sl3-yang-R},
$R(u) = 1 + \hbar\,P/u$. Quasi-triangularity at leading
order in $\hbar$ is the identity
\begin{equation}\label{eq:sl3-quasi-triang}
R(u{-}v)\,\Delta(T_{ij}(u))
\;=\;
\Delta^{\mathrm{op}}(T_{ij}(u))\,R(u{-}v)
\end{equation}
for all $i,j\in\{1,2,3\}$, where
$\Delta^{\mathrm{op}} = \sigma\circ\Delta$ is the opposite
coproduct ($\sigma$ the tensor flip). This identity is
equivalent to the RTT relation~\eqref{eq:sl3-rtt}:
substituting $\Delta(T_{ij}(u)) = \sum_a T_{ia}(u)\otimes T_{aj}(u)$
into the left-hand side of~\eqref{eq:sl3-quasi-triang} gives
\[
\sum_a R(u{-}v)\bigl(T_{ia}(u)\otimes T_{aj}(u)\bigr)
\;=\;
\sum_a \bigl(T_{aj}(u)\otimes T_{ia}(u)\bigr)R(u{-}v),
\]
which is precisely the RTT relation with the roles of the
two tensor factors encoding the coproduct. The
quasi-triangularity of the Yangian is therefore not an
additional axiom but a restatement of the defining RTT
relation in the Hopf-algebraic language.

\medskip
\noindent\emph{(6) Verification of
Definition~\ref{def:e1-chiral-quantum-group}.}
The four axioms of an $E_1$-chiral quantum group
(Definition~\ref{def:e1-chiral-quantum-group}) are
verified for $Y_\hbar(\mathfrak{sl}_3)$:
\begin{enumerate}[label=\textup{(\roman*)},nosep]
\item \textbf{Chiral coproduct.}\enspace
  $\Delta(T_{ij}(u)) = \sum_a T_{ia}(u)\otimes T_{aj}(u)$
  is coassociative (part~(4) above). The coproduct is
  a genuine Hopf-theoretic datum on
  $Y_\hbar(\mathfrak{sl}_3)$ itself, not the
  deconcatenation coproduct on $\Barchord(Y_\hbar)$.
\item \textbf{$R$-matrix.}\enspace
  $R(u) = 1 + \hbar\,P/u$ in
  $\operatorname{End}(\mathbb{C}^3\otimes\mathbb{C}^3)$
  satisfies the quantum YBE~\eqref{eq:sl3-qybe}.
\item \textbf{Quasi-triangularity.}\enspace
  $R\,\Delta = \Delta^{\mathrm{op}}\,R$
  (equation~\eqref{eq:sl3-quasi-triang}), equivalent to
  the RTT relation.
\item \textbf{Antipode.}\enspace
  $S(T(u)) = T(u)^{-1}$, satisfying
  $m\circ(S\otimes\operatorname{id})\circ\Delta = \eta\circ\varepsilon$
  by the matrix identity $T^{-1}\cdot T = I$.
\end{enumerate}
Therefore $Y_\hbar(\mathfrak{sl}_3)$ is an $E_1$-chiral
quantum group in the sense of
Definition~\ref{def:e1-chiral-quantum-group}. Its module
category $\operatorname{Mod}_{Y_\hbar(\mathfrak{sl}_3)}$ is
braided monoidal ($E_2$ in $\mathrm{Cat}$), with the braiding
given by the action of $R(u{-}v)$ on tensor products of
evaluation modules.

\medskip
\noindent\emph{(7) Comparison with $\mathfrak{sl}_2$ and the
rank-independence of the framework.}
The passage from $\mathfrak{sl}_2$ to $\mathfrak{sl}_3$
modifies only the representation-theoretic data:
\[
\renewcommand{\arraystretch}{1.15}
\begin{array}{c|c|c}
& Y_\hbar(\mathfrak{sl}_2)
& Y_\hbar(\mathfrak{sl}_3) \\
\hline
\text{Fundamental rep} & \mathbb{C}^2 & \mathbb{C}^3 \\
R\text{-matrix size} & 4\times 4 & 9\times 9 \\
\text{Simple roots} & 1\;(\alpha) & 2\;(\alpha_1,\alpha_2) \\
\text{Positive roots} & 1 & 3\;(\alpha_1,\alpha_2,\alpha_1{+}\alpha_2) \\
\text{Drinfeld generators} & E,F,H & E_1,E_2,F_1,F_2,K_1,K_2 \\
\text{Gauss factors} & 2\times 2 & 3\times 3 \\
\hbar & 1/(k{+}2) & 1/(k{+}3) \\
\text{Serre relations} & \text{none} & [E_i,[E_i,E_j]] = 0\;(i{\neq}j) \\
\text{Casimir }\Omega & 3\text{-term} & 8\text{-term}
\end{array}
\]
The $E_1$-chiral quantum group structure
(Definition~\ref{def:e1-chiral-quantum-group}) is
\emph{identical} in form: $R(u) = 1 + \hbar\,P/u$,
$\Delta(T) = T\,\dot\otimes\,T$,
$R\,\Delta = \Delta^{\mathrm{op}}\,R$,
$S(T) = T^{-1}$. The Yang $R$-matrix, the matrix
coproduct, quasi-triangularity, and the antipode
are rank-independent. The Lie-algebraic content enters
only through:
\begin{enumerate}[label=\textup{(\alph*)},nosep]
\item the size of the permutation matrix $P$
  ($N^2\times N^2$ for $\mathfrak{sl}_N$);
\item the Gauss decomposition pattern (lower-triangular
  $\times$ diagonal $\times$ upper-triangular, with composite
  root entries);
\item the Serre relations (trivial for $\mathfrak{sl}_2$,
  cubic for $\mathfrak{sl}_3$, with the degree growing as
  $1-A_{ij}$);
\item the dual Coxeter number $h^\vee$ in
  $\hbar = 1/(k{+}h^\vee)$.
\end{enumerate}
The framework extends uniformly to
$Y_\hbar(\mathfrak{sl}_N)$ for all $N\ge 2$, and more
generally to $Y_\hbar(\fg)$ for any finite-dimensional simple
$\fg$ (types $A$--$G$), using the RTT presentation with the
Yang $R$-matrix for the fundamental representation. The only
rank-dependent complexity is in the Gauss decomposition
(which has $|\Phi^+|$ off-diagonal entries in each triangular
factor) and the Serre relations (whose degree equals
$1-A_{ij}$).
\end{example}

\subsection{The annular bar complex and genus-one curvature}

The strip ordered bar complex
$\Barchord(\cA)$
on $\FM_k(\C)\times\Conf_k^<(\R)$ computes
genus-zero data. Replacing the strip~$\R$ by the
circle $S^1=\R/\Z$ wraps the topological direction and
produces a genus-one analogue: the
\emph{annular bar complex}.

\begin{definition}[Annular bar complex]%
\label{def:annular-bar}%
Let $\cA$ be a strongly admissible augmented $E_1$-chiral
algebra and $C=\Barchord(\cA)$ its ordered
bar coalgebra, with $C_\Delta$ the diagonal bicomodule.
The \emph{annular bar complex} is
\[
B^{\mathrm{ann}}_\bullet(\cA)
\;:=\;
C_\bullet\!\bigl(
\overline{\FM}{}^{\mathrm{ann}}_{r,s};\,
\omega_{\mathrm{bar}}
\bigr)
\;\cong\;
\bigoplus_{n\ge 0}
\mathrm{cyc}^n(C,C_\Delta),
\]
where $\mathrm{cyc}^n(C,C_\Delta):=
C_\Delta\widehat\otimes_C C^{\widehat\otimes\,n}$ is the
cyclic cotensor.
\end{definition}

\begin{theorem}[Annular bar differential;
\ClaimStatusProvedHere]%
\label{thm:annular-bar-differential}%
The differential on $B^{\mathrm{ann}}_\bullet(\cA)$
decomposes as
\[
d^{\mathrm{ann}}
\;=\;
d_{\mathrm{int}}+d_{\mathrm{res}}
 +d_\Delta+d_{\mathrm{wrap}},
\]
where:
\begin{enumerate}[label=\textup{(\alph*)}]
\item $d_{\mathrm{int}}$: internal differential on each
 tensor factor.
\item $d_{\mathrm{res}}$: collision residue differential
 (OPE residues at interior collisions).
\item $d_\Delta$: comultiplication differential at
 non-wrapping boundary collisions.
\item $d_{\mathrm{wrap}}$: the wrap-around differential
 at the seam:
\[
d_{\mathrm{wrap}}\langle a_1|\cdots|a_n\rangle
\;=\;
\sum_{(a_n)}
(-1)^{\epsilon}\,
\langle a_1|\cdots|a_{n-1}|a_n'|a_n''\rangle
\;+\;
\sum_{(a_1)}
(-1)^{\epsilon'}\,
\langle a_1'|a_1''|a_2|\cdots|a_n\rangle.
\]
\end{enumerate}
The identity $(d^{\mathrm{ann}})^2=0$ follows from
$\partial^2=0$ on the manifold with corners
$\overline{\FM}{}^{\mathrm{ann}}_{r,s}$.
\end{theorem}

The wrap-around differential $d_{\mathrm{wrap}}$ is the
only component without a strip analogue. Algebraically,
it provides the coface maps $\delta_0$ and $\delta_n$ that
close the two-sided cobar complex into the coHochschild
complex: without it the complex computes
$\mathrm{Cotor}^{C^e}(C_\Delta,C_\Delta)$; with it, the
result is the self-cotrace
$C_\Delta\square_{C^e}^{\mathbf{L}}
C_\Delta\simeq
\operatorname{coHH}_\bullet^{\mathrm{ch}}(C,C_\Delta)$.
The annular bar complex thereby computes
$\HH_\bullet^{\mathrm{ch}}(\cA)$ (chiral associative
Hochschild homology), connecting the ordered sector's
genus-one structure to the modular operad through
Theorem~\ref{thm:HH-coHH-homology}.

\begin{remark}[Genus-one $E_1$ vs $E_\infty$ surplus]%
\label{rem:genus-one-e1-einfty-surplus}%
\index{annular bar complex!E1 vs E-infty surplus@$E_1$ vs $E_\infty$ surplus}%
The ordered annular bar complex
$B^{\mathrm{ann}}_\bullet(\cA)$ and the
symmetric genus-one bar complex
$B^{g=1}_\bullet(\cA)$ stand in the same
ordered-to-unordered relation as their genus-zero
counterparts, but the surplus at genus one is
qualitatively richer.

At genus zero, the surplus is the pure braid group
representation (Computation~\ref{comp:e1-einfty-dimension-surplus}).
At genus one, the annular wrap-around introduces
the $A$-cycle monodromy as new datum. The ordered
annular bar complex $B^{\mathrm{ann}}_\bullet(\cA)$
carries both the $A$-cycle and $B$-cycle monodromies
independently; the symmetric genus-one complex
retains only their $\Sigma_n$-invariant combination,
which is the modular $S$-matrix.

For affine Kac--Moody algebras: the ordered annular
bar cohomology produces the \emph{elliptic quantum group}
$E_{q,p}(\fg)$ of Felder--Varchenko type, with the
$A$-cycle monodromy parameter~$p$ and the $B$-cycle
spectral parameter~$q$. The symmetric genus-one bar
cohomology collapses to the Verlinde ring
$K^0(\mathrm{Rep}_k(\widehat{\fg}))$ with fusion
numbers controlled by the $S$-matrix. The surplus at
genus one is thus the elliptic deformation parameter~$p$:
the datum that promotes a modular tensor category
(the Verlinde output) to an elliptic quantum group
(the ordered annular output).

For Virasoro: the ordered annular bar sees the full
genus-one curvature $\kappa(\mathrm{Vir}_c)\cdot\omega_1$
(where $\omega_1$ is the Hodge class on $\cM_{1,1}$)
through the wrap-around differential $d_{\mathrm{wrap}}$.
The interplay between the $A$-cycle monodromy and the
quartic-pole OPE (class~$\mathbf{M}$) produces
genuinely entangled genus-one data: the two colours of
$\SCchtop$ cannot be separated at genus one when the
Lie bracket coefficient $c_0 = \{L_{(0)}L\} = 2L$ is
nonzero
(Theorem~\ref{thm:elliptic-spectral-dichotomy}).
The Heisenberg ($c_0 = 0$) is the unique standard
family where the genus-one ordered and symmetric
theories decouple.
\end{remark}

\begin{remark}[Gauge-gravity complexity dichotomy]%
\label{rem:gauge-gravity-yangian-dichotomy}%
\index{gauge-gravity dichotomy!Yangian perspective|textbf}%
The genus-one entanglement has a sharper avatar in the $\Ainf$ structure. The open-colour Koszul dual distinguishes gauge and gravitational theories by two independent invariants:
\[
\begin{array}{c|cc}
& \Ainf\text{ products } m_k & \text{Coproduct } \Delta_z \\
\hline
\text{Gauge (CS)} & m_k = 0 \;(k \ge 3) &
 \text{nontrivial Yangian} \\
\text{Gravity (Vir)} & m_k \neq 0 \;\forall\, k &
 \text{strictly primitive}
\end{array}
\]
For gauge theories (Chern--Simons type, class~$\mathbf{L}$), the open-colour dual is a genuine Yangian $Y_\hbar(\fg)$ with a nontrivial spectral coproduct $\Delta_z$ but no higher $\Ainf$ products: the bar complex is formal and the Koszul dual is a quadratic algebra. For gravitational theories (Virasoro type, class~$\mathbf{M}$), the two columns swap: the higher products $m_k$ are all nonzero (an infinite $\Ainf$ tower), but the coproduct $\Delta_z$ degenerates to $\Delta_z(x) = x \otimes 1 + 1 \otimes x$ because the Virasoro algebra admits no finite-type Hopf structure.

Drinfeld--Sokolov reduction transports class~$\mathbf{L}$ to class~$\mathbf{M}$: it preserves Koszulness but destroys formality. The resolvent tree formula
$m_k^{\cW} = \sum_{\mathrm{Catalan}}(m_2^{\mathrm{aff}}
\circ h_{\mathrm{DS}})$
generates the infinite tower of higher products from the
single quadratic product of the affine input, with each tree
indexed by a Catalan binary bracketing of the DS homotopy transfer.
See Remark~\ref{rem:gravitational-yangian-explicit-brackets}
for explicit bracket computations.
\end{remark}

\begin{definition}[$E_1$-chiral quantum group]%
\label{def:e1-chiral-quantum-group}%
\index{E1-chiral quantum group@$E_1$-chiral quantum group|textbf}%
\index{quantum group!E1-chiral@$E_1$-chiral|textbf}%
An \textbf{$E_1$-chiral quantum group} on a curve $X$ is an
$E_1$-chiral algebra $\cA$ on $X$ equipped with:
\begin{enumerate}
\item a chiral coproduct
  $\Delta : \cA \to \cA \widehat\otimes \cA$
  (coassociative, in the appropriate completed tensor product
  of chiral algebras),
\item an $R$-matrix $R(z) \in \cA^{\widehat\otimes\, 2}$,
\item quasi-triangularity:
  $R\,\Delta(a) = \Delta^{\mathrm{op}}(a)\,R$
  for all $a \in \cA$,
\item an antipode $S: \cA \to \cA$ satisfying the Hopf
  compatibility,
\end{enumerate}
such that $\mathrm{Mod}_\cA$ is a braided monoidal category.
\end{definition}

\begin{remark}[Coproduct, $R$-matrix, and $E_1$ provenance]%
\label{rem:e1-chiral-qg-provenance}%
\index{coproduct!chiral vs deconcatenation}%
\index{R-matrix!classical shadow vs quantum lift}%
Several clarifications are essential to prevent
conflation of algebraically distinct structures.

The chiral coproduct $\Delta$ in
Definition~\ref{def:e1-chiral-quantum-group} lives on $\cA$
itself: it is an independent Hopf-theoretic datum, not to be
confused with the deconcatenation coproduct on the bar
complex $\Barchord(\cA)$, which is structural (encoding
interval-cutting in $\Conf_k^<(\R)$).

The $r$-matrix obtained from collision residues of the OPE
(Construction~\ref{constr:r-matrix-monodromy}) is the
classical shadow: it satisfies the CYBE and governs the
leading term $R(z) = 1 + r(z) + O(\hbar^2)$. The full
$R(z)$ is the quantum lift, satisfying the quantum
Yang--Baxter equation.

Quantum groups are $E_1$ objects; their module categories
carry the $E_2$ structure in $\mathrm{Cat}$. This is the
standard Lurie--Francis delooping: an $E_n$-algebra in
$\mathrm{Cat}$ is an $(n{-}1)$-braided monoidal category,
so $E_1$-quantum groups produce $E_2$ (braided monoidal)
module categories.

For $E_\infty$-chiral algebras (vertex algebras in the sense
of \S\ref{sec:HT-operad}), both the coproduct and $R$-matrix
are canonically determined by the local OPE data: the
coproduct is the vertex coalgebra structure of~\cite{JKL26},
and $R(z)$ is derived from monodromy of the OPE connection.
For genuinely $E_1$-chiral algebras, both $\Delta$ and $R(z)$
are independent structure not recoverable from a local OPE
; the discriminant is provenance, not value.
See~\cite{GRZ23} for the prototype in the context of deformed
double current algebras.
\end{remark}

\begin{lemma}[Coproduct--bracket compatibility; \ClaimStatusProvedHere]%
\label{lem:coproduct-bracket-compat}%
\index{coproduct!bracket compatibility|textbf}%
\index{E1-chiral quantum group@$E_1$-chiral quantum group!coproduct--bracket compatibility}%
Let $(\cA, \Delta, R(z), S)$ be an $E_1$-chiral quantum group
(Definition~\textup{\ref{def:e1-chiral-quantum-group}}). Then
$\Delta$ preserves all chiral operations: for every $a, b \in \cA$
and every $n \ge 0$,
\begin{equation}\label{eq:coproduct-bracket-compat}
\Delta\bigl(a_{(n)}b\bigr)
\;=\;
\Delta(a)_{(n)}\,\Delta(b),
\end{equation}
where the right-hand side uses the $n$-th product on
$\cA \widehat\otimes \cA$ induced by the tensor product chiral
algebra structure.
\end{lemma}

\begin{proof}
By Definition~\ref{def:e1-chiral-quantum-group},
$\Delta \colon \cA \to \cA \widehat\otimes \cA$ is a chiral
coproduct, meaning a morphism of chiral algebras (the target
carries the tensor product chiral algebra structure). A morphism
of chiral algebras preserves all $n$-th products by definition:
$\Delta(a_{(n)}b) = \Delta(a)_{(n)}\,\Delta(b)$. This is the
content of~\eqref{eq:coproduct-bracket-compat}.

For $E_\infty$-chiral algebras (vertex algebras), the $n$-th
product on $\cA \widehat\otimes \cA$ is determined by the
Leibniz rule
\[
(a_1 \otimes a_2)_{(n)}\,(b_1 \otimes b_2)
\;=\;
\sum_{j \in \Z}
(-1)^{|a_2||b_1|}
\bigl(a_1{}_{(j)}\,b_1\bigr) \otimes
\bigl(a_2{}_{(n-j)}\,b_2\bigr),
\]
and~\eqref{eq:coproduct-bracket-compat} expands into a
system of identities among the structure constants of
the OPE. For $V_k(\fg)$ with $\fg$ simple and
$\Delta(J^a(z)) = J^a(z) \otimes 1 + 1 \otimes J^a(z)$
(primitive coproduct on generators), the identity reduces to
\[
\Delta\bigl(J^a{}_{(n)}\,J^b\bigr)
\;=\;
\begin{cases}
f^{ab}{}_c\bigl(J^c\otimes 1 + 1\otimes J^c\bigr)
& n = 0,\\[2pt]
k\,\kappa^{ab}\bigl(1\otimes 1\bigr)
& n = 1,\\[2pt]
0 & n \ge 2,
\end{cases}
\]
which is immediate from $\Delta$ applied termwise to the
Kac--Moody OPE $J^a{}_{(0)}\,J^b = f^{ab}{}_c\,J^c$,
$J^a{}_{(1)}\,J^b = k\,\kappa^{ab}$, and the primitivity
of $\Delta$ on generators.
\end{proof}

\begin{remark}[Extensions: reductive and product algebras]%
\label{rem:coproduct-bracket-compat-extensions}%
\index{coproduct!bracket compatibility!reductive extension}%
The verification above extends without modification:
\begin{enumerate}[label=\textup{(\roman*)},nosep]
\item For $\fg$ reductive (e.g.\ $\mathfrak{gl}_n$):
  $H^3(\mathfrak{gl}_n, \C) = \C$, so the Kac--Moody
  vertex algebra $V_k(\mathfrak{gl}_n)$ carries a unique
  (up to scalar) invariant inner product and the primitive
  coproduct on generators satisfies
  \eqref{eq:coproduct-bracket-compat} by the same termwise
  argument.
\item For products of simple factors
  $\fg = \fg_1 \times \cdots \times \fg_r$: the coproduct
  and bracket decompose along the direct sum, and
  compatibility holds factor by factor.
\item For genuinely $E_1$-chiral algebras (e.g.\ Yangians),
  the coproduct is not primitive on generators
  (cf.\ Computation~\textup{\ref{comp:yangian-e1-shadow}});
  the compatibility~\eqref{eq:coproduct-bracket-compat}
  is then a nontrivial constraint satisfied by the RTT
  relations. In the Yangian case, this is the statement that
  the matrix coproduct
  $\Delta(T_{ij}(u)) = \sum_a T_{ia}(u) \otimes T_{aj}(u)$
  is an algebra homomorphism.
\end{enumerate}
\end{remark}

\begin{theorem}[Categorical level shift for chiral quantum groups;
\ClaimStatusProvedHere]%
\label{thm:categorical-level-shift}%
\index{categorical level shift|textbf}%
\index{E1-chiral quantum group@$E_1$-chiral quantum group!categorical level shift}%
\index{E2-category@$E_2$-category!from $E_1$-quantum group}%
Let $(\cA, \Delta, R(z), S)$ be an $E_1$-chiral quantum group on
a curve~$X$
\textup{(}Definition~\textup{\ref{def:e1-chiral-quantum-group}}\textup{)}.
Assume $\Mod_\cA$ is compactly generated as a stable
$\infty$-category.  Then:
\begin{enumerate}[label=\textup{(\roman*)}]
\item\label{item:cls-E1}
$\cA$ is an $E_1$-algebra: associative and non-commutative.
\item\label{item:cls-E2-Mod}
The category $\mathrm{Mod}_\cA$ of $\cA$-modules is
$E_2$ in $\mathsf{Cat}$: it is a braided monoidal category
with monoidal structure induced by $\Delta$ and braiding
induced by $R(z)$.
\item\label{item:cls-E2-center}
The chiral derived center
$Z^{\mathrm{der}}_{\mathrm{ch}}(\cA)
= C^\bullet_{\mathrm{ch}}(\cA, \cA)$
carries a natural $E_2$-algebra structure.
\item\label{item:cls-identification}
There is a canonical identification
\begin{equation}\label{eq:center-end-identification}
\mathrm{End}_{\mathrm{Fun}(\mathrm{Mod}_\cA,\,
\mathrm{Mod}_\cA)}(\id_{\mathrm{Mod}_\cA})
\;\simeq\;
Z^{\mathrm{der}}_{\mathrm{ch}}(\cA),
\end{equation}
and this identification intertwines the $E_2$-structure on
$\mathrm{Mod}_\cA$ from~\ref{item:cls-E2-Mod} with the
$E_2$-structure on
$Z^{\mathrm{der}}_{\mathrm{ch}}(\cA)$
from~\ref{item:cls-E2-center}.
\end{enumerate}
In particular, the $E_2$ structure lives on
$\mathrm{Mod}_\cA$ and on $Z^{\mathrm{der}}_{\mathrm{ch}}(\cA)$,
never on $\cA$ itself.
\end{theorem}

\begin{proof}
\ref{item:cls-E1}.
By definition, an $E_1$-chiral quantum group is an
$E_1$-chiral algebra
(Definition~\ref{def:e1-chiral-quantum-group}), hence
an $E_1$-algebra.  The non-commutativity is intrinsic: for
genuinely $E_1$-chiral algebras, the multiplication
$\mu \colon \cA \otimes \cA \to \cA$ does not factor
through the symmetric product, and the $R$-matrix
$R(z) \neq \tau$ is independent structure.

\medskip
\noindent\ref{item:cls-E2-Mod}.
The chiral coproduct $\Delta \colon \cA \to
\cA \widehat\otimes \cA$ makes $\cA$ a bialgebra, and
the standard construction equips $\mathrm{Mod}_\cA$ with
a monoidal structure: given $\cA$-modules $M$ and $N$,
their tensor product is
$M \otimes_\Delta N := (M \otimes N)$ with $\cA$-action via
$\Delta$.  Coassociativity of $\Delta$ gives associativity
of this tensor product (up to coherent homotopy).  The
$R$-matrix $R(z) \in \cA^{\widehat\otimes\,2}$ provides a
natural isomorphism
$\beta \colon M \otimes_\Delta N \xrightarrow{\sim}
N \otimes_\Delta M$
via $\beta(m \otimes n) = R(z) \cdot (n \otimes m)$.
Quasi-triangularity
$R\,\Delta(a) = \Delta^{\mathrm{op}}(a)\,R$ ensures that
$\beta$ is a morphism of $\cA$-modules.  The Yang--Baxter
equation for $R(z)$
(Theorem~\ref{thm:YBE}) is equivalent to the hexagon
axioms for the braiding on the triple tensor product,
so $(\mathrm{Mod}_\cA, \otimes_\Delta, \beta)$ is a
braided monoidal category.  In the language of the
$E_n$-delooping hypothesis
(Lurie~\cite{HA}, \S5.1.2): an $E_n$-algebra in
$\mathsf{Cat}$ is an $(n{-}1)$-braided monoidal category;
a braided monoidal category is therefore an $E_2$-algebra
in~$\mathsf{Cat}$.

\medskip
\noindent\ref{item:cls-E2-center}.
This is the Deligne conjecture applied to
$E_1$-algebras, proved by
Kontsevich~\cite{Kon03} and
Tamarkin~\cite{Tamarkin00}:
for any $E_1$-algebra~$\cA$, the Hochschild cochain complex
$C^\bullet(\cA, \cA)$ carries a natural $E_2$-algebra
structure.  The cup product provides the $E_1$ factor;
the Gerstenhaber bracket $[-,-]_G$ provides the additional
$E_1$ factor; together they assemble into $E_2$
(Theorem~\ref{thm:tamarkin-higher-structure}).
In the chiral setting, the relevant complex is the
chiral Hochschild cochain complex
$C^\bullet_{\mathrm{ch}}(\cA, \cA)$, computed via the
chiral endomorphism operad $\mathrm{End}^{\mathrm{ch}}_\cA$
with spectral parameters from $\FM_k(\C)$.
The same $E_2$-structure persists: the factorisation
algebra on $\FM_k(\C)$ provides the little-discs operations
that upgrade the cup product and bracket to a full
$E_2$-algebra
(Proposition~\ref{prop:derived-center-E2}).

\medskip
\noindent\ref{item:cls-identification}.
The derived center of $\cA$, defined as
$Z^{\mathrm{der}}_{\mathrm{ch}}(\cA)
= \RHom_{\mathrm{Fun}(\mathrm{Mod}_\cA,\,
\mathrm{Mod}_\cA)}(\id, \id)$,
is tautologically the endomorphism algebra of the identity
functor on $\mathrm{Mod}_\cA$.  The identification with the
chiral Hochschild cochain complex
$C^\bullet_{\mathrm{ch}}(\cA, \cA)$ is
Proposition~\ref{prop:derived-center-E2}.
It remains to verify that the two $E_2$-structures
coincide.  On one side, the $E_2$-structure on
$\mathrm{End}(\id_{\mathrm{Mod}_\cA})$ is induced by
the braided monoidal structure of~\ref{item:cls-E2-Mod}:
the braiding on $\mathrm{Mod}_\cA$ provides natural
transformations $\id \to \id$ that generate the $E_2$
operations on the endomorphism algebra.  On the other side,
the $E_2$-structure on
$C^\bullet_{\mathrm{ch}}(\cA, \cA)$
comes from the Deligne conjecture
(\ref{item:cls-E2-center}).  These agree because both are
the universal $E_2$-structure on the center of an
$E_1$-algebra in $\mathsf{Cat}$: the center functor
$Z \colon \mathrm{Alg}_{E_1}(\mathsf{Cat}) \to
\mathrm{Alg}_{E_2}(\mathsf{Cat})$
is the left adjoint to the forgetful functor (cf.\
Ben-Zvi--Francis--Nadler~\cite{BZFN10}, \S4), and
universality forces the two constructions to yield the
same $E_2$-algebra.
\end{proof}

\begin{remark}[$E_2$ lives on modules and center, not on $\cA$]%
\label{rem:E2-never-on-A}%
\index{E2-category@$E_2$-category!never on the algebra}%
Theorem~\ref{thm:categorical-level-shift} makes precise
the principle that quantum groups are $E_1$ objects whose
$E_2$ structure is visible only after categorical
delooping.  The algebra $\cA$ itself is non-commutative
and carries no braiding; the $R$-matrix acts on pairs of
modules, not on elements of~$\cA$.  Conflating the
$E_1$-level (the algebra) with the $E_2$-level (its
modules or its center) is a persistent source of error;
see Remark~\textup{\ref{rem:E2-category-E1-algebra}}
and Remark~\textup{\ref{rem:r-matrix-coproduct-warning}}
for the specific mechanisms.
\end{remark}

\section{Conventions: notions, Yangian types, level prefix}
\label{sec:ordered-kd-core-conventions}

\subsection{Five $E_1$-chiral notions, ordered KD reading}
\label{subsec:ordered-kd-core-five-notions}

\begin{convention}[Ordered KD core uses Notion~B + Notion~E hybrid;
\ClaimStatusProvedElsewhere]
\label{conv:ordered-kd-core-five-notions}
The ordered Koszul-duality theory of this chapter — the bar coalgebra
$\Barchord(\cA)$, the diagonal bicomodule $C_\Delta$, and the
duality involution (Theorem~\ref{thm:duality-involution}) — uses
two of the five $E_1$-chiral notions catalogued in Volume~I:
\begin{enumerate}[label=\textup{(\Alph*)}]
\setcounter{enumi}{1}
\item \textbf{Notion~B} ($A_\infty$ in $\mathrm{End}^{\mathrm{ch}}_\cA$):
the algebra side, where $\cA$ has $\{m_k\}_{k\ge 1}$ chain-level
chiral structure maps. The ordered bar
$\Barchord(\cA) = T^c(s^{-1}\bar\cA)$ takes Notion-B input.
\setcounter{enumi}{4}
\item \textbf{Notion~E} (factorization on
$\mathrm{Ran}^{\mathrm{ord}}(X)$): the coalgebra side, where the
deconcatenation coproduct on $\Barchord(\cA)$ presents the
ordered factorization-coalgebra structure. The diagonal
bicomodule $C_\Delta$ is the standard module of this Notion-E
object.
\end{enumerate}
The duality functor $\Barchord \dashv \Cobar$ converts Notion-B
input to Notion-E output and back; this is the $E_1$-chiral analogue
of the ordered associative Loday--Vallette theory. Notion~D
($A_\infty$ in $E_1$-chiral) appears only when the ordered
diagonal is itself an $A_\infty$-module, which happens for
gravitational Yangian targets at $A = \mathrm{Vir}_c$.
\end{convention}

\subsection{Four Yangian types and the chiral Yangian}
\label{subsec:ordered-kd-core-yangian-types}

\begin{convention}[Ordered chiral Koszul duality is the universal
source of the chiral Yangian $Y(\fg)^{\mathrm{ch}}$;
\ClaimStatusProvedElsewhere]
\label{conv:ordered-kd-core-yangian-types}
Four Yangian types appear across the programme: classical
$Y_\hbar(\fg)$ (E$_1$-top on $\R$); dg-shifted
$Y^{\mathrm{dg}}_\hbar(\fg)$ (point/formal disk); chiral
$Y(\fg)^{\mathrm{ch}}$ ($E_1$-chiral on a curve $X$); spectral
factorization on $\mathbb{A}^1_u$. The ordered chiral Koszul
duality of this chapter produces all four as different functorial
specialisations:
\begin{enumerate}[label=\textup{(\arabic*)}]
\item Restriction to a point $x\in X$: $Y(\fg)^{\mathrm{ch}}|_x$
recovers the classical $Y_\hbar(\fg)$ as the mode algebra; the
filtration on $\Cobar$ collapses to the standard PBW filtration
on $Y_\hbar(\fg)$.
\item Restriction to a formal disk $D\subset X$: the completed
ordered cobar $\widehat{\Cobar}|_D$ produces the dg-shifted
Yangian $Y^{\mathrm{dg}}_\hbar(\fg)$ with explicit $A_\infty$
correction operations $m_k$ ($k\ge 3$).
\item Globally on $X$: $Y(\fg)^{\mathrm{ch}} :=
\Cobar(\Barchord(U^{\mathrm{ch}}(\fg[t])))$ is the chiral
Yangian as $D$-module on $X$; this is the central object of the
chapter.
\item Spectral factorization: passing to the spectral line
$\mathbb{A}^1_u$ via the Koszul-locus evaluation modules gives the
Latyntsev factorization quantum group; the dictionary is
formal-disk--to--spectral-line transport along the universal
family.
\end{enumerate}
\end{convention}

\begin{remark}[Two definitions of chiral Yangian; equivalence open at
chain level]
\label{rem:ordered-kd-core-two-definitions}
Two definitions of chiral Yangian circulate. The weaker one
(def:e1-chiral-yangian) requires only an $E_1$-factorization
algebra on $X$ with RTT-presentation; the stronger one
(def:chiral-yangian-datum) requires the modular Maurer--Cartan
tower $(A, S(z), \Delta^{\mathrm{ch}}, \{m_k\})$ over the genus
hierarchy. The ordered chiral Koszul-duality construction of this
chapter manifestly produces the weaker version; the lift to the
stronger version is conditional on the Vol~II curved-Dunn $H^2=0$
result for $g\ge 2$ (closed for the non-logarithmic standard
landscape via the irregular-singular KZB framework, conjectural for
logarithmic $W(p)$ pending the Adamovi\'c--Milas character bound).
The duality involution (Theorem~\ref{thm:duality-involution}) is
associator-dependent: the bar-side $\kappa$ and shadow tower are
associator-free; the ordered cobar coproduct depends on the chosen
Drinfeld associator $\Phi$.
\end{remark}

\subsection{Level-prefix discipline for ordered Yangians}
\label{subsec:ordered-kd-core-rmatrix-level-prefix}

\begin{remark}[Level-prefix discipline; \ClaimStatusProvedElsewhere]
\label{rem:ordered-kd-core-rmatrix-level-prefix}
The ordered KD core's spectral $r$-matrix carries an explicit
prefix in every appearance. For $\cA = U^{\mathrm{ch}}(\fg_k)$:
\begin{enumerate}[label=\textup{(\roman*)}]
\item Trace-form: $r(z) = k\,\Omega/z$, vanishing at $k=0$.
\item KZ normalisation: $r(z) = \Omega/((k+h^\vee)z)$, equal to
$\Omega/(h^\vee z)$ at $k=0$.
\item Bridge identity: $k\,\Omega_{\mathrm{tr}} = \Omega/(k+h^\vee)$
at generic $k$.
\end{enumerate}
For Heisenberg $\cH_k$, the simple-pole $r(z) = k/z$ gives the
ordered bar differential
$d(s^{-1}\alpha\otimes s^{-1}\alpha) = \pm k\, s^{-1}(1)$
already noted in the chapter opening; the level prefix $k$ is
load-bearing because the $R$-matrix is $R(z) = z^k$, nontrivial
unless $k=0$.
\end{remark}

\subsection{Averaging map and the ordered--symmetric reduction}
\label{subsec:ordered-kd-core-averaging}

\begin{proposition}[Averaging recovers Vol~I modular shadow;
\ClaimStatusProvedElsewhere]
\label{prop:ordered-kd-core-averaging-modular}
The averaging map
$\mathrm{av}\colon \mathfrak{g}^{E_1}_\cA \to
\mathfrak{g}^{\mathrm{mod}}_\cA$ takes the ordered $r$-matrix
$r(z)$ to the modular characteristic $\kappa(\cA)$ at degree $2$,
the binary collision projection of the universal MC element
$\Theta_\cA$. For $\cA$ abelian or scalar (Heisenberg, Virasoro),
$\mathrm{av}(r(z)) = \kappa(\cA)$; for $\cA$ non-abelian
Kac--Moody, $\mathrm{av}(r(z)) + \dim(\fg)/2 = \kappa(V_k(\fg))$
(the Sugawara shift). This is the lossy
$\Sigma_n$-coinvariant projection of Vol~I E$_1$-primacy directive,
applied to the ordered KD core.
\end{proposition}

\subsection{Theorem~A as Koszul reflection}
\label{subsec:ordered-kd-core-koszul-reflection-anchor}

\begin{remark}[Ordered chiral KD specialises the unified
Koszul reflection; \ClaimStatusProvedElsewhere]
\label{rem:ordered-kd-core-koszul-reflection}
The ordered bar--cobar adjunction
$\Barchord \dashv \Cobar$ presented in this chapter is the
ordered $E_1$-chiral specialisation of the Vol~I unified
Koszul reflection theorem
(\cref{V1-thm:koszul-reflection}): the cofree conilpotent
ordered bar coalgebra $\Barchord(\cA) = T^c(s^{-1}\bar\cA)$ is the
left adjoint of $\Cobar$ on the conilpotent-complete locus, and
the counit $\Cobar(\Barchord(\cA)) \to \cA$ is a chain-level
quasi-isomorphism on the Koszul locus. The unordered reduction
along $\mathrm{av}$ recovers the symmetric $E_\infty$-bar
$B^\Sigma(\cA)$ of Vol~I via R-twisted $\Sigma_n$-descent
(prop:R-twisted-sigma-n-descent of the Theorem~A$^{\infty,2}$
reconstitution).
\end{remark}

\begin{remark}[Bridge to the Vol~II climax]
\label{rem:ordered-kd-core-uhm-bridge}
The ordered cobar $Y(\fg)^{\mathrm{ch}} =
\Cobar(\Barchord(U^{\mathrm{ch}}(\fg[t])))$ is the open-colour
boundary input to
\cref{thm:universal-holography-master} (Vol~II
\cref{ch:programme-climax-platonic}): the universal holography functor
$\Phi_{\mathrm{hol}}$ takes the chiral Yangian $Y(\fg)^{\mathrm{ch}}$
and produces the bulk $3$d HT theory $T_Y$ whose Wilson lines
recover the original line-operator algebra. The ordered
deconcatenation coproduct on $\Barchord$ is the open-colour
shadow of the bulk SC$^{\mathrm{ch,top}}$ braiding; the duality
involution $\cA \leftrightarrow \cA^!$
(Theorem~\ref{thm:duality-involution}) is the bulk--boundary swap
of $\Phi_{\mathrm{hol}}$.
\end{remark}

\begin{remark}[Bridge to the universal conductor]
\label{rem:ordered-kd-core-conductor-bridge}
The ordered KD core sees the Trinity ghost-charge conductor
$K(\cA) := c(\cA) + c(\cA^!) = -c_{\mathrm{ghost}}(\BRST)$ of Vol~I
\cref{V1-chap:universal-conductor} through the duality involution.
The scalar complementarity invariant
$\varkappa(\cA) := \kappa(\cA) + \kappa(\cA^!)$ is a DISTINCT
family-specific quantity ($\varkappa = 0$ for KM/free, $\varkappa = 13$
for Vir, $\varkappa = 98/3$ for BP, $\varkappa = 250/3$ for $\cW_3$, as
catalogued in Vol~I landscape census), related to the Trinity conductor
by $\varkappa = \varrho_{\cA}\cdot K$ with
$\varrho_{\mathrm{KM}} = \varrho_{\mathrm{free}} = 0$,
$\varrho_{\cW_N} = H_N - 1$, $\varrho_{\mathrm{BP}} = 1/6$. The
Trinity takes values $K(\Vir) = 26$, $K(\cW_3) = 100$,
$K(\mathrm{BP}) = 196$, $K(V_k(\fg)) = 2\dim\fg$. The cohomological
reflection $\cA \mapsto \cA^!$ implemented by $\Cobar\circ
\Barchord$ realises both stratifications on the ordered side. The Vol~I climax binds $\kappa$ (degree-$2$ shadow)
to the conductor $K$ via genus-tower complementarity; the
ordered chiral KD core supplies the chain-level lift.
\end{remark}

\section{Wave 13 DNA perspective: core E_1-chiral Hall--Drinfeld double}
\label{sec:oackd-core-wave13-DNA}

\begin{remark}[E_1-native at $d = 3$ with E_2 via Drinfeld centre]
\label{rem:oackd-core-DNA-E1-E2}
At the Vol~III K3 chiral bialgebra $\mathbf{H}_{\Delta_5}$ ($d = 3$), the native operadic level
is $E_1$ (CY-to-chiral correspondence, $\Phi_3$). The braided $E_2$-structure lives on
the Drinfeld centre
\[
  \mathcal{Z}(\mathrm{Rep}^{E_1}(\mathbf{H}_{\Delta_5})) \;\simeq\; \mathrm{Rep}^{E_2}(Z^{\mathrm{der}}_{\mathrm{ch}}(\mathbf{H}_{\Delta_5}))^{\mathrm{centred}},
\]
with half-braiding $\sigma_{V_u}(V_v) = R_{\mathrm{Sieg,dyn}}(u - v)$. This is the $d \ge 3$
pattern of Chapter~\texttt{cy\_to\_chiral} property (U$3$). Attributing $E_2$-structure to
$\mathbf{H}_{\Delta_5}$ itself is a category error.
\end{remark}

\begin{remark}[Hall--Drinfeld double as fourth quasi-Hopf taxonomic slot]
\label{rem:oackd-core-DNA-hall-drinfeld}
The Vol~III Wave~13 platonic chapter pins the core quantum-group deformation of
$\mathfrak{g}_{\Delta_5}$ as the K3 chiral Hall--Drinfeld double
$\mathcal{D}_\hbar(\mathcal{Y}^{\mathrm{Hall}}_\hbar(\mathrm{CoHA}_{K3 \times E}))$, with
Schiffmann--Vasserot $2012$ shuffle (Davison $2017$ CY-$3$ extension). Fourth taxonomic
slot beyond Drinfeld--Jimbo, Yangians, RTT, and Manin-triple quasi-Hopf. Manin pair
(not triple) because signature-$(2,1)$ Cartan forbids Lagrangian decomposition.
$H^2(\mathfrak{g}_{\Delta_5})^{\mathbb{Z}/2, K(1)} \cong \mathbb{C} \cdot \Delta_5$: the Igusa
cusp form IS the gauge-class invariant.
\end{remark}

\section{Wave 13 DNA perspective (deep saturation)}
\label{sec:oackd-wave13-DNA-deep}

\begin{remark}[Ordered associative $E_1$-bar on $\mathbf H_{\Delta_5}$]
\label{rem:oackd-DNA-deep-1}
The ordered associative $E_1$-bar core of this chapter applies to the
Wave~13 K3 bialgebra to produce
\[
B^{E_1,\mathrm{ord}}_{\mathrm{ch}}(\mathbf H_{\Delta_5}) \;=\; T^c(s^{-1}\overline{\mathbf H_{\Delta_5}})
\]
in the Hall subcategory of $\mathrm{CoHA}_{K3\times E}$ (note the
$E_1$-ordered direction: the tensor is taken in a \emph{non-symmetric}
monoidal context, since the Hall category is non-symmetric by
Schiffmann-Vasserot 2012). This is the source of the non-abelian
chiral structure. Cross-ref Vol.~I,
\S\ref{sec:barconstr-wave13-DNA}.
\end{remark}

\begin{remark}[Averaging and the symmetric shadow]
\label{rem:oackd-DNA-deep-2}
The averaging morphism
$\mathrm{av}:\mathfrak g^{E_1}_{\Delta_5}\to\mathfrak g^{\mathrm{mod}}_{\Delta_5}$
collapses the $E_1$-ordered Hall-non-commutativity of
$\mathbf H_{\Delta_5}$ onto its symmetric shadow. Under the averaging,
the $24$-fold Miki generators are antisymmetrised and the BKM
Lie-algebra structure emerges from the Hall Yangian. The Wave~13
$E_1$-ordered structure is therefore strictly finer than any
symmetric chiral algebra: it carries the non-abelian Hall data
(associator, $R$-matrix, Miki generators) that averaging erases.
Cross-ref Vol.~I, \S\ref{sec:e1mkd-wave13-DNA}.
\end{remark}
